\chapter{Operational Equations, Syntactic Transition--Output Quotients, and Algebraic Eilenberg Duality}
\label{ch:syntactic-transition-output-eilenberg}

\section{Overview and standing assumptions}
\label{sec:syntactic-to-overview-standing-assumptions}

\subsection{Orientation}
\label{subsec:syntactic-to-orientation}

The preceding chapters developed residual semantics at the level of states and
behaviours. Every Mealy morphism
\[
    P:\mathbb A\rightsquigarrow\mathbb B
\]
has a unique strict semantic homomorphism
\[
    \beta_P:
    P\to
    \mathbf{Beh}^{\mathrm{can}}_{\mathbb A,\mathbb B},
\]
and its regular image
\[
    \operatorname{Beh}(P)
    \hookrightarrow
    \mathbf{Beh}^{\mathrm{can}}_{\mathbb A,\mathbb B}
\]
is the local coequational theory of behaviours realized by \(P\). The
semantic image factorization is
\[
\begin{tikzcd}[column sep=large]
    P
        \arrow[r, two heads, "\eta_P"]
        \arrow[dr, "\beta_P"']
    &
    \operatorname{Beh}(P)
        \arrow[d, hook, "m_P"]
    \\
    &
    \mathbf{Beh}^{\mathrm{can}}_{\mathbb A,\mathbb B}.
\end{tikzcd}
\]

For a behaviour specification
\[
    k:K\to\mathbf{Beh}_{\mathbb A,\mathbb B,0},
\]
the residual evaluation map
\[
    d_K:D_K\to\mathbf{Beh}_{\mathbb A,\mathbb B,0}
\]
has regular image
\[
    \operatorname{Der}_0(K)
    \hookrightarrow
    \mathbf{Beh}_{\mathbb A,\mathbb B,0},
\]
which is the local coequational theory freely generated by the specified
behaviours:
\[
\begin{tikzcd}[column sep=large]
    D_K
        \arrow[r, two heads]
        \arrow[dr, "d_K"']
    &
    \operatorname{Der}_0(K)
        \arrow[d, hook]
    \\
    &
    \mathbf{Beh}_{\mathbb A,\mathbb B,0}.
\end{tikzcd}
\]

The present chapter develops the corresponding algebra of input operations.
Its central idea is that a Mealy structure on a carrier
\[
    A_0\xleftarrow{p}P\xrightarrow{q}B_0
\]
is a representation of the opposite input category
\[
    \mathbb A^{\mathrm{op}}
\]
by fibrewise output-comparison Kleisli operations on \(P\). An input arrow
\[
    a:u\to v
\]
acts on a state \(x\in P_v\) by
\[
    x\longmapsto
    \bigl(\delta(a,x),\omega(a,x)\bigr),
\]
where
\[
    \delta(a,x)\in P_u,
    \qquad
    \omega(a,x):
    q(\delta(a,x))\to q(x).
\]
The action is displayed by
\[
\begin{tikzcd}[column sep=huge]
    P_v
        \arrow[r, dashed,
        "{x\mapsto(\delta(a,x),\omega(a,x))}"]
    &
    P_u,
\end{tikzcd}
\]
where the dashed arrow denotes a Kleisli map rather than an ordinary map in
\(\mathcal E/B_0\).

For a local coequational theory
\[
    V\hookrightarrow
    \mathbf{Beh}^{\mathrm{can}}_{\mathbb A,\mathbb B},
\]
this gives an internal representation
\[
    \chi_V:
    \mathbb A^{\mathrm{op}}
    \to
    \operatorname{KlEnd}_{\mathbb B}(V).
\]
Its kernel congruence
\[
    R_V:=\ker(\chi_{V,1})
\]
consists of the parallel input arrows having the same transition--output
effect on every behaviour in \(V\). Its regular image
\[
    \operatorname{Op}_{\mathbb A,\mathbb B}(V)
    :=
    \operatorname{RegIm}(\chi_V)
\]
is the corresponding faithful transition--output quotient:
\[
\begin{tikzcd}[column sep=large, row sep=large]
    \mathbb A^{\mathrm{op}}
        \arrow[rr, "\chi_V"]
        \arrow[dr, two heads, "q_V"']
    &&
    \operatorname{KlEnd}_{\mathbb B}(V)
    \\
    &
    \operatorname{Op}_{\mathbb A,\mathbb B}(V)
        \arrow[ur, hook, "i_V"']
    &
\end{tikzcd}
\]

Conversely, an identity-on-objects operation congruence
\[
    R\rightrightarrows A_1
\]
determines a one-step equation predicate
\[
    \operatorname{Eq}(R)
    \hookrightarrow
    \mathbf{Beh}_{\mathbb A,\mathbb B,0}.
\]
The largest local coequational theory satisfying \(R\) is the derivative
safety core
\[
    V_R
    :=
    \operatorname{Safe}_{\partial}
    \bigl(\operatorname{Eq}(R)\bigr).
\]
The resulting operational polarity is
\[
    V\leq V_R
    \quad\Longleftrightarrow\quad
    R\leq R_V.
\]
Diagrammatically,
\[
\begin{tikzcd}[column sep=huge]
    \mathsf{Loc}_{\mathbb A,\mathbb B}
        \arrow[r, shift left=1.2ex, "{V\mapsto R_V}"]
    &
    \mathsf{Cong}_{\mathrm{ioo}}(\mathbb A^{\mathrm{op}})^{\mathrm{op}}
        \arrow[l, shift left=1.2ex, "{R\mapsto V_R}"]
\end{tikzcd}
\]
is a Galois connection. Its fixed points yield the local algebraic Eilenberg
duality.

For a specification \(K\), the syntactic transition--output category is
\[
    \operatorname{SynTO}_{\mathbb A,\mathbb B}(K)
    :=
    \operatorname{Op}_{\mathbb A,\mathbb B}
    \bigl(\operatorname{Der}_0(K)\bigr).
\]
It is the faithful identity-on-objects quotient of
\(\mathbb A^{\mathrm{op}}\) determined by all operational equations valid in
the residual behaviours generated by \(K\):
\[
\begin{tikzcd}[column sep=large]
    \mathbb A^{\mathrm{op}}
        \arrow[r, two heads]
    &
    \operatorname{SynTO}_{\mathbb A,\mathbb B}(K)
        \arrow[r, hook]
    &
    \operatorname{KlEnd}_{\mathbb B}
    \bigl(\operatorname{Der}_0(K)\bigr).
\end{tikzcd}
\]

The final parts of the chapter globalize this correspondence along input
functors and specialize it to monoids, stateless functors, Boolean languages,
syntactic monoids, and the classical Eilenberg correspondence
\cite{Eilenberg1974,AdamekMiliusMyersUrbat2019TOCL,Salamanca2017Unveiling}.

\subsection{Standing assumptions and notation}
\label{subsec:syntactic-to-standing-assumptions}

Throughout this chapter we use the assumptions and notation of
Chapters~\ref{ch:behaviour-categories-myhill-nerode} and
\ref{ch:coequational-subobjects-birkhoff}. Thus \(\mathcal E\) is a
Grothendieck topos with chosen pullbacks. In particular, \(\mathcal E\) is
Barr-exact and locally cartesian closed, regular epimorphisms are stable under
pullback, internal equivalence relations are effective, and every slice has
small limits, small colimits, and complete subobject lattices
\cite{Johnstone2002,Borceux1994,Jacobs1999CategoricalLogic}.

The input and output internal categories are
\[
    \mathbb A=(A_0,A_1,s_A,t_A,e_A,m_A),
    \qquad
    \mathbb B=(B_0,B_1,s_B,t_B,e_B,m_B).
\]
When \(\mathbb A\) and \(\mathbb B\) are fixed, write
\[
    Z:=A_0\times B_0.
\]
A Mealy carrier is an object of
\[
    \mathcal E/Z,
\]
displayed as a span
\[
\begin{tikzcd}
    &
    P
        \arrow[dl, "p"']
        \arrow[dr, "q"]
    &
    \\
    A_0
    &&
    B_0.
\end{tikzcd}
\]

Composition in \(\mathbb A\) is written
\[
    m_A(a,b)=b\circ a
\]
for composable arrows
\[
    a:u\to v,
    \qquad
    b:v\to w.
\]
The target-to-source execution convention is
\[
    \delta(b\circ a,x)
    =
    \delta(a,\delta(b,x))
\]
and
\[
    \omega(b\circ a,x)
    =
    \omega(b,x)\circ
    \omega(a,\delta(b,x)).
\]
The comparison arrows form the triangle
\[
\begin{tikzcd}[column sep=huge]
    q\bigl(\delta(a,\delta(b,x))\bigr)
        \arrow[
            r,
            "{\omega\bigl(a,\delta(b,x)\bigr)}"'
        ]
        \arrow[
            rr,
            bend left=25,
            "{\omega(b\circ a,x)}"
        ]
    &
    q\bigl(\delta(b,x)\bigr)
        \arrow[
            r,
            "{\omega(b,x)}"'
        ]
    &
    q(x).
\end{tikzcd}
\]

We regard \(A_1\), as the arrow object of \(\mathbb A^{\mathrm{op}}\), as an
object over
\[
    A_0\times A_0
\]
by
\[
    (t_A,s_A):A_1\to A_0\times A_0.
\]
Thus an arrow
\[
    a:u\to v
\]
of \(\mathbb A\) is read as an arrow
\[
    a^{\mathrm{op}}:v\to u
\]
of \(\mathbb A^{\mathrm{op}}\).

We write
\[
    \mathsf{Loc}_{\mathbb A,\mathbb B}
\]
for the complete lattice of local coequational theories
\[
    V\hookrightarrow
    \mathbf{Beh}^{\mathrm{can}}_{\mathbb A,\mathbb B},
\]
ordered by inclusion.

We write
\[
    \ker(f)
\]
for the kernel pair, or kernel congruence, of a morphism \(f\). Thus the
notation does not refer to a pointed kernel.

Whenever local theories or operation congruences are written with
\(\cong\), this denotes a canonical isomorphism of chosen representatives.
After passage to the corresponding subobject posets, these canonical
isomorphisms are regarded as equalities.

\section{Mealy morphisms as bicategorical representations}
\label{sec:to-bicategorical-representations}

\subsection{Internal categories as monads in spans}
\label{subsec:to-internal-categories-monads-spans}

Recall the bicategory
\[
    \mathbf{Span}(\mathcal E).
\]
Its objects are objects of \(\mathcal E\), its \(1\)-cells
\[
    X\rightsquigarrow Y
\]
are spans
\[
\begin{tikzcd}
    &
    S
        \arrow[dl, "\ell"']
        \arrow[dr, "r"]
    &
    \\
    X
    &&
    Y,
\end{tikzcd}
\]
and its \(2\)-cells are maps of spans. Horizontal composition is defined by
pullback.

An internal category
\[
    \mathbb A
\]
is equivalently a monad on \(A_0\) in \(\mathbf{Span}(\mathcal E)\), whose
underlying \(1\)-cell is
\[
\begin{tikzcd}
    &
    A_1
        \arrow[dl, "s_A"']
        \arrow[dr, "t_A"]
    &
    \\
    A_0
    &&
    A_0.
\end{tikzcd}
\]
Likewise, \(\mathbb B\) is a monad on \(B_0\).

A Mealy carrier
\[
    A_0\xleftarrow{p}P\xrightarrow{q}B_0
\]
is a \(1\)-cell
\[
    P:A_0\rightsquigarrow B_0.
\]
The composite
\[
    \mathbb A\odot P
\]
has apex
\[
    A_1\times_{t_A,A_0,p}P,
\]
defined by
\[
\begin{tikzcd}
    A_1\times_{A_0}P
        \arrow[r, "\pi_P"]
        \arrow[d, "\pi_{A_1}"']
        \arrow[dr, phantom, "\lrcorner", very near start]
    &
    P
        \arrow[d, "p"]
    \\
    A_1
        \arrow[r, "t_A"']
    &
    A_0.
\end{tikzcd}
\]
The composite
\[
    P\odot\mathbb B
\]
has apex
\[
    P\times_{q,B_0,s_B}B_1,
\]
defined by
\[
\begin{tikzcd}
    P\times_{B_0}B_1
        \arrow[r, "\pi_{B_1}"]
        \arrow[d, "\pi_P"']
        \arrow[dr, phantom, "\lrcorner", very near start]
    &
    B_1
        \arrow[d, "s_B"]
    \\
    P
        \arrow[r, "q"']
    &
    B_0.
\end{tikzcd}
\]

\begin{definition}[Right-lax monad morphism structure]
\label{def:to-right-lax-monad-morphism}
    A \textbf{right-lax monad morphism structure} from \(\mathbb A\) to
    \(\mathbb B\) on the span \(P\) is a \(2\)-cell
    \[
        \alpha_P:
        \mathbb A\odot P
        \Longrightarrow
        P\odot\mathbb B
    \]
    satisfying the following unit and multiplication axioms:
    \[
        \alpha_P
        \circ
        (\eta_{\mathbb A}\odot1_P)
        =
        1_P\odot\eta_{\mathbb B},
    \]
    and
    \[
        \alpha_P
        \circ
        (\mu_{\mathbb A}\odot1_P)
        =
        (1_P\odot\mu_{\mathbb B})
        \circ
        (\alpha_P\odot1_{\mathbb B})
        \circ
        (1_{\mathbb A}\odot\alpha_P),
    \]
    with the canonical associators and unitors of
    \(\mathbf{Span}(\mathcal E)\) understood.
\end{definition}

The multiplication axiom compares the two paths in
\[
\begin{tikzcd}[column sep=large, row sep=large]
    \mathbb A\odot\mathbb A\odot P
        \arrow[r, "\mu_{\mathbb A}\odot1_P"]
        \arrow[d, "1_{\mathbb A}\odot\alpha_P"']
    &
    \mathbb A\odot P
        \arrow[d, "\alpha_P"]
    \\
    \mathbb A\odot P\odot\mathbb B
        \arrow[r, "\alpha_P\odot1_{\mathbb B}"']
    &
    P\odot\mathbb B\odot\mathbb B
        \arrow[d, "1_P\odot\mu_{\mathbb B}"]
    \\
    &
    P\odot\mathbb B.
\end{tikzcd}
\]

\begin{remark}[Terminology]
\label{rem:to-lax-terminology}
    The term ``right-lax'' refers here to the direction
    \[
        \mathbb A\odot P
        \Longrightarrow
        P\odot\mathbb B
    \]
    under the execution-order convention for horizontal composition used in
    this chapter. Under the opposite convention for bicategorical
    composition, the same structure is often called a colax monad morphism or
    monad opmorphism.
\end{remark}

A map of spans
\[
    \alpha_P:
    A_1\times_{t_A,A_0,p}P
    \longrightarrow
    P\times_{q,B_0,s_B}B_1
\]
is equivalently a pair
\[
    \delta:
    A_1\times_{t_A,A_0,p}P\to P,
    \qquad
    \omega:
    A_1\times_{t_A,A_0,p}P\to B_1
\]
satisfying
\[
    p\delta=s_A\pi_{A_1},
\]
\[
    s_B\omega=q\delta,
\]
and
\[
    t_B\omega=q\pi_P.
\]
The defining map is
\[
\begin{tikzcd}[column sep=large]
    A_1\times_{A_0}P
        \arrow[r, "{\langle\delta,\omega\rangle}"]
        \arrow[dr, "\delta"']
    &
    P\times_{B_0}B_1
        \arrow[d, "\pi_P"]
    \\
    &
    P,
\end{tikzcd}
\]
and the output typing equations form
\[
\begin{tikzcd}[column sep=large]
    A_1\times_{A_0}P
        \arrow[r, "\omega"]
        \arrow[d, "\delta"']
        \arrow[dr, phantom, "\circlearrowright"]
    &
    B_1
        \arrow[d, "s_B"]
    \\
    P
        \arrow[r, "q"']
    &
    B_0
\end{tikzcd}
\qquad
\begin{tikzcd}[column sep=large]
    A_1\times_{A_0}P
        \arrow[r, "\omega"]
        \arrow[d, "\pi_P"']
        \arrow[dr, phantom, "\circlearrowright"]
    &
    B_1
        \arrow[d, "t_B"]
    \\
    P
        \arrow[r, "q"']
    &
    B_0.
\end{tikzcd}
\]

\begin{theorem}[Mealy structures as right-lax monad morphisms]
\label{thm:to-mealy-lax-monad-morphisms}
    Let
    \[
        A_0\xleftarrow{p}P\xrightarrow{q}B_0
    \]
    be a carrier span. To give a Mealy structure on \(P\) is equivalent to
    giving a right-lax monad morphism structure
    \[
        \alpha_P:
        \mathbb A\odot P
        \Longrightarrow
        P\odot\mathbb B.
    \]
\end{theorem}

\begin{proof}
    \leavevmode
    \begin{enumerate}
        \item \textbf{Extract the transition--output pair.}
        A \(2\)-cell \(\alpha_P\) is equivalently a pair
        \[
            (\delta,\omega)
        \]
        with
        \[
            \delta:
            A_1\times_{A_0}P\to P,
            \qquad
            \omega:
            A_1\times_{A_0}P\to B_1,
        \]
        satisfying the Mealy typing equations. It is represented by
        \[
        \begin{tikzcd}[column sep=large]
            A_1\times_{A_0}P
                \arrow[r, "{\langle\delta,\omega\rangle}"]
            &
            P\times_{B_0}B_1.
        \end{tikzcd}
        \]

        \item \textbf{Read the unit axiom.}
        Compatibility with the units of \(\mathbb A\) and \(\mathbb B\)
        gives, for \(x\in P_v\),
        \[
            \delta(e_Av,x)=x,
            \qquad
            \omega(e_Av,x)=e_B(qx).
        \]
        The comparison is
        \[
        \begin{tikzcd}[column sep=large]
            P
                \arrow[r, "{\langle e_Ap,1_P\rangle}"]
                \arrow[dr, "{\langle1_P,e_Bq\rangle}"']
            &
            A_1\times_{A_0}P
                \arrow[d, "{\langle\delta,\omega\rangle}"]
            \\
            &
            P\times_{B_0}B_1.
        \end{tikzcd}
        \]

        \item \textbf{Read the multiplication axiom.}
        For composable arrows
        \[
            a:u\to v,
            \qquad
            b:v\to w,
            \qquad
            x\in P_w,
        \]
        the multiplication axiom gives
        \[
            \delta(b\circ a,x)
            =
            \delta(a,\delta(b,x))
        \]
        and
        \[
            \omega(b\circ a,x)
            =
            \omega(b,x)\circ
            \omega(a,\delta(b,x)).
        \]
        The output equation is the commutativity of
        \[
        \begin{tikzcd}[column sep=large]
            q\bigl(\delta(a,\delta(b,x))\bigr)
                \arrow[
                    r,
                    "{\omega\bigl(a,\delta(b,x)\bigr)}"
                ]
                \arrow[
                    rr,
                    bend left=20,
                    "{\omega(b\circ a,x)}"'
                ]
            &
            q\bigl(\delta(b,x)\bigr)
                \arrow[
                    r,
                    "{\omega(b,x)}"
                ]
            &
            q(x).
        \end{tikzcd}
        \]

        \item \textbf{Reconstruct the bicategorical structure.}
        Conversely, a transition--output pair satisfying the Mealy typing,
        unit, and multiplication laws defines a map of spans
        \[
            \alpha_P:
            \mathbb A\odot P
            \Longrightarrow
            P\odot\mathbb B.
        \]
        The displayed unit and multiplication equations are precisely the
        right-lax monad morphism axioms. Hence the two constructions are
        mutually inverse.
    \end{enumerate}
\end{proof}

\begin{remark}[Bicategorical role of the chapter]
\label{rem:to-bicategorical-role}
    The bicategorical formulation identifies the transition--output structure
    as a representation-type structure carried by the span \(P\). The
    remaining constructions internalize the corresponding endomorphism
    object and take the faithful image of the induced representation.
\end{remark}

\subsection{The output-comparison Kleisli monad}
\label{subsec:to-output-comparison-kleisli-monad}

Right composition by the monad \(\mathbb B\) induces the
output-comparison monad
\[
    T_{\mathbb B}
    :=
    (t_B)_!(s_B)^*:
    \mathcal E/B_0\to\mathcal E/B_0.
\]
For
\[
    q_X:X\to B_0,
\]
the pullback part is
\[
\begin{tikzcd}
    X\times_{q_X,B_0,s_B}B_1
        \arrow[r, "\pi_{B_1}"]
        \arrow[d, "\pi_X"']
        \arrow[dr, phantom, "\lrcorner", very near start]
    &
    B_1
        \arrow[d, "s_B"]
    \\
    X
        \arrow[r, "q_X"']
    &
    B_0.
\end{tikzcd}
\]
Postcomposition with \(t_B\) gives
\[
    T_{\mathbb B}X
    =
    X\times_{q_X,B_0,s_B}B_1
    \xrightarrow{t_B\pi_{B_1}}
    B_0.
\]
A generalized element is a pair
\[
    (x,\alpha),
    \qquad
    \alpha:q_X(x)\to b.
\]

The unit is
\[
    \eta_X(x)
    =
    \bigl(x,e_B(q_Xx)\bigr),
\]
represented by
\[
\begin{tikzcd}[column sep=large]
    X
        \arrow[r, "\eta_X"]
        \arrow[dr, "q_X"']
    &
    T_{\mathbb B}X
        \arrow[d, "t_B\pi_{B_1}"]
    \\
    &
    B_0.
\end{tikzcd}
\]
A generalized element of \(T_{\mathbb B}^2X\) is represented by
\[
    (x,\alpha,\beta),
    \qquad
    \alpha:q_Xx\to b,
    \qquad
    \beta:b\to c,
\]
and multiplication is
\[
    \mu_X(x,\alpha,\beta)
    =
    (x,\beta\circ\alpha).
\]
The underlying comparison chain is
\[
\begin{tikzcd}[column sep=large]
    q_Xx
        \arrow[r, "\alpha"]
        \arrow[rr, bend right=20, "\beta\circ\alpha"']
    &
    b
        \arrow[r, "\beta"]
    &
    c.
\end{tikzcd}
\]

More generally, for any object \(S\), base extension gives a monad
\[
    T_{\mathbb B}^S
    :=
    (1_S\times t_B)_!
    (1_S\times s_B)^*
\]
on
\[
    \mathcal E/(S\times B_0).
\]

\begin{proposition}[Kleisli normal form]
\label{prop:to-kleisli-normal-form}
    Let
    \[
        q_X:X\to B_0,
        \qquad
        q_Y:Y\to B_0
    \]
    be objects of \(\mathcal E/B_0\). A Kleisli map
    \[
        X\rightsquigarrow Y
    \]
    for \(T_{\mathbb B}\) is equivalently a pair
    \[
        F:X\to Y,
        \qquad
        \alpha:X\to B_1
    \]
    satisfying
    \[
        s_B\alpha=q_YF,
        \qquad
        t_B\alpha=q_X.
    \]
    Thus
    \[
        \alpha_x:q_Y(Fx)\to q_X(x).
    \]

    If
    \[
        (F,\alpha):X\rightsquigarrow Y,
        \qquad
        (G,\gamma):Y\rightsquigarrow Z,
    \]
    then their Kleisli composite is
    \[
        (G,\gamma)\odot(F,\alpha)
        =
        (GF,\alpha\circ_F\gamma),
    \]
    where
    \[
        (\alpha\circ_F\gamma)_x
        =
        \alpha_x\circ\gamma_{Fx}.
    \]
\end{proposition}

\begin{proof}
    \leavevmode
    \begin{enumerate}
        \item \textbf{Unpack a Kleisli map.}
        A Kleisli map is a morphism
        \[
            X\to T_{\mathbb B}Y
            =
            Y\times_{q_Y,B_0,s_B}B_1
        \]
        in the slice \(\mathcal E/B_0\). Hence it is a pair
        \[
            (F,\alpha)
        \]
        fitting into
        \[
        \begin{tikzcd}
            X
                \arrow[r, "{\langle F,\alpha\rangle}"]
                \arrow[d, "F"']
            &
            Y\times_{B_0}B_1
                \arrow[d, "\pi_Y"]
            \\
            Y
                \arrow[r, equal]
            &
            Y.
        \end{tikzcd}
        \]
        The pullback condition gives
        \[
            s_B\alpha=q_YF,
        \]
        while the slice condition gives
        \[
            t_B\alpha=q_X.
        \]

        \item \textbf{Compute the comparison arrow.}
        For each generalized element \(x\), the typing equations give
        \[
        \begin{tikzcd}[column sep=large]
            q_Y(Fx)
                \arrow[r, "\alpha_x"]
            &
            q_X(x).
        \end{tikzcd}
        \]

        \item \textbf{Compute Kleisli composition.}
        Given
        \[
            (F,\alpha):X\rightsquigarrow Y,
            \qquad
            (G,\gamma):Y\rightsquigarrow Z,
        \]
        the comparison chain is
        \[
        \begin{tikzcd}[column sep=large]
            q_Z(GFx)
                \arrow[r, "\gamma_{Fx}"]
            &
            q_Y(Fx)
                \arrow[r, "\alpha_x"]
            &
            q_X(x).
        \end{tikzcd}
        \]
        Applying the multiplication of \(T_{\mathbb B}\) composes these
        arrows, so the resulting comparison is
        \[
            \alpha_x\circ\gamma_{Fx}.
        \]
        The underlying state map is \(GF\), proving the formula.
    \end{enumerate}
\end{proof}

\subsection{The internal Kleisli endomorphism category}
\label{subsec:to-internal-klend}

Let
\[
    A_0\xleftarrow{p}P\xrightarrow{q}B_0
\]
be a carrier. Put
\[
    A_0^{[2]}:=A_0\times A_0
\]
with projections
\[
    \ell,r:A_0^{[2]}\to A_0.
\]
A point
\[
    (v,u)\in A_0^{[2]}
\]
is read as the source and target of an operation
\[
    v\to u.
\]

Define the source-state and target-state families
\[
    P_{\mathrm{src}}
    :=
    A_0^{[2]}
    \times_{\ell,A_0,p}
    P
\]
and
\[
    P_{\mathrm{tgt}}
    :=
    A_0^{[2]}
    \times_{r,A_0,p}
    P.
\]
Write
\[
    \pi_{\mathrm{src}}:
    P_{\mathrm{src}}\to A_0^{[2]},
    \qquad
    \overline\pi_{\mathrm{src}}:
    P_{\mathrm{src}}\to P
\]
and
\[
    \pi_{\mathrm{tgt}}:
    P_{\mathrm{tgt}}\to A_0^{[2]},
    \qquad
    \overline\pi_{\mathrm{tgt}}:
    P_{\mathrm{tgt}}\to P
\]
for the pullback projections. The defining squares are
\[
\begin{tikzcd}
    P_{\mathrm{src}}
        \arrow[r, "\overline\pi_{\mathrm{src}}"]
        \arrow[d, "\pi_{\mathrm{src}}"']
        \arrow[dr, phantom, "\lrcorner", very near start]
    &
    P
        \arrow[d, "p"]
    \\
    A_0^{[2]}
        \arrow[r, "\ell"']
    &
    A_0
\end{tikzcd}
\qquad
\begin{tikzcd}
    P_{\mathrm{tgt}}
        \arrow[r, "\overline\pi_{\mathrm{tgt}}"]
        \arrow[d, "\pi_{\mathrm{tgt}}"']
        \arrow[dr, phantom, "\lrcorner", very near start]
    &
    P
        \arrow[d, "p"]
    \\
    A_0^{[2]}
        \arrow[r, "r"']
    &
    A_0.
\end{tikzcd}
\]
Over \((v,u)\), the two families are respectively \(P_v\) and \(P_u\).

Regard \(P_{\mathrm{src}}\) and \(P_{\mathrm{tgt}}\) as objects of
\[
    \mathcal E/(A_0^{[2]}\times B_0)
\]
by the structure maps
\[
    \widehat q_{\mathrm{src}}
    :=
    \left\langle
        \pi_{\mathrm{src}},
        q\overline\pi_{\mathrm{src}}
    \right\rangle:
    P_{\mathrm{src}}
    \to
    A_0^{[2]}\times B_0
\]
and
\[
    \widehat q_{\mathrm{tgt}}
    :=
    \left\langle
        \pi_{\mathrm{tgt}},
        q\overline\pi_{\mathrm{tgt}}
    \right\rangle:
    P_{\mathrm{tgt}}
    \to
    A_0^{[2]}\times B_0.
\]

Apply the base-extended output-comparison monad
\[
    T_{\mathbb B}^{[2]}
    :=
    (1_{A_0^{[2]}}\times t_B)_!
    (1_{A_0^{[2]}}\times s_B)^*
\]
on
\[
    \mathcal E/(A_0^{[2]}\times B_0)
\]
to \(P_{\mathrm{tgt}}\). Its underlying object is
\[
    T_{\mathbb B}^{[2]}(P_{\mathrm{tgt}})
    =
    P_{\mathrm{tgt}}
    \times_{
        q\overline\pi_{\mathrm{tgt}},\,
        B_0,\,
        s_B
    }
    B_1.
\]
Write
\[
    \pi_{P_{\mathrm{tgt}}}:
    T_{\mathbb B}^{[2]}(P_{\mathrm{tgt}})
    \to P_{\mathrm{tgt}}
\]
and
\[
    \pi_{B_1}:
    T_{\mathbb B}^{[2]}(P_{\mathrm{tgt}})
    \to B_1
\]
for its projections. Its structure map to
\(A_0^{[2]}\times B_0\) is
\[
    \widehat q_T
    :=
    \left\langle
        \pi_{\mathrm{tgt}}\pi_{P_{\mathrm{tgt}}},
        t_B\pi_{B_1}
    \right\rangle.
\]

A generalized element of
\(T_{\mathbb B}^{[2]}(P_{\mathrm{tgt}})\) is written
\[
    (v,u,y,\alpha)
\]
with
\[
    y\in P_u,
    \qquad
    \alpha:q(y)\to b,
\]
and it lies over
\[
    ((v,u),b)\in A_0^{[2]}\times B_0.
\]

Define
\[
    E_{[2]}
    :=
    P_{\mathrm{src}}
    \times_{
        \widehat q_{\mathrm{src}},\,
        A_0^{[2]}\times B_0,\,
        \widehat q_T
    }
    T_{\mathbb B}^{[2]}(P_{\mathrm{tgt}}).
\]
Write
\[
    \varepsilon:
    E_{[2]}\to P_{\mathrm{src}}
\]
for the first pullback projection. The defining square is
\[
\begin{tikzcd}[column sep=large]
    E_{[2]}
        \arrow[r]
        \arrow[d, "\varepsilon"']
        \arrow[dr, phantom, "\lrcorner", very near start]
    &
    T_{\mathbb B}^{[2]}(P_{\mathrm{tgt}})
        \arrow[d, "\widehat q_T"]
    \\
    P_{\mathrm{src}}
        \arrow[r, "\widehat q_{\mathrm{src}}"']
    &
    A_0^{[2]}\times B_0.
\end{tikzcd}
\]

A generalized element of \(E_{[2]}\) over
\[
    (v,u,x)\in P_{\mathrm{src}},
    \qquad
    x\in P_v,
\]
is a pair
\[
    y\in P_u,
    \qquad
    \alpha:q(y)\to q(x).
\]

Since \(\mathcal E\) is locally cartesian closed, the pullback functor
\[
    \pi_{\mathrm{src}}^*:
    \mathcal E/A_0^{[2]}
    \to
    \mathcal E/P_{\mathrm{src}}
\]
has a right adjoint
\[
    \Pi_{\pi_{\mathrm{src}}}.
\]
Define
\[
    \operatorname{KlEnd}_{\mathbb B}(P)_1
    :=
    \Pi_{\pi_{\mathrm{src}}}
    \bigl(
        \varepsilon:E_{[2]}\to P_{\mathrm{src}}
    \bigr).
\]
This is an object over \(A_0^{[2]}\).

The counit of the dependent-product adjunction gives an evaluation map
\[
    \mathrm{ev}:
    \operatorname{KlEnd}_{\mathbb B}(P)_1
    \times_{A_0^{[2]}}
    P_{\mathrm{src}}
    \to
    E_{[2]}
\]
over \(P_{\mathrm{src}}\), displayed by
\[
\begin{tikzcd}[column sep=large]
    \operatorname{KlEnd}_{\mathbb B}(P)_1
    \times_{A_0^{[2]}}P_{\mathrm{src}}
        \arrow[r, "\mathrm{ev}"]
        \arrow[d]
    &
    E_{[2]}
        \arrow[d, "\varepsilon"]
    \\
    P_{\mathrm{src}}
        \arrow[r, equal]
    &
    P_{\mathrm{src}}.
\end{tikzcd}
\]

A generalized element
\[
    \varphi\in
    \operatorname{KlEnd}_{\mathbb B}(P)_1
\]
over \((v,u)\) therefore represents a family
\[
    x\in P_v
    \longmapsto
    \bigl(
        \varphi_0(x),
        \alpha^\varphi_x
    \bigr),
\]
where
\[
    \varphi_0(x)\in P_u
\]
and
\[
    \alpha^\varphi_x:
    q(\varphi_0(x))\to q(x).
\]
Equivalently, \(\varphi\) represents a Kleisli map
\[
    P_v\rightsquigarrow P_u.
\]

\begin{proposition}[Internal Kleisli endomorphism category]
\label{prop:to-internal-klend}
    Let
    \[
        \pi_{\operatorname{KlEnd}}:
        \operatorname{KlEnd}_{\mathbb B}(P)_1
        \to
        A_0\times A_0
    \]
    denote the structure map of the dependent product
    \[
        \operatorname{KlEnd}_{\mathbb B}(P)_1
        =
        \Pi_{\pi_{\mathrm{src}}}
        \bigl(
            \varepsilon:E_{[2]}\to P_{\mathrm{src}}
        \bigr).
    \]
    Then the object object \(A_0\) and arrow object
    \[
        \operatorname{KlEnd}_{\mathbb B}(P)_1
    \]
    form an internal category
    \[
        \operatorname{KlEnd}_{\mathbb B}(P).
    \]
    Its source and target maps are
    \[
        s_{\operatorname{KlEnd}}
        :=
        \ell\pi_{\operatorname{KlEnd}},
        \qquad
        t_{\operatorname{KlEnd}}
        :=
        r\pi_{\operatorname{KlEnd}}.
    \]
    Identities are induced by the unit of \(T_{\mathbb B}\), and composition is
    induced by parameterized Kleisli composition.
\end{proposition}

\begin{proof}
    \leavevmode
    \begin{enumerate}
        \item \textbf{Define source and target.}
        By construction,
        \[
            \operatorname{KlEnd}_{\mathbb B}(P)_1
            =
            \Pi_{\pi_{\mathrm{src}}}
            \bigl(
                E_{[2]}\to P_{\mathrm{src}}
            \bigr)
        \]
        is an object over
        \[
            A_0\times A_0
        \]
        with structure map
        \[
            \pi_{\operatorname{KlEnd}}:
            \operatorname{KlEnd}_{\mathbb B}(P)_1
            \to
            A_0\times A_0.
        \]
        Define
        \[
            s_{\operatorname{KlEnd}}
            :=
            \ell\pi_{\operatorname{KlEnd}},
            \qquad
            t_{\operatorname{KlEnd}}
            :=
            r\pi_{\operatorname{KlEnd}}.
        \]
        Thus an operation lying over \((v,u)\) is regarded as an internal
        arrow
        \[
            v\to u
        \]
        and represents a Kleisli map
        \[
            P_v\rightsquigarrow P_u.
        \]

        \item \textbf{Construct identities by transposition.}
        Regard \(A_0\) as an object over \(A_0\times A_0\) through the
        diagonal
        \[
            \Delta:A_0\to A_0\times A_0.
        \]
        Over \(v\in A_0\), consider the parameterized Kleisli identity
        \[
            x
            \longmapsto
            \bigl(
                x,e_B(qx)
            \bigr),
            \qquad
            x\in P_v.
        \]
        This defines a map from the pullback of \(P_{\mathrm{src}}\) along
        \(\Delta\) into the corresponding pullback of \(E_{[2]}\), over the
        pulled-back source-state family.

        Transposition through
        \[
            \pi_{\mathrm{src}}^*
            \dashv
            \Pi_{\pi_{\mathrm{src}}}
        \]
        therefore defines
        \[
            e_{\operatorname{KlEnd}}:
            A_0
            \to
            \operatorname{KlEnd}_{\mathbb B}(P)_1
        \]
        over the diagonal:
        \[
        \begin{tikzcd}
            A_0
                \arrow[r, "e_{\operatorname{KlEnd}}"]
                \arrow[dr, "\Delta"']
            &
            \operatorname{KlEnd}_{\mathbb B}(P)_1
                \arrow[d, "\pi_{\operatorname{KlEnd}}"]
            \\
            &
            A_0\times A_0.
        \end{tikzcd}
        \]
        Consequently,
        \[
            s_{\operatorname{KlEnd}}
            e_{\operatorname{KlEnd}}
            =
            1_{A_0}
            =
            t_{\operatorname{KlEnd}}
            e_{\operatorname{KlEnd}}.
        \]

        \item \textbf{Form the object of composable operations.}
        Put
        \[
            \operatorname{KlEnd}_{\mathbb B}(P)_2
            :=
            \operatorname{KlEnd}_{\mathbb B}(P)_1
            \times_{
                t_{\operatorname{KlEnd}},\,
                A_0,\,
                s_{\operatorname{KlEnd}}
            }
            \operatorname{KlEnd}_{\mathbb B}(P)_1.
        \]
        A generalized element consists of operations
        \[
            \varphi:P_v\rightsquigarrow P_u,
            \qquad
            \psi:P_u\rightsquigarrow P_w.
        \]
        Their composite must be an operation
        \[
            P_v\rightsquigarrow P_w.
        \]

        The composability square is
        \[
        \begin{tikzcd}
            \operatorname{KlEnd}_{\mathbb B}(P)_2
                \arrow[r, "\pi_2"]
                \arrow[d, "\pi_1"']
                \arrow[dr, phantom, "\lrcorner", very near start]
            &
            \operatorname{KlEnd}_{\mathbb B}(P)_1
                \arrow[d, "s_{\operatorname{KlEnd}}"]
            \\
            \operatorname{KlEnd}_{\mathbb B}(P)_1
                \arrow[r, "t_{\operatorname{KlEnd}}"']
            &
            A_0.
        \end{tikzcd}
        \]

        \item \textbf{Evaluate the two composable operations.}
        Pulling back the universal evaluation map along the first projection
        from
        \(\operatorname{KlEnd}_{\mathbb B}(P)_2\)
        evaluates \(\varphi\) on a source state \(x\in P_v\):
        \[
            x
            \longmapsto
            \bigl(
                \varphi_0(x),
                \alpha^\varphi_x
            \bigr),
        \]
        where
        \[
            \varphi_0(x)\in P_u
        \]
        and
        \[
            \alpha^\varphi_x:
            q(\varphi_0x)\to q(x).
        \]

        Pulling back the universal evaluation map along the second projection
        and substituting the state \(\varphi_0(x)\) evaluates \(\psi\):
        \[
            \varphi_0(x)
            \longmapsto
            \bigl(
                \psi_0(\varphi_0x),
                \alpha^\psi_{\varphi_0x}
            \bigr),
        \]
        where
        \[
            \alpha^\psi_{\varphi_0x}:
            q(\psi_0(\varphi_0x))
            \to
            q(\varphi_0x).
        \]

        Thus the two evaluations produce the composable comparison chain
        \[
        \begin{tikzcd}[column sep=large]
            q(\psi_0(\varphi_0x))
                \arrow[r, "\alpha^\psi_{\varphi_0x}"]
            &
            q(\varphi_0x)
                \arrow[r, "\alpha^\varphi_x"]
            &
            q(x).
        \end{tikzcd}
        \]

        \item \textbf{Apply Kleisli multiplication.}
        Applying the multiplication of the output-comparison monad
        \(T_{\mathbb B}\) composes the two comparison arrows. The resulting
        evaluated operation is
        \[
            x
            \longmapsto
            \left(
                \psi_0(\varphi_0x),
                \alpha^\varphi_x
                \circ
                \alpha^\psi_{\varphi_0x}
            \right).
        \]
        Equivalently,
        \[
            (\psi\odot\varphi)(x)
            =
            \left(
                \psi_0(\varphi_0x),
                \alpha^\varphi_x
                \circ
                \alpha^\psi_{\varphi_0x}
            \right).
        \]

        This construction is natural in the parameter object
        \(\operatorname{KlEnd}_{\mathbb B}(P)_2\) and defines a map from the
        outer source-state family to the corresponding outer operation-data
        object.

        \item \textbf{Transpose the evaluated composite.}
        Transposition through the dependent-product adjunction for the outer
        source-state projection yields a unique map
        \[
            m_{\operatorname{KlEnd}}:
            \operatorname{KlEnd}_{\mathbb B}(P)_2
            \to
            \operatorname{KlEnd}_{\mathbb B}(P)_1
        \]
        whose base map sends the outer endpoints
        \[
            (v,u,w)\longmapsto(v,w).
        \]
        By construction, its evaluation is precisely parameterized Kleisli
        composition. In particular,
        \[
            s_{\operatorname{KlEnd}}
            m_{\operatorname{KlEnd}}
            =
            s_{\operatorname{KlEnd}}\pi_1
        \]
        and
        \[
            t_{\operatorname{KlEnd}}
            m_{\operatorname{KlEnd}}
            =
            t_{\operatorname{KlEnd}}\pi_2.
        \]

        \item \textbf{Verify associativity after evaluation.}
        Consider the two maps from the object of composable triples to
        \(\operatorname{KlEnd}_{\mathbb B}(P)_1\) obtained from the two
        possible parenthesizations of \(m_{\operatorname{KlEnd}}\).

        Under dependent-product transposition, their evaluations are the two
        parenthesizations of the Kleisli composite of three operations. These
        evaluated maps are equal by associativity of Kleisli composition,
        equivalently by associativity of the multiplication of
        \(T_{\mathbb B}\).

        Dependent-product transposition is a bijection of hom-sets.
        Therefore equality of the evaluated transposes implies equality of
        the original maps. Hence \(m_{\operatorname{KlEnd}}\) is associative.

        \item \textbf{Verify the unit laws after evaluation.}
        The evaluations of
        \[
            m_{\operatorname{KlEnd}}
            (e_{\operatorname{KlEnd}},1)
            \qquad\text{and}\qquad
            m_{\operatorname{KlEnd}}
            (1,e_{\operatorname{KlEnd}})
        \]
        are respectively left and right Kleisli composition with the
        Kleisli identity
        \[
            x\longmapsto
            \bigl(
                x,e_B(qx)
            \bigr).
        \]
        They therefore agree with the universal evaluation map by the two
        unit laws of \(T_{\mathbb B}\).

        Bijectivity of dependent-product transposition again implies the
        corresponding equalities before evaluation. Thus the left and right
        internal-category unit laws hold.

        Hence
        \[
            \operatorname{KlEnd}_{\mathbb B}(P)
        \]
        is an internal category.
    \end{enumerate}
\end{proof}

\begin{proposition}[Internal representation property]
\label{prop:to-klend-representation-property}
    Let
    \[
        u:U\to A_0\times A_0
    \]
    be an object of
    \[
        \mathcal E/(A_0\times A_0).
    \]
    Put
    \[
        P_{\mathrm{src},U}
        :=
        U
        \times_{A_0\times A_0}
        P_{\mathrm{src}},
    \]
    and write
    \[
        \pi_U:
        P_{\mathrm{src},U}
        \to
        P_{\mathrm{src}}
    \]
    and
    \[
        \pi_{\mathrm{src},U}:
        P_{\mathrm{src},U}
        \to
        U
    \]
    for the pullback projections.

    Define the operation-data object over \(P_{\mathrm{src},U}\) by
    \[
        E_{[2],U}
        :=
        P_{\mathrm{src},U}
        \times_{P_{\mathrm{src}}}
        E_{[2]},
    \]
    where the pullback is taken along
    \[
        \pi_U:
        P_{\mathrm{src},U}\to P_{\mathrm{src}}
    \]
    and
    \[
        \varepsilon:
        E_{[2]}\to P_{\mathrm{src}}.
    \]
    Thus there is a pullback square
    \[
    \begin{tikzcd}[column sep=large]
        E_{[2],U}
            \arrow[r]
            \arrow[d, "\varepsilon_U"']
            \arrow[dr, phantom, "\lrcorner", very near start]
        &
        E_{[2]}
            \arrow[d, "\varepsilon"]
        \\
        P_{\mathrm{src},U}
            \arrow[r, "\pi_U"']
        &
        P_{\mathrm{src}}.
    \end{tikzcd}
    \]
    Equivalently, there is a canonical isomorphism
    \[
        E_{[2],U}
        \cong
        U\times_{A_0\times A_0}E_{[2]}.
    \]

    Maps
    \[
        \varphi:
        U\to
        \operatorname{KlEnd}_{\mathbb B}(P)_1
    \]
    over \(A_0\times A_0\) correspond naturally to sections
    \[
        \widetilde\varphi:
        P_{\mathrm{src},U}
        \to
        E_{[2],U}
    \]
    of \(\varepsilon_U\), equivalently to maps
    \[
        P_{\mathrm{src},U}
        \to
        E_{[2]}
    \]
    over \(P_{\mathrm{src}}\).
\end{proposition}

\begin{proof}
    By definition,
    \[
        \operatorname{KlEnd}_{\mathbb B}(P)_1
        =
        \Pi_{\pi_{\mathrm{src}}}
        \bigl(
            \varepsilon:E_{[2]}\to P_{\mathrm{src}}
        \bigr).
    \]
    The dependent-product adjunction
    \[
        \pi_{\mathrm{src}}^*
        \dashv
        \Pi_{\pi_{\mathrm{src}}}
    \]
    gives a natural bijection
    \[
    \begin{aligned}
        &\operatorname{Hom}_{\mathcal E/(A_0\times A_0)}
        \left(
            U,
            \Pi_{\pi_{\mathrm{src}}}E_{[2]}
        \right)
        \\
        &\qquad\cong
        \operatorname{Hom}_{\mathcal E/P_{\mathrm{src}}}
        \left(
            P_{\mathrm{src}}
            \times_{A_0\times A_0}
            U,
            E_{[2]}
        \right).
    \end{aligned}
    \]
    The pullback
    \[
        P_{\mathrm{src}}
        \times_{A_0\times A_0}
        U
    \]
    is canonically \(P_{\mathrm{src},U}\). Hence maps
    \[
        U\to
        \operatorname{KlEnd}_{\mathbb B}(P)_1
    \]
    over \(A_0\times A_0\) correspond naturally to maps
    \[
        P_{\mathrm{src},U}
        \to
        E_{[2]}
    \]
    over \(P_{\mathrm{src}}\).

    By the universal property of the pullback defining \(E_{[2],U}\), such a
    map is equivalently a section
    \[
        \widetilde\varphi:
        P_{\mathrm{src},U}
        \to
        E_{[2],U}
    \]
    satisfying
    \[
        \varepsilon_U\widetilde\varphi
        =
        1_{P_{\mathrm{src},U}}.
    \]
    Naturality follows from the naturality of the dependent-product
    adjunction and of pullback.
\end{proof}

\begin{lemma}[Epimorphic extensionality]
\label{lem:to-regular-epi-extensionality}
    Let
    \[
        \varphi,\psi:
        U\to
        \operatorname{KlEnd}_{\mathbb B}(P)_1
    \]
    be maps over \(A_0\times A_0\), corresponding to evaluated maps
    \[
        \widetilde\varphi,\widetilde\psi:
        P_{\mathrm{src},U}
        \rightrightarrows
        E_{[2],U}.
    \]
    Let
    \[
        e:E\twoheadrightarrow P_{\mathrm{src},U}
    \]
    be an epimorphism. If
    \[
        \widetilde\varphi e
        =
        \widetilde\psi e,
    \]
    then
    \[
        \varphi=\psi.
    \]
\end{lemma}

\begin{proof}
    \leavevmode
    \begin{enumerate}
        \item \textbf{Cancel the epimorphism.}
        Since \(e\) is epic,
        \[
            \widetilde\varphi e
            =
            \widetilde\psi e
            \quad\Longrightarrow\quad
            \widetilde\varphi
            =
            \widetilde\psi.
        \]
        This cancellation is represented by
        \[
        \begin{tikzcd}[column sep=large]
            E
                \arrow[r, two heads, "e"]
                \arrow[dr, bend right=15,
                "\widetilde\varphi e=\widetilde\psi e"']
            &
            P_{\mathrm{src},U}
                \arrow[d, shift left=.7ex, "\widetilde\varphi"]
                \arrow[d, shift right=.7ex, "\widetilde\psi"']
            \\
            &
            E_{[2],U}.
        \end{tikzcd}
        \]

        \item \textbf{Transpose the equality.}
        The dependent-product correspondence of
        Proposition~\ref{prop:to-klend-representation-property} is bijective.
        Transposing
        \[
            \widetilde\varphi=\widetilde\psi
        \]
        therefore gives
        \[
            \varphi=\psi.
        \]
    \end{enumerate}
\end{proof}

\subsection{Mealy actions as internal representations}
\label{subsec:to-mealy-actions-representations}

\begin{theorem}[Transition--output representation theorem]
\label{thm:to-transition-output-representation}
    Let
    \[
        A_0\xleftarrow{p}P\xrightarrow{q}B_0
    \]
    be a carrier. To give a Mealy structure on \(P\) is equivalent to giving an
    identity-on-objects internal functor
    \[
        \chi_P:
        \mathbb A^{\mathrm{op}}
        \to
        \operatorname{KlEnd}_{\mathbb B}(P).
    \]
    Under this correspondence, an arrow
    \[
        a:u\to v
    \]
    of \(\mathbb A\), regarded as
    \[
        a^{\mathrm{op}}:v\to u,
    \]
    is sent to the Kleisli map
    \[
        P_v\rightsquigarrow P_u,
        \qquad
        x\longmapsto
        \bigl(\delta(a,x),\omega(a,x)\bigr).
    \]
\end{theorem}

\begin{proof}
    \leavevmode
    \begin{enumerate}
        \item \textbf{Transpose the Mealy action.}
        Given a Mealy structure, the typing equations make
        \[
            x\longmapsto
            \bigl(\delta(a,x),\omega(a,x)\bigr)
        \]
        a Kleisli map
        \[
            P_v\rightsquigarrow P_u
        \]
        for every arrow
        \[
            a:u\to v.
        \]
        By the internal representation property, these operations transpose
        to an arrow map
        \[
            \chi_{P,1}:
            A_1\to
            \operatorname{KlEnd}_{\mathbb B}(P)_1
        \]
        over
        \[
            (t_A,s_A):A_1\to A_0\times A_0:
        \]
        \[
        \begin{tikzcd}
            A_1
                \arrow[r, "\chi_{P,1}"]
                \arrow[dr, "{(t_A,s_A)}"']
            &
            \operatorname{KlEnd}_{\mathbb B}(P)_1
                \arrow[d,
                "{(s_{\operatorname{KlEnd}},
                   t_{\operatorname{KlEnd}})}"]
            \\
            &
            A_0\times A_0.
        \end{tikzcd}
        \]
        Together with the identity object map \(1_{A_0}\), this is the
        underlying internal graph morphism
        \[
            \mathbb A^{\mathrm{op}}
            \to
            \operatorname{KlEnd}_{\mathbb B}(P).
        \]

        \item \textbf{Verify identities.}
        The Mealy unit equations
        \[
            \delta(e_Av,x)=x,
            \qquad
            \omega(e_Av,x)=e_B(qx)
        \]
        say exactly that
        \[
            \chi_P(e_Av)
        \]
        is the Kleisli identity on \(P_v\). Hence \(\chi_P\) preserves
        identities.

        \item \textbf{Fix the order of composable arrows.}
        Let
        \[
            a:u\to v,
            \qquad
            b:v\to w
        \]
        be composable arrows of \(\mathbb A\). In
        \(\mathbb A^{\mathrm{op}}\) they become
        \[
            b^{\mathrm{op}}:w\to v,
            \qquad
            a^{\mathrm{op}}:v\to u.
        \]
        Thus the corresponding composable pair in the arrow object of
        \(\mathbb A^{\mathrm{op}}\) is represented, on the underlying arrow
        object \(A_1\), by the ordered pair
        \[
            (b,a)
            \in
            A_1\times_{s_A,A_0,t_A}A_1.
        \]
        Its composite is
        \[
            a^{\mathrm{op}}\circ b^{\mathrm{op}}
            =
            (b\circ a)^{\mathrm{op}},
        \]
        so the underlying arrow of \(\mathbb A\) is \(b\circ a\).

        \item \textbf{Verify composition.}
        Functoriality requires
        \[
            \chi_P(b\circ a)
            =
            \chi_P(a)\odot\chi_P(b).
        \]
        With the ordered pair \((b,a)\) made explicit, the functoriality
        square is
        \[
        \begin{tikzcd}[column sep=large, row sep=large]
            A_1\times_{s_A,A_0,t_A}A_1
                \arrow[r, "m_{\mathbb A^{\mathrm{op}}}"]
                \arrow[d,
                "{\langle
                    \chi_{P,1}\pi_1,
                    \chi_{P,1}\pi_2
                  \rangle}"']
            &
            A_1
                \arrow[d, "\chi_{P,1}"]
            \\
            \operatorname{KlEnd}_{\mathbb B}(P)_2
                \arrow[r, "m_{\operatorname{KlEnd}}"']
            &
            \operatorname{KlEnd}_{\mathbb B}(P)_1.
        \end{tikzcd}
        \]
        Here
        \[
            m_{\mathbb A^{\mathrm{op}}}(b,a)
            =
            b\circ a
        \]
        on the underlying arrow object, while
        \[
            m_{\operatorname{KlEnd}}
            \bigl(
                \chi_P(b),
                \chi_P(a)
            \bigr)
            =
            \chi_P(a)\odot\chi_P(b).
        \]

        Evaluation at \(x\in P_w\) first applies \(\chi_P(b)\):
        \[
            x
            \longmapsto
            \bigl(
                \delta(b,x),
                \omega(b,x)
            \bigr),
        \]
        and then applies \(\chi_P(a)\):
        \[
            \delta(b,x)
            \longmapsto
            \left(
                \delta(a,\delta(b,x)),
                \omega(a,\delta(b,x))
            \right).
        \]
        Kleisli composition therefore gives
        \[
            x
            \longmapsto
            \left(
                \delta(a,\delta(b,x)),
                \omega(b,x)\circ
                \omega(a,\delta(b,x))
            \right).
        \]
        Equality with the operation associated to \(b\circ a\) is exactly
        the pair of Mealy multiplication equations
        \[
            \delta(b\circ a,x)
            =
            \delta(a,\delta(b,x))
        \]
        and
        \[
            \omega(b\circ a,x)
            =
            \omega(b,x)\circ
            \omega(a,\delta(b,x)).
        \]
        Hence \(\chi_P\) preserves composition.

        \item \textbf{Recover the Mealy maps.}
        Conversely, let
        \[
            \chi_P:
            \mathbb A^{\mathrm{op}}
            \to
            \operatorname{KlEnd}_{\mathbb B}(P)
        \]
        be an identity-on-objects internal functor. Evaluating its arrow map
        on the source-state family yields maps
        \[
            \delta:
            A_1\times_{t_A,A_0,p}P\to P,
            \qquad
            \omega:
            A_1\times_{t_A,A_0,p}P\to B_1.
        \]
        The fact that the evaluated operation associated to
        \(a:u\to v\) is a Kleisli map
        \[
            P_v\rightsquigarrow P_u
        \]
        gives the typing equations
        \[
            p\delta=s_A\pi_{A_1},
        \]
        \[
            s_B\omega=q\delta,
        \]
        and
        \[
            t_B\omega=q\pi_P.
        \]

        Preservation of identities gives
        \[
            \delta(e_Av,x)=x,
            \qquad
            \omega(e_Av,x)=e_B(qx).
        \]
        Preservation of composition, evaluated on the ordered composable pair
        \((b,a)\), gives
        \[
            \delta(b\circ a,x)
            =
            \delta(a,\delta(b,x))
        \]
        and
        \[
            \omega(b\circ a,x)
            =
            \omega(b,x)\circ
            \omega(a,\delta(b,x)).
        \]
        Thus \((\delta,\omega)\) is a Mealy structure.

        \item \textbf{Verify mutual inverse correspondence.}
        Starting from a Mealy structure, transposing its evaluated operations
        and then evaluating the transpose recovers the original transition
        and output maps by the counit of the dependent-product adjunction.

        Conversely, starting from the internal functor \(\chi_P\), evaluating
        and then transposing recovers its original arrow map by the unit of
        the same adjunction. Hence the two constructions are mutually
        inverse.
    \end{enumerate}
\end{proof}

\begin{lemma}[Restriction of representations]
\label{lem:to-restriction-representations}
    Let
    \[
        i:V\hookrightarrow W
    \]
    be a strict monomorphism of Mealy morphisms over
    \((\mathbb A,\mathbb B)\). Then the representation \(\chi_V\) is the
    restriction of \(\chi_W\) to the state fibres of \(V\).
\end{lemma}

\begin{proof}
    \leavevmode
    \begin{enumerate}
        \item \textbf{Restrict transitions.}
        Strictness gives
        \[
            i\delta_V
            =
            \delta_W(1_{A_1}\times i),
        \]
        represented by
        \[
        \begin{tikzcd}
            A_1\times_{A_0}V
                \arrow[r, "1\times i"]
                \arrow[d, "\delta_V"']
            &
            A_1\times_{A_0}W
                \arrow[d, "\delta_W"]
            \\
            V
                \arrow[r, hook, "i"']
            &
            W.
        \end{tikzcd}
        \]

        \item \textbf{Restrict output comparisons.}
        Strictness also gives
        \[
            \omega_V
            =
            \omega_W(1_{A_1}\times i).
        \]
        Hence every input arrow acts on \(V\) by restricting its action on
        \(W\), which is precisely the asserted restriction of
        representations.
    \end{enumerate}
\end{proof}

\section{Operational equations and behavioural satisfaction}
\label{sec:to-operational-equations-satisfaction}

\subsection{Parallel pairs and operation congruences}
\label{subsec:to-parallel-pairs-congruences}

Define the object of parallel arrows of \(\mathbb A^{\mathrm{op}}\) by
\[
    \operatorname{Par}(\mathbb A)
    :=
    A_1
    \times_{A_0\times A_0}
    A_1,
\]
where each copy of \(A_1\) is regarded over \(A_0\times A_0\) by
\[
    (t_A,s_A).
\]
The defining pullback is
\[
\begin{tikzcd}
    \operatorname{Par}(\mathbb A)
        \arrow[r, "\pi_2"]
        \arrow[d, "\pi_1"']
        \arrow[dr, phantom, "\lrcorner", very near start]
    &
    A_1
        \arrow[d, "{(t_A,s_A)}"]
    \\
    A_1
        \arrow[r, "{(t_A,s_A)}"']
    &
    A_0\times A_0.
\end{tikzcd}
\]

\begin{definition}[Identity-on-objects operation congruence]
\label{def:to-operation-congruence}
    An \textbf{identity-on-objects operation congruence} on
    \(\mathbb A^{\mathrm{op}}\) is an internal equivalence relation
    \[
        \rho_0,\rho_1:
        R\rightrightarrows A_1
    \]
    in
    \[
        \mathcal E/(A_0\times A_0),
    \]
    where \(A_1\) is regarded over \(A_0\times A_0\) by
    \[
        (t_A,s_A):A_1\to A_0\times A_0,
    \]
    and which is compatible with composition in
    \(\mathbb A^{\mathrm{op}}\).

    We write
    \[
        \mathsf{Cong}_{\mathrm{ioo}}
        (\mathbb A^{\mathrm{op}})
    \]
    for the poset of operation congruences, ordered by inclusion.
\end{definition}

Because \(\rho_0\) and \(\rho_1\) are morphisms over
\(A_0\times A_0\), the relation \(R\) relates only parallel arrows.

No separate compatibility condition involving identity arrows is required.
Indeed, reflexivity of the equivalence relation already gives
\[
    e_{\mathbb A^{\mathrm{op}}}(v)
    \;R\;
    e_{\mathbb A^{\mathrm{op}}}(v)
\]
for every \(v\in A_0\).

Compatibility with composition may be stated as follows. Let
\[
    R_2
    :=
    R
    \times_{
        t_{\mathbb A^{\mathrm{op}}}\rho_0,\,
        A_0,\,
        s_{\mathbb A^{\mathrm{op}}}\rho_0
    }
    R.
\]
Since \(R\) relates parallel arrows, the same pullback is obtained by using
\(\rho_1\) in place of \(\rho_0\).

The two componentwise composition maps
\[
    R_2
    \rightrightarrows
    A_1,
\]
given internally by
\[
    (x\,R\,x',\,y\,R\,y')
    \longmapsto
    y\circ_{\mathbb A^{\mathrm{op}}}x
\]
and
\[
    (x\,R\,x',\,y\,R\,y')
    \longmapsto
    y'\circ_{\mathbb A^{\mathrm{op}}}x',
\]
must together factor through
\[
    R\hookrightarrow\operatorname{Par}(\mathbb A).
\]
Equivalently, whenever
\[
    x\,R\,x',
    \qquad
    y\,R\,y'
\]
and the corresponding pairs are composable, one has
\[
    y\circ_{\mathbb A^{\mathrm{op}}}x
    \;R\;
    y'\circ_{\mathbb A^{\mathrm{op}}}x'.
\]

The poset
\[
    \mathsf{Cong}_{\mathrm{ioo}}
    (\mathbb A^{\mathrm{op}})
\]
is a complete lattice. Meets are intersections of congruences. The join of a
family is the intersection of all operation congruences containing the union
of the given relations.

Since \(\mathcal E/(A_0\times A_0)\) is exact, every operation congruence is
effective as a relation on \(A_1\). In
Lemma~\ref{lem:to-quotient-operation-congruence}, compatibility with
identities and composition will equip the effective arrow quotient with an
internal category structure.

\subsection{Satisfaction of operational equations}
\label{subsec:to-satisfaction-operational-equations}

Let
\[
    H\in
    \mathbf{Beh}_{\mathbb A,\mathbb B,0}
\]
be a residual behaviour over the input object \(v\). Let
\[
    a,a':u\to v
\]
be parallel arrows of \(\mathbb A\).

\begin{definition}[Transition--output satisfaction]
\label{def:to-transition-output-satisfaction}
    We write
    \[
        H\models_{\mathrm{to}} a\equiv a'
    \]
    when
    \[
        \partial_aH=\partial_{a'}H
    \]
    and
    \[
        \omega_{\mathbf{Beh}}(a,H)
        =
        \omega_{\mathbf{Beh}}(a',H).
    \]
    The second equality is taken in the common hom-object determined by the
    first equality of derivatives.
\end{definition}

Equivalently, the two evaluated maps into the appropriate pullback of the
operation-data object \(E_{[2]}\) are equal. In this formulation, the
transition and output components are compared simultaneously and all typing
is automatic.

The satisfaction condition compares the two Kleisli operations
\[
\begin{tikzcd}[column sep=huge]
    \{H\}
        \arrow[r, shift left=.7ex, dashed,
        "{(\partial_aH,\omega_{\mathbf{Beh}}(a,H))}"]
        \arrow[r, shift right=.7ex, dashed,
        "{(\partial_{a'}H,\omega_{\mathbf{Beh}}(a',H))}"']
    &
    \mathbf{Beh}_u.
\end{tikzcd}
\]
Thus an operational equation asserts equality both of the transformed
behaviour and of the output comparison produced at the current residual state.

For a local coequational theory
\[
    V\hookrightarrow
    \mathbf{Beh}^{\mathrm{can}}_{\mathbb A,\mathbb B},
\]
define a relation \(R_V\) internally by
\[
    a\,R_V\,a'
    \quad\Longleftrightarrow\quad
    \forall H\in V_v,\;
    H\models_{\mathrm{to}}a\equiv a'.
\]

\begin{theorem}[Valid equations form the kernel congruence]
\label{thm:to-valid-equations-kernel-congruence}
    For every local coequational theory \(V\),
    \[
        R_V
        =
        \ker(\chi_{V,1})
    \]
    in
    \[
        \mathcal E/(A_0\times A_0).
    \]
    In particular, \(R_V\) is an identity-on-objects operation congruence on
    \(\mathbb A^{\mathrm{op}}\).
\end{theorem}

\begin{proof}
    \leavevmode
    \begin{enumerate}
        \item \textbf{Translate equality through evaluation.}
        By Proposition~\ref{prop:to-klend-representation-property}, equality
        \[
            \chi_V(a)=\chi_V(a')
        \]
        is equivalent to equality of the evaluated maps on the source-state
        family. The comparison is
        \[
        \begin{tikzcd}[column sep=large]
            V_v
                \arrow[r, shift left=.7ex, dashed, "\chi_V(a)"]
                \arrow[r, shift right=.7ex, dashed, "\chi_V(a')"']
            &
            V_u.
        \end{tikzcd}
        \]

        \item \textbf{Separate state and output components.}
        Equality of the evaluated Kleisli maps is exactly
        \[
            \partial_aH=\partial_{a'}H
        \]
        and
        \[
            \omega_{\mathbf{Beh}}(a,H)
            =
            \omega_{\mathbf{Beh}}(a',H)
        \]
        for every \(H\in V_v\). Hence the kernel pair of \(\chi_{V,1}\) is
        precisely \(R_V\).

        \item \textbf{Use functoriality.}
        Since \(\chi_V\) is an internal functor, its kernel relation is
        compatible with identities and composition.

        \item \textbf{Use exactness.}
        The slice \(\mathcal E/(A_0\times A_0)\) is exact, so the kernel
        relation is an effective equivalence relation. It is therefore an
        identity-on-objects operation congruence.
    \end{enumerate}
\end{proof}

\subsection{The one-step equation predicate}
\label{subsec:to-one-step-equation-predicate}

Let
\[
    R\in
    \mathsf{Cong}_{\mathrm{ioo}}
    (\mathbb A^{\mathrm{op}})
\]
with projections
\[
    \rho_0,\rho_1:
    R\rightrightarrows A_1.
\]
Because these projections are morphisms over
\(A_0\times A_0\), one has
\[
    (t_A,s_A)\rho_0
    =
    (t_A,s_A)\rho_1.
\]
Write
\[
    \tau_R:
    R\to A_0\times A_0
\]
for this common map:
\[
    \tau_R
    :=
    (t_A,s_A)\rho_0
    =
    (t_A,s_A)\rho_1.
\]
Write also
\[
    \sigma_R
    :=
    t_A\rho_0
    =
    t_A\rho_1:
    R\to A_0.
\]
This is the common source of the represented arrows in
\(\mathbb A^{\mathrm{op}}\).

Let
\[
    p_{\mathbf{Beh}}:
    \mathbf{Beh}_{\mathbb A,\mathbb B,0}
    \to A_0
\]
be the input anchor of the canonical behaviour carrier. Form
\[
    X_R
    :=
    R
    \times_{\sigma_R,A_0,p_{\mathbf{Beh}}}
    \mathbf{Beh}_{\mathbb A,\mathbb B,0}.
\]
Write
\[
    \pi_R^{R}:X_R\to R,
    \qquad
    \pi_R^{\mathbf{Beh}}:
    X_R\to
    \mathbf{Beh}_{\mathbb A,\mathbb B,0}
\]
for the pullback projections. Thus a generalized element of \(X_R\) is a
triple
\[
    (a,a',H)
\]
where
\[
    a\,R\,a'
\]
and \(H\) lies over the common source of \(a^{\mathrm{op}}\) and
\((a')^{\mathrm{op}}\).

Apply the construction of
Subsection~\ref{subsec:to-internal-klend} to the canonical behaviour carrier.
Write
\[
    \mathbf{Beh}_{\mathrm{src}}
    :=
    (A_0\times A_0)
    \times_{\ell,A_0,p_{\mathbf{Beh}}}
    \mathbf{Beh}_{\mathbb A,\mathbb B,0}
\]
and write
\[
    \varepsilon_{\mathbf{Beh}}:
    E_{[2]}^{\mathbf{Beh}}
    \to
    \mathbf{Beh}_{\mathrm{src}}
\]
for the corresponding operation-data projection.

There is a canonical map
\[
    \xi_R:
    X_R\to\mathbf{Beh}_{\mathrm{src}}
\]
defined by
\[
    \xi_R
    :=
    \left\langle
        \tau_R\pi_R^R,
        \pi_R^{\mathbf{Beh}}
    \right\rangle.
\]
Internally,
\[
    \xi_R(a,a',H)
    =
    \bigl(
        (t_Aa,s_Aa),
        H
    \bigr).
\]
The pullback condition defining \(X_R\) ensures that this map is well typed,
since
\[
    p_{\mathbf{Beh}}(H)=t_Aa.
\]

Define the operation-data object over \(X_R\) by
\[
    \mathcal O_R
    :=
    X_R
    \times_{
        \xi_R,\,
        \mathbf{Beh}_{\mathrm{src}},\,
        \varepsilon_{\mathbf{Beh}}
    }
    E_{[2]}^{\mathbf{Beh}}.
\]
The defining pullback is
\[
\begin{tikzcd}[column sep=large]
    \mathcal O_R
        \arrow[r]
        \arrow[d]
        \arrow[dr, phantom, "\lrcorner", very near start]
    &
    E_{[2]}^{\mathbf{Beh}}
        \arrow[d, "\varepsilon_{\mathbf{Beh}}"]
    \\
    X_R
        \arrow[r, "\xi_R"']
    &
    \mathbf{Beh}_{\mathrm{src}}.
\end{tikzcd}
\]

A generalized element of \(\mathcal O_R\) over
\[
    (a,a',H)\in X_R
\]
is a pair
\[
    (K,\alpha)
\]
where \(K\) lies over the common target
\[
    s_Aa=s_Aa'
\]
and
\[
    \alpha:q_{\mathbf{Beh}}(K)\to q_{\mathbf{Beh}}(H).
\]

Evaluation of the canonical representation along
\[
    \rho_0\pi_R^R
    \qquad\text{and}\qquad
    \rho_1\pi_R^R
\]
gives two sections
\[
    \operatorname{act}_0,
    \operatorname{act}_1:
    X_R\rightrightarrows\mathcal O_R.
\]
Internally, these sections are
\[
    \operatorname{act}_0(a,a',H)
    =
    \left(
        \partial_aH,
        \omega_{\mathbf{Beh}}(a,H)
    \right)
\]
and
\[
    \operatorname{act}_1(a,a',H)
    =
    \left(
        \partial_{a'}H,
        \omega_{\mathbf{Beh}}(a',H)
    \right).
\]

Let
\[
    e_R:E_R\hookrightarrow X_R
\]
be their equalizer:
\[
\begin{tikzcd}[column sep=large]
    E_R
        \arrow[r, hook, "e_R"]
    &
    X_R
        \arrow[r, shift left=.7ex, "\operatorname{act}_0"]
        \arrow[r, shift right=.7ex, "\operatorname{act}_1"']
    &
    \mathcal O_R.
\end{tikzcd}
\]

Regard \(X_R\) as an object of \(\mathcal E/Z\) through
\[
    X_R
    \xrightarrow{\pi_R^{\mathbf{Beh}}}
    \mathbf{Beh}_{\mathbb A,\mathbb B,0}
    \to
    Z.
\]
Then
\[
    \pi_R^{\mathbf{Beh}}:
    X_R\to
    \mathbf{Beh}_{\mathbb A,\mathbb B,0}
\]
is a morphism in \(\mathcal E/Z\). Since \(\mathcal E/Z\) is a topos, the
pullback map on subobjects along \(\pi_R^{\mathbf{Beh}}\) has a right
adjoint
\[
    \forall_{\pi_R^{\mathbf{Beh}}}.
\]

\begin{definition}[One-step equation predicate]
\label{def:to-equation-predicate}
    The \textbf{one-step equation predicate}
    \[
        \operatorname{Eq}(R)
        \hookrightarrow
        \mathbf{Beh}_{\mathbb A,\mathbb B,0}
    \]
    is
    \[
        \operatorname{Eq}(R)
        :=
        \forall_{\pi_R^{\mathbf{Beh}}}(E_R)
    \]
    in the subobject lattice of
    \(\mathbf{Beh}_{\mathbb A,\mathbb B,0}\) in \(\mathcal E/Z\).
\end{definition}

Equivalently, in the internal language,
\[
    H\in\operatorname{Eq}(R)
\]
if and only if, for every pair
\[
    a,a':u\to p_{\mathbf{Beh}}(H)
\]
with
\[
    a\,R\,a',
\]
one has
\[
    \partial_aH=\partial_{a'}H
\]
and
\[
    \omega_{\mathbf{Beh}}(a,H)
    =
    \omega_{\mathbf{Beh}}(a',H).
\]
Thus \(H\in\operatorname{Eq}(R)\) precisely when every equation in \(R\) is
valid at the current residual state \(H\).

The subobject
\[
    \operatorname{Eq}(R)
\]
need not be derivative-closed. An equation may hold at \(H\) while failing
at a residual derivative of \(H\). The invariant theory determined by \(R\)
is therefore obtained by applying the derivative safety core.

\subsection{The largest invariant theory satisfying a congruence}
\label{subsec:to-largest-invariant-theory}

Recall from Chapter~\ref{ch:coequational-subobjects-birkhoff} that
\[
    \operatorname{Safe}_{\partial}
\]
is right adjoint to the inclusion of derivative-closed subobjects:
\[
\begin{tikzcd}[column sep=huge]
    \mathsf{Loc}_{\mathbb A,\mathbb B}
        \arrow[r, shift left=1.2ex, "I_\partial"]
    &
    \operatorname{Sub}_{\mathcal E/Z}(\mathbf{Beh}_0)
        \arrow[l, shift left=1.2ex, "\operatorname{Safe}_\partial"]
\end{tikzcd}
\]
with
\[
    I_\partial\dashv\operatorname{Safe}_\partial.
\]
Thus, for every local coequational theory \(V\) and every subobject
\[
    C\hookrightarrow\mathbf{Beh}_0,
\]
one has
\[
    V\leq\operatorname{Safe}_{\partial}(C)
    \quad\Longleftrightarrow\quad
    V\leq C.
\]

\begin{definition}[Theory of a congruence]
\label{def:to-theory-of-congruence}
    For an operation congruence \(R\), define
    \[
        V_R
        :=
        \operatorname{Safe}_{\partial}
        \bigl(\operatorname{Eq}(R)\bigr).
    \]
\end{definition}

Thus \(V_R\) is the greatest derivative-closed subobject contained in the
one-step model predicate \(\operatorname{Eq}(R)\):
\[
\begin{tikzcd}[column sep=large]
    V_R
        \arrow[r, hook]
        \arrow[dr, hook]
    &
    \operatorname{Eq}(R)
        \arrow[d, hook]
    \\
    &
    \mathbf{Beh}_0.
\end{tikzcd}
\]

\begin{proposition}[Contextual form of \(V_R\)]
\label{prop:to-contextual-form-VR}
    Reasoning in the internal language, a residual behaviour \(H\) lies in
    \(V_R\) if and only if, for every residual context
    \[
        r:v\to p_{\mathbf{Beh}}(H)
    \]
    and every \(R\)-related pair
    \[
        a,a':u\to v,
    \]
    one has
    \[
        \partial_{r\circ a}H
        =
        \partial_{r\circ a'}H
    \]
    and
    \[
        \omega_{\mathbf{Beh}}(a,\partial_rH)
        =
        \omega_{\mathbf{Beh}}(a',\partial_rH).
    \]
    The latter arrows are precisely the output comparisons assigned by \(H\)
    to the residual morphisms
    \[
        r\circ a\to r,
        \qquad
        r\circ a'\to r.
    \]
\end{proposition}

\begin{proof}
    \leavevmode
    \begin{enumerate}
        \item \textbf{Expand safety-core membership.}
        Membership
        \[
            H\in V_R
        \]
        means that every derivative
        \[
            \partial_rH
        \]
        lies in \(\operatorname{Eq}(R)\).

        \item \textbf{Apply the one-step equation predicate.}
        At the derivative \(\partial_rH\), the relation \(a\,R\,a'\) gives
        \[
            \partial_a(\partial_rH)
            =
            \partial_{a'}(\partial_rH)
        \]
        and
        \[
            \omega_{\mathbf{Beh}}(a,\partial_rH)
            =
            \omega_{\mathbf{Beh}}(a',\partial_rH).
        \]

        \item \textbf{Compose derivatives.}
        The derivative composition law gives
        \[
            \partial_a(\partial_rH)
            =
            \partial_{r\circ a}H
        \]
        and similarly for \(a'\). The residual contexts fit into
        \[
        \begin{tikzcd}[column sep=large]
            u
                \arrow[r, shift left=.6ex, "a"]
                \arrow[r, shift right=.6ex, "a'"']
            &
            v
                \arrow[r, "r"]
            &
            p_{\mathbf{Beh}}(H).
        \end{tikzcd}
        \]
        Substituting these identities yields the claimed contextual
        condition.
    \end{enumerate}
\end{proof}

\begin{theorem}[Operational soundness and completeness]
\label{thm:to-operational-soundness-completeness}
    For every local coequational theory \(V\) and every operation congruence
    \(R\),
    \[
        V\leq V_R
        \quad\Longleftrightarrow\quad
        R\leq R_V.
    \]
\end{theorem}

\begin{proof}
    \leavevmode
    \begin{enumerate}
        \item \textbf{Use derivative invariance of \(V\).}
        Since \(V\) is derivative-closed, the safety-core adjunction gives
        \[
            V\leq V_R
            \quad\Longleftrightarrow\quad
            V\leq\operatorname{Eq}(R).
        \]

        \item \textbf{Interpret the predicate inclusion.}
        The inclusion
        \[
            V\leq\operatorname{Eq}(R)
        \]
        says that every equation in \(R\) is valid at every state of \(V\).

        \item \textbf{Identify the kernel congruence.}
        By Theorem~\ref{thm:to-valid-equations-kernel-congruence}, the
        equations valid at every state of \(V\) are precisely \(R_V\).
        Therefore
        \[
            V\leq\operatorname{Eq}(R)
            \quad\Longleftrightarrow\quad
            R\leq R_V.
        \]
    \end{enumerate}
\end{proof}

\begin{corollary}[Largest recognized theory]
\label{cor:to-largest-recognized-theory}
    The local coequational theory \(V_R\) is the largest local theory
    satisfying every operational equation in \(R\).
\end{corollary}

\begin{corollary}[Join description]
\label{cor:to-VR-join-description}
    One has
    \[
        V_R
        =
        \bigvee
        \left\{
            S\in\mathsf{Loc}_{\mathbb A,\mathbb B}
            \mid
            R\leq R_S
        \right\}.
    \]
\end{corollary}

\begin{proof}
    \leavevmode
    \begin{enumerate}
        \item \textbf{Identify the indexed family.}
        By Theorem~\ref{thm:to-operational-soundness-completeness},
        \[
            R\leq R_S
            \quad\Longleftrightarrow\quad
            S\leq V_R.
        \]

        \item \textbf{Take the join.}
        The displayed family is therefore exactly the family of local
        subtheories of \(V_R\). Their join in the complete lattice
        \(\mathsf{Loc}_{\mathbb A,\mathbb B}\) is \(V_R\).
    \end{enumerate}
\end{proof}

\section{Faithful operational images}
\label{sec:to-faithful-operational-images}

\subsection{Identity-on-objects quotients and regular images}
\label{subsec:to-ioo-regular-images}

\begin{definition}[Identity-on-objects regular epimorphism]
\label{def:to-ioo-regular-epi}
    Let
    \[
        F:\mathbb C\to\mathbb D
    \]
    be an internal functor whose object map is the identity on a common object
    object \(C_0=D_0\). We call \(F\) an
    \textbf{identity-on-objects regular epimorphism} if its arrow map is a
    regular epimorphism in
    \[
        \mathcal E/(C_0\times C_0).
    \]
    Identity-on-objects monomorphisms are defined similarly.
\end{definition}

Thus an identity-on-objects regular epimorphism is represented by
\[
\begin{tikzcd}
    C_1
        \arrow[r, two heads, "F_1"]
        \arrow[d, "{(s_C,t_C)}"']
    &
    D_1
        \arrow[d, "{(s_D,t_D)}"]
    \\
    C_0\times C_0
        \arrow[r, equal]
    &
    C_0\times C_0.
\end{tikzcd}
\]

\begin{lemma}[Quotient by an operation congruence]
\label{lem:to-quotient-operation-congruence}
    Let
    \[
        R\rightrightarrows C_1
    \]
    be an identity-on-objects operation congruence on an internal category
    \(\mathbb C\) with object object \(C_0\). Let
    \[
        q:C_1\twoheadrightarrow Q_1
    \]
    be the effective quotient of \(R\) in
    \[
        \mathcal E/(C_0\times C_0).
    \]
    Then \(Q_1\), with object object \(C_0\), carries a unique internal
    category structure such that
    \[
        q_R:
        \mathbb C\twoheadrightarrow\mathbb C/R
    \]
    is an identity-on-objects internal functor.
\end{lemma}

\begin{proof}
    \leavevmode
    \begin{enumerate}
        \item \textbf{Define source, target, and identities.}
        Since \(q\) is a morphism in
        \(\mathcal E/(C_0\times C_0)\), the structure map
        \[
            Q_1\to C_0\times C_0
        \]
        supplies source and target maps
        \[
            s_Q,t_Q:Q_1\rightrightarrows C_0.
        \]
        Define
        \[
            e_Q:=q e_C.
        \]
        The identity square is
        \[
        \begin{tikzcd}
            C_0
                \arrow[r, "e_C"]
                \arrow[dr, "e_Q"']
            &
            C_1
                \arrow[d, two heads, "q"]
            \\
            &
            Q_1.
        \end{tikzcd}
        \]

        \item \textbf{Construct the regular-epimorphic cover of composable pairs.}
        Put
        \[
            C_2
            :=
            C_1\times_{t_C,C_0,s_C}C_1,
            \qquad
            Q_2
            :=
            Q_1\times_{t_Q,C_0,s_Q}Q_1.
        \]
        The induced map
        \[
            q_2:C_2\to Q_2
        \]
        factors as
        \[
        \begin{tikzcd}[column sep=large]
            C_1\times_{C_0}C_1
                \arrow[r, two heads, "q\times1"]
            &
            Q_1\times_{t_Q,C_0,s_C}C_1
                \arrow[r, two heads, "1\times q"]
            &
            Q_1\times_{C_0}Q_1.
        \end{tikzcd}
        \]
        Each arrow is a pullback of \(q\), hence a regular epimorphism.
        Therefore \(q_2\) is a regular epimorphism.

        \item \textbf{Descend composition.}
        We claim that
        \[
            q m_C:C_2\to Q_1
        \]
        descends through \(q_2\). Reasoning internally, if
        \[
            q(a)=q(a'),
            \qquad
            q(b)=q(b'),
        \]
        then
        \[
            a\,R\,a',
            \qquad
            b\,R\,b'.
        \]
        Compatibility of \(R\) with composition gives
        \[
            m_C(a,b)\;R\;m_C(a',b'),
        \]
        and hence
        \[
            q m_C(a,b)=q m_C(a',b').
        \]
        Thus \(q m_C\) coequalizes the kernel pair of \(q_2\). There is a
        unique map
        \[
            m_Q:Q_2\to Q_1
        \]
        satisfying
        \[
            m_Qq_2=q m_C:
        \]
        \[
        \begin{tikzcd}[column sep=large]
            C_2
                \arrow[r, two heads, "q_2"]
                \arrow[d, "m_C"']
            &
            Q_2
                \arrow[d, dashed, "m_Q"]
            \\
            C_1
                \arrow[r, two heads, "q"']
            &
            Q_1.
        \end{tikzcd}
        \]

        \item \textbf{Verify associativity.}
        Let \(C_3\) and \(Q_3\) be the corresponding objects of composable
        triples. The induced map
        \[
            q_3:C_3\twoheadrightarrow Q_3
        \]
        is a composite of pullbacks of \(q\), hence a regular epimorphism.
        The two associativity composites on \(Q_3\) become equal after
        precomposition with \(q_3\), because they become \(q\) applied to the
        two associative composites in \(\mathbb C\):
        \[
        \begin{tikzcd}[column sep=large]
            C_3
                \arrow[r, two heads, "q_3"]
                \arrow[d, shift left=.8ex]
                \arrow[d, shift right=.8ex]
            &
            Q_3
                \arrow[d, shift left=.8ex]
                \arrow[d, shift right=.8ex]
            \\
            C_1
                \arrow[r, two heads, "q"']
            &
            Q_1.
        \end{tikzcd}
        \]
        Since \(q_3\) is epic, associativity holds in \(Q\).

        \item \textbf{Verify the unit laws.}
        The left and right unit equations in \(Q\) become the corresponding
        unit equations in \(\mathbb C\) after precomposition with
        \[
            q:C_1\twoheadrightarrow Q_1.
        \]
        Epimorphicity of \(q\) therefore proves the unit laws.

        \item \textbf{Prove uniqueness.}
        The maps \(q\), \(q_2\), and the higher-arity covers are epic.
        Consequently, source, target, identities, and composition are uniquely
        determined by the requirement that \(q_R\) be a functor.
    \end{enumerate}
\end{proof}

The resulting quotient has the expected universal property.

\begin{proposition}[Universal property of the operation quotient]
\label{prop:to-operation-quotient-universal-property}
    Let \(R\) be an operation congruence on \(\mathbb C\). If
    \[
        F:\mathbb C\to\mathbb D
    \]
    is an identity-on-objects internal functor whose arrow map coequalizes
    \(R\), then there is a unique identity-on-objects internal functor
    \[
        \overline F:
        \mathbb C/R\to\mathbb D
    \]
    satisfying
    \[
        \overline Fq_R=F.
    \]
\end{proposition}

\begin{proof}
    \leavevmode
    \begin{enumerate}
        \item \textbf{Descend the arrow map.}
        The quotient property in
        \(\mathcal E/(C_0\times C_0)\) gives a unique arrow map
        \[
            \overline F_1:Q_1\to D_1
        \]
        satisfying
        \[
            \overline F_1q=F_1:
        \]
        \[
        \begin{tikzcd}[column sep=large]
            C_1
                \arrow[r, two heads, "q"]
                \arrow[dr, "F_1"']
            &
            Q_1
                \arrow[d, dashed, "\overline F_1"]
            \\
            &
            D_1.
        \end{tikzcd}
        \]

        \item \textbf{Check identities.}
        One has
        \[
        \begin{aligned}
            \overline F_1e_Q
            &=
            \overline F_1q e_C
            \\
            &=
            F_1e_C
            \\
            &=
            e_D.
        \end{aligned}
        \]

        \item \textbf{Check composition.}
        The two candidate composites
        \[
            \overline F_1m_Q,
            \qquad
            m_D(\overline F_1\times\overline F_1)
        \]
        become equal after precomposition with the regular epimorphism
        \[
            q_2:C_2\twoheadrightarrow Q_2,
        \]
        because \(F\) preserves composition. Therefore they are equal.

        \item \textbf{Prove uniqueness.}
        Any other factorization has the same arrow map after precomposition
        with the epimorphism \(q\). Hence it equals \(\overline F\).
    \end{enumerate}
\end{proof}

\begin{lemma}[Creation of identity-on-objects regular images]
\label{lem:to-create-ioo-regular-image}
    Let
    \[
        F:\mathbb C\to\mathbb D
    \]
    be an identity-on-objects internal functor. Factor its arrow map in
    \[
        \mathcal E/(C_0\times C_0)
    \]
    as
    \[
        C_1
        \xrightarrow{r}
        I_1
        \xrightarrow{m}
        D_1,
    \]
    where \(r\) is a regular epimorphism and \(m\) is a monomorphism. Then
    \(I_1\), with object object \(C_0\), carries a unique internal category
    structure for which
    \[
        \mathbb C
        \xrightarrow{r}
        I
        \xrightarrow{m}
        \mathbb D
    \]
    are internal functors.

    Moreover,
    \[
        I\cong\mathbb C/\ker(F_1)
    \]
    canonically as an identity-on-objects internal category.
\end{lemma}

\begin{proof}
    \leavevmode
    \begin{enumerate}
        \item \textbf{Form the kernel congruence.}
        Since \(F\) is a functor,
        \[
            \ker(F_1)\rightrightarrows C_1
        \]
        is an operation congruence.

        \item \textbf{Take its effective quotient.}
        By Lemma~\ref{lem:to-quotient-operation-congruence}, the effective
        quotient
        \[
            C_1\twoheadrightarrow C_1/\ker(F_1)
        \]
        carries a unique internal category structure.

        \item \textbf{Identify the quotient with the regular image.}
        Because \(m\) is monic,
        \[
            \ker(F_1)=\ker(r).
        \]
        Exactness of the slice identifies the quotient of \(C_1\) by this
        kernel pair with the regular image \(I_1\). The comparison is
        \[
        \begin{tikzcd}[column sep=large]
            C_1
                \arrow[r, two heads, "r"]
                \arrow[dr, two heads]
            &
            I_1
                \arrow[d, "\cong"]
                \arrow[r, hook, "m"]
            &
            D_1
            \\
            &
            C_1/\ker(F_1).
        \end{tikzcd}
        \]

        \item \textbf{Transport the category structure.}
        Transporting the quotient category structure across the canonical
        isomorphism gives the required internal category structure on \(I_1\).

        \item \textbf{Verify uniqueness.}
        The regular epimorphism \(r\) determines the identities and
        composition of \(I\), while the monomorphism \(m\) detects the
        required equations. Hence the structure is unique.
    \end{enumerate}
\end{proof}

\subsection{The transition--output image}
\label{subsec:to-transition-output-image}

\begin{definition}[Transition--output image]
\label{def:to-transition-output-image}
    Let
    \[
        V\hookrightarrow
        \mathbf{Beh}^{\mathrm{can}}_{\mathbb A,\mathbb B}
    \]
    be a local coequational theory. Its \textbf{transition--output image} is
    the identity-on-objects regular image
    \[
        \operatorname{Op}_{\mathbb A,\mathbb B}(V)
        :=
        \operatorname{RegIm}
        \left(
            \chi_V:
            \mathbb A^{\mathrm{op}}
            \to
            \operatorname{KlEnd}_{\mathbb B}(V)
        \right).
    \]
\end{definition}

The defining factorization is
\[
\begin{tikzcd}[column sep=large, row sep=large]
    \mathbb A^{\mathrm{op}}
        \arrow[rr, "\chi_V"]
        \arrow[dr, two heads, "q_V"']
    &&
    \operatorname{KlEnd}_{\mathbb B}(V)
    \\
    &
    \operatorname{Op}_{\mathbb A,\mathbb B}(V)
        \arrow[ur, hook, "i_V"']
    &
\end{tikzcd}
\]
where the factorization is taken on arrow objects in
\[
    \mathcal E/(A_0\times A_0).
\]

\begin{theorem}[Transition--output quotient theorem]
\label{thm:to-transition-output-quotient}
    There is a canonical isomorphism of identity-on-objects internal categories
    \[
        \operatorname{Op}_{\mathbb A,\mathbb B}(V)
        \cong
        \mathbb A^{\mathrm{op}}/R_V.
    \]
    Under this isomorphism, the regular-epimorphic part \(q_V\) corresponds to
    the effective quotient map
    \[
        q_{R_V}:
        \mathbb A^{\mathrm{op}}
        \twoheadrightarrow
        \mathbb A^{\mathrm{op}}/R_V.
    \]
\end{theorem}

\begin{proof}
    \leavevmode
    \begin{enumerate}
        \item \textbf{Identify the kernel congruence.}
        By Theorem~\ref{thm:to-valid-equations-kernel-congruence},
        \[
            R_V=\ker(\chi_{V,1}).
        \]

        \item \textbf{Apply regular-image creation.}
        Lemma~\ref{lem:to-create-ioo-regular-image} identifies the regular
        image of \(\chi_V\) with the quotient of
        \(\mathbb A^{\mathrm{op}}\) by this kernel congruence:
        \[
        \begin{tikzcd}[column sep=large]
            \mathbb A^{\mathrm{op}}
                \arrow[r, two heads]
                \arrow[dr, two heads, "q_V"']
            &
            \mathbb A^{\mathrm{op}}/R_V
                \arrow[d, "\cong"]
            \\
            &
            \operatorname{Op}_{\mathbb A,\mathbb B}(V).
        \end{tikzcd}
        \]
    \end{enumerate}
\end{proof}

\subsection{Descent of operational equality}
\label{subsec:to-operational-descent}

\begin{proposition}[Operational descent along strict regular epimorphisms]
\label{prop:to-operational-descent}
    Let
    \[
        e:P\twoheadrightarrow Q
    \]
    be a strict regular epimorphism of Mealy morphisms over
    \((\mathbb A,\mathbb B)\). If two parallel input-arrow families induce the
    same transition--output operation on \(P\), then they induce the same
    transition--output operation on \(Q\).

    Equivalently,
    \[
        \ker(\chi_{P,1})
        \leq
        \ker(\chi_{Q,1}).
    \]
\end{proposition}

\begin{proof}
    \leavevmode
    \begin{enumerate}
        \item \textbf{Choose a generalized parallel pair.}
        Let
        \[
            a,a':U\to A_1
        \]
        be parallel input-arrow families, and suppose their evaluated
        operations on \(P\) agree.

        \item \textbf{Pull back the state quotient.}
        Pulling \(e\) back to the corresponding source-state families gives a
        regular epimorphism
        \[
            e_U:
            P_{\mathrm{src},U}
            \twoheadrightarrow
            Q_{\mathrm{src},U},
        \]
        in the pullback square
        \[
        \begin{tikzcd}
            P_{\mathrm{src},U}
                \arrow[r, two heads, "e_U"]
                \arrow[d]
                \arrow[dr, phantom, "\lrcorner", very near start]
            &
            Q_{\mathrm{src},U}
                \arrow[d]
            \\
            P
                \arrow[r, two heads, "e"']
            &
            Q.
        \end{tikzcd}
        \]

        \item \textbf{Use strictness.}
        Strictness gives
        \[
            e\,\delta_P(a,-)
            =
            \delta_Q(a,-)\,e_U
        \]
        and
        \[
            \omega_P(a,-)
            =
            \omega_Q(a,-)\,e_U,
        \]
        and similarly for \(a'\). The transition square is
        \[
        \begin{tikzcd}
            P_{\mathrm{src},U}
                \arrow[r, two heads, "e_U"]
                \arrow[d, "{\delta_P(a,-)}"']
            &
            Q_{\mathrm{src},U}
                \arrow[d, "{\delta_Q(a,-)}"]
            \\
            P
                \arrow[r, two heads, "e"']
            &
            Q.
        \end{tikzcd}
        \]

        \item \textbf{Descend transition equality.}
        Since
        \[
            \delta_P(a,-)=\delta_P(a',-),
        \]
        one has
        \[
        \begin{aligned}
            \delta_Q(a,-)\,e_U
            &=
            e\,\delta_P(a,-)
            \\
            &=
            e\,\delta_P(a',-)
            \\
            &=
            \delta_Q(a',-)\,e_U.
        \end{aligned}
        \]
        Epimorphicity of \(e_U\) gives
        \[
            \delta_Q(a,-)=\delta_Q(a',-).
        \]

        \item \textbf{Descend output equality.}
        Similarly,
        \[
        \begin{aligned}
            \omega_Q(a,-)\,e_U
            &=
            \omega_P(a,-)
            \\
            &=
            \omega_P(a',-)
            \\
            &=
            \omega_Q(a',-)\,e_U.
        \end{aligned}
        \]
        Hence
        \[
            \omega_Q(a,-)=\omega_Q(a',-).
        \]

        \item \textbf{Conclude kernel inclusion.}
        The two induced operations on \(Q\) are equal, so every pair in
        \(\ker(\chi_{P,1})\) lies in \(\ker(\chi_{Q,1})\).
    \end{enumerate}
\end{proof}

\subsection{Recognition and faithful recognition}
\label{subsec:to-recognition-faithful-recognition}

\begin{definition}[Recognition by an operation quotient]
\label{def:to-recognition-operation-quotient}
    Let
    \[
        q_R:
        \mathbb A^{\mathrm{op}}
        \twoheadrightarrow
        \mathbb A^{\mathrm{op}}/R
    \]
    be an operation quotient. We say that \(q_R\)
    \textbf{recognizes} a local coequational theory \(V\) if there is an
    identity-on-objects internal functor
    \[
        \overline\chi_V:
        \mathbb A^{\mathrm{op}}/R
        \to
        \operatorname{KlEnd}_{\mathbb B}(V)
    \]
    satisfying
    \[
        \overline\chi_Vq_R=\chi_V.
    \]
\end{definition}

Recognition is the commutative triangle
\[
\begin{tikzcd}[column sep=large]
    \mathbb A^{\mathrm{op}}
        \arrow[r, two heads, "q_R"]
        \arrow[dr, "\chi_V"']
    &
    \mathbb A^{\mathrm{op}}/R
        \arrow[d, dashed, "\overline\chi_V"]
    \\
    &
    \operatorname{KlEnd}_{\mathbb B}(V).
\end{tikzcd}
\]

\begin{proposition}[Recognition criterion]
\label{prop:to-recognition-criterion}
    For an operation quotient \(q_R\) and a local theory \(V\), the following
    are equivalent:
    \begin{enumerate}
        \item \(q_R\) recognizes \(V\);
        \item
        \[
            R\leq R_V;
        \]
        \item
        \[
            V\leq V_R.
        \]
    \end{enumerate}
\end{proposition}

\begin{proof}
    \leavevmode
    \begin{enumerate}
        \item \textbf{Translate factorization into kernel containment.}
        The representation \(\chi_V\) descends through \(q_R\) exactly when
        its arrow map coequalizes \(R\). This is equivalent to
        \[
            R\leq\ker(\chi_{V,1}).
        \]

        \item \textbf{Identify the valid-equation congruence.}
        By Theorem~\ref{thm:to-valid-equations-kernel-congruence},
        \[
            \ker(\chi_{V,1})=R_V.
        \]
        Hence recognition is equivalent to
        \[
            R\leq R_V.
        \]

        \item \textbf{Apply operational soundness and completeness.}
        Theorem~\ref{thm:to-operational-soundness-completeness} gives
        \[
            R\leq R_V
            \quad\Longleftrightarrow\quad
            V\leq V_R.
        \]
    \end{enumerate}
\end{proof}

\begin{definition}[Faithful recognition]
\label{def:to-faithful-recognition}
    A quotient
    \[
        q_R:
        \mathbb A^{\mathrm{op}}
        \twoheadrightarrow
        \mathbb A^{\mathrm{op}}/R
    \]
    \textbf{faithfully recognizes} \(V\) if it recognizes \(V\) and the
    induced functor
    \[
        \overline\chi_V:
        \mathbb A^{\mathrm{op}}/R
        \to
        \operatorname{KlEnd}_{\mathbb B}(V)
    \]
    is monic on arrows.
\end{definition}

\begin{theorem}[Universal faithful recognition]
\label{thm:to-universal-faithful-recognition}
    Let
    \[
        q_R:
        \mathbb A^{\mathrm{op}}
        \twoheadrightarrow
        \mathbb A^{\mathrm{op}}/R
    \]
    recognize a local theory \(V\). Then there is a unique
    identity-on-objects regular epimorphism
    \[
        h:
        \mathbb A^{\mathrm{op}}/R
        \twoheadrightarrow
        \operatorname{Op}_{\mathbb A,\mathbb B}(V)
    \]
    satisfying
    \[
        hq_R=q_V
    \]
    and
    \[
        i_Vh=\overline\chi_V.
    \]

    Moreover, \(q_R\) faithfully recognizes \(V\) if and only if
    \[
        R=R_V.
    \]
    Consequently, \(\operatorname{Op}_{\mathbb A,\mathbb B}(V)\) is, up to
    unique isomorphism, the unique faithful recognizing quotient of
    \(\mathbb A^{\mathrm{op}}\).
\end{theorem}

\begin{proof}
    \leavevmode
    \begin{enumerate}
        \item \textbf{Compare the kernel congruences.}
        Recognition gives
        \[
            R\leq R_V.
        \]

        \item \textbf{Induce the quotient comparison.}
        The quotient by \(R_V\) factors uniquely through the quotient by \(R\):
        \[
        \begin{tikzcd}[column sep=large, row sep=large]
            \mathbb A^{\mathrm{op}}
                \arrow[r, two heads, "q_R"]
                \arrow[dr, two heads, "q_{R_V}"']
            &
            \mathbb A^{\mathrm{op}}/R
                \arrow[d, dashed, "h"]
            \\
            &
            \mathbb A^{\mathrm{op}}/R_V.
        \end{tikzcd}
        \]

        \item \textbf{Show that the comparison is regular epic.}
        Factor \(h_1\) as
        \[
            h_1=mr
        \]
        with \(r\) regular epic and \(m\) monic. Since
        \[
            h_1q_R=q_{R_V}
        \]
        is a regular epimorphism, the monomorphic factor \(m\) is an
        isomorphism. Hence \(h_1\) is a regular epimorphism.

        \item \textbf{Identify the target with the transition--output image.}
        Theorem~\ref{thm:to-transition-output-quotient} gives
        \[
            \mathbb A^{\mathrm{op}}/R_V
            \cong
            \operatorname{Op}_{\mathbb A,\mathbb B}(V).
        \]
        Under this identification, the quotient comparison satisfies
        \[
            hq_R=q_V.
        \]

        \item \textbf{Compare the represented actions.}
        The two functors
        \[
            i_Vh,
            \qquad
            \overline\chi_V
        \]
        agree after precomposition with the regular epimorphism \(q_R\):
        \[
        \begin{tikzcd}[column sep=large]
            \mathbb A^{\mathrm{op}}
                \arrow[r, two heads, "q_R"]
                \arrow[dr, "\chi_V"']
            &
            \mathbb A^{\mathrm{op}}/R
                \arrow[d, shift left=.7ex, "i_Vh"]
                \arrow[d, shift right=.7ex,
                "\overline\chi_V"']
            \\
            &
            \operatorname{KlEnd}_{\mathbb B}(V).
        \end{tikzcd}
        \]
        Hence
        \[
            i_Vh=\overline\chi_V.
        \]

        \item \textbf{Characterize faithful recognition.}
        The induced representation factors as
        \[
        \begin{tikzcd}[column sep=large]
            \mathbb A^{\mathrm{op}}/R
                \arrow[r, two heads, "h"]
            &
            \mathbb A^{\mathrm{op}}/R_V
                \arrow[r, hook, "i_V"]
            &
            \operatorname{KlEnd}_{\mathbb B}(V).
        \end{tikzcd}
        \]
        Since \(i_V\) is monic, the composite is monic exactly when \(h\) is
        monic. Since \(h\) is already regular epic, this occurs exactly when
        \(h\) is an isomorphism.

        \item \textbf{Translate isomorphism into equality of congruences.}
        The quotient comparison \(h\) is an isomorphism exactly when
        \[
            \ker(q_R)=\ker(q_{R_V}),
        \]
        equivalently
        \[
            R=R_V.
        \]
        Uniqueness follows from the epimorphicity of \(q_R\).
    \end{enumerate}
\end{proof}

\begin{proposition}[Contravariance under inclusion]
\label{prop:to-op-contravariance}
    If
    \[
        V\leq W
    \]
    are local coequational theories, then
    \[
        R_W\leq R_V
    \]
    and there is a canonical identity-on-objects regular epimorphism
    \[
        \operatorname{Op}_{\mathbb A,\mathbb B}(W)
        \twoheadrightarrow
        \operatorname{Op}_{\mathbb A,\mathbb B}(V).
    \]
\end{proposition}

\begin{proof}
    \leavevmode
    \begin{enumerate}
        \item \textbf{Restrict valid equations.}
        Equality of operations on every state of \(W\) implies equality on
        the subobject \(V\). Hence
        \[
            R_W\leq R_V.
        \]

        \item \textbf{Pass to effective quotients.}
        The inclusion of congruences induces
        \[
        \begin{tikzcd}[column sep=large]
            \mathbb A^{\mathrm{op}}
                \arrow[r, two heads]
                \arrow[dr, two heads]
            &
            \mathbb A^{\mathrm{op}}/R_W
                \arrow[d, two heads]
            \\
            &
            \mathbb A^{\mathrm{op}}/R_V.
        \end{tikzcd}
        \]

        \item \textbf{Identify the quotient categories.}
        Apply Theorem~\ref{thm:to-transition-output-quotient} to identify the
        source and target with
        \[
            \operatorname{Op}(W)
            \quad\text{and}\quad
            \operatorname{Op}(V).
        \]
    \end{enumerate}
\end{proof}

\section{Local algebraic Eilenberg duality}
\label{sec:to-local-algebraic-eilenberg}

\subsection{Semantic closure operators}
\label{subsec:to-semantic-closure-operators}

\begin{definition}[Saturation and semantic closure]
\label{def:to-saturation-semantic-closure}
    For a local coequational theory \(V\), define its
    \textbf{operation saturation} by
    \[
        c(V):=V_{R_V}.
    \]
    We call \(V\) \textbf{operation-saturated} if
    \[
        c(V)=V.
    \]

    For an operation congruence \(R\), define its
    \textbf{semantic closure} by
    \[
        d(R):=R_{V_R}.
    \]
    We call \(R\) \textbf{semantically closed} if
    \[
        d(R)=R.
    \]
\end{definition}

The two composites are
\[
\begin{tikzcd}[column sep=huge]
    \mathsf{Loc}
        \arrow[r, "{V\mapsto R_V}"]
    &
    \mathsf{Cong}^{\mathrm{op}}
        \arrow[r, "{R\mapsto V_R}"]
    &
    \mathsf{Loc},
\end{tikzcd}
\]
and
\[
\begin{tikzcd}[column sep=huge]
    \mathsf{Cong}
        \arrow[r, "{R\mapsto V_R}"]
    &
    \mathsf{Loc}^{\mathrm{op}}
        \arrow[r, "{V\mapsto R_V}"]
    &
    \mathsf{Cong}.
\end{tikzcd}
\]

\begin{proposition}[Closure operators]
\label{prop:to-closure-operators}
    The assignments
    \[
        c:
        \mathsf{Loc}_{\mathbb A,\mathbb B}
        \to
        \mathsf{Loc}_{\mathbb A,\mathbb B}
    \]
    and
    \[
        d:
        \mathsf{Cong}_{\mathrm{ioo}}(\mathbb A^{\mathrm{op}})
        \to
        \mathsf{Cong}_{\mathrm{ioo}}(\mathbb A^{\mathrm{op}})
    \]
    are monotone, inflationary, and idempotent.
\end{proposition}

\begin{proof}
    \leavevmode
    \begin{enumerate}
        \item \textbf{Establish antitonicity of the polarity.}
        One has
        \[
            V\leq W
            \Longrightarrow
            R_W\leq R_V,
        \]
        and
        \[
            R\leq S
            \Longrightarrow
            V_S\leq V_R.
        \]
        Hence both composites \(c\) and \(d\) are monotone.

        \item \textbf{Prove inflationarity of \(c\).}
        Apply operational soundness and completeness to the reflexive
        inequality
        \[
            R_V\leq R_V.
        \]
        This gives
        \[
            V\leq V_{R_V}=c(V).
        \]

        \item \textbf{Prove inflationarity of \(d\).}
        Apply operational soundness and completeness to
        \[
            V_R\leq V_R.
        \]
        This gives
        \[
            R\leq R_{V_R}=d(R).
        \]

        \item \textbf{Prove idempotence of \(c\).}
        Since
        \[
            V\leq c(V),
        \]
        antitonicity gives
        \[
            R_{c(V)}\leq R_V.
        \]
        Inflationarity of \(d\), applied to \(R_V\), gives
        \[
            R_V
            \leq
            d(R_V)
            =
            R_{V_{R_V}}
            =
            R_{c(V)}.
        \]
        Hence
        \[
            R_{c(V)}=R_V,
        \]
        and therefore
        \[
        \begin{aligned}
            c^2(V)
            &=
            V_{R_{c(V)}}
            \\
            &=
            V_{R_V}
            \\
            &=
            c(V).
        \end{aligned}
        \]

        \item \textbf{Prove idempotence of \(d\).}
        Since
        \[
            R\leq d(R),
        \]
        antitonicity gives
        \[
            V_{d(R)}\leq V_R.
        \]
        Inflationarity of \(c\), applied to \(V_R\), gives
        \[
            V_R
            \leq
            c(V_R)
            =
            V_{R_{V_R}}
            =
            V_{d(R)}.
        \]
        Hence
        \[
            V_{d(R)}=V_R,
        \]
        and therefore
        \[
        \begin{aligned}
            d^2(R)
            &=
            R_{V_{d(R)}}
            \\
            &=
            R_{V_R}
            \\
            &=
            d(R).
        \end{aligned}
        \]
    \end{enumerate}
\end{proof}

\begin{corollary}[Galois triangle identities]
\label{cor:to-galois-triangle-identities}
    For every local theory \(V\) and operation congruence \(R\),
    \[
        R_{V_{R_V}}=R_V
    \]
    and
    \[
        V_{R_{V_R}}=V_R.
    \]
    In particular, every congruence of the form \(R_V\) is semantically
    closed, and every theory of the form \(V_R\) is operation-saturated.
\end{corollary}

\subsection{Complete-lattice identities}
\label{subsec:to-complete-lattice-identities}

\begin{theorem}[Galois formulas for joins]
\label{thm:to-galois-join-formulas}
    For every small family of local coequational theories \((V_i)_i\),
    \[
        R_{\bigvee_iV_i}
        =
        \bigwedge_iR_{V_i}.
    \]
    For every small family of operation congruences \((R_j)_j\),
    \[
        V_{\bigvee_jR_j}
        =
        \bigwedge_jV_{R_j}.
    \]
\end{theorem}

\begin{proof}
    \leavevmode
    \begin{enumerate}
        \item \textbf{Test the first equality against an arbitrary congruence.}
        For every operation congruence \(R\),
        \[
        \begin{aligned}
            R\leq R_{\bigvee_iV_i}
            &\Longleftrightarrow
            \bigvee_iV_i\leq V_R
            \\
            &\Longleftrightarrow
            \forall i,\;V_i\leq V_R
            \\
            &\Longleftrightarrow
            \forall i,\;R\leq R_{V_i}
            \\
            &\Longleftrightarrow
            R\leq\bigwedge_iR_{V_i}.
        \end{aligned}
        \]
        Hence
        \[
            R_{\bigvee_iV_i}
            =
            \bigwedge_iR_{V_i}.
        \]

        \item \textbf{Test the second equality against an arbitrary theory.}
        For every local theory \(V\),
        \[
        \begin{aligned}
            V\leq V_{\bigvee_jR_j}
            &\Longleftrightarrow
            \bigvee_jR_j\leq R_V
            \\
            &\Longleftrightarrow
            \forall j,\;R_j\leq R_V
            \\
            &\Longleftrightarrow
            \forall j,\;V\leq V_{R_j}
            \\
            &\Longleftrightarrow
            V\leq\bigwedge_jV_{R_j}.
        \end{aligned}
        \]
        Hence
        \[
            V_{\bigvee_jR_j}
            =
            \bigwedge_jV_{R_j}.
        \]
    \end{enumerate}
\end{proof}

\begin{corollary}[Fixed-point lattices]
\label{cor:to-fixed-point-lattices}
    Operation-saturated local theories form a complete lattice with
    \[
        \bigwedge\nolimits_{\mathrm{sat}}V_i
        =
        \bigwedge_iV_i
    \]
    and
    \[
        \bigvee\nolimits_{\mathrm{sat}}V_i
        =
        c\left(\bigvee_iV_i\right).
    \]
    Semantically closed operation congruences form a complete lattice with
    \[
        \bigwedge\nolimits_{\mathrm{cl}}R_j
        =
        \bigwedge_jR_j
    \]
    and
    \[
        \bigvee\nolimits_{\mathrm{cl}}R_j
        =
        d\left(\bigvee_jR_j\right).
    \]
\end{corollary}

\begin{proof}
    \leavevmode
    \begin{enumerate}
        \item \textbf{Use the general fixed-point construction.}
        The fixed points of any closure operator on a complete lattice are
        closed under arbitrary ambient meets.

        \item \textbf{Compute fixed-point joins.}
        The join of fixed points is obtained by taking the ambient join and
        applying the closure operator. Applying this general fact to \(c\) and
        \(d\) gives the displayed formulas.
    \end{enumerate}
\end{proof}

\subsection{The local algebraic Eilenberg theorem}
\label{subsec:to-local-eilenberg-theorem}

\begin{theorem}[Local algebraic Eilenberg duality]
\label{thm:to-local-algebraic-eilenberg}
    The assignments
    \[
        V\longmapsto R_V,
        \qquad
        R\longmapsto V_R
    \]
    restrict to an order-reversing isomorphism of complete lattices between:
    \begin{enumerate}
        \item operation-saturated local coequational theories over
        \((\mathbb A,\mathbb B)\);
        \item semantically closed identity-on-objects operation congruences on
        \(\mathbb A^{\mathrm{op}}\).
    \end{enumerate}
\end{theorem}

\begin{proof}
    \leavevmode
    \begin{enumerate}
        \item \textbf{Send saturated theories to closed congruences.}
        If \(V\) is operation-saturated, then
        \[
            V=V_{R_V}.
        \]
        By the Galois triangle identity,
        \[
            R_{V_{R_V}}=R_V,
        \]
        so \(R_V\) is semantically closed.

        \item \textbf{Send closed congruences to saturated theories.}
        If \(R\) is semantically closed, then
        \[
            R=R_{V_R}.
        \]
        Again by the Galois triangle identity,
        \[
            V_{R_{V_R}}=V_R,
        \]
        so \(V_R\) is operation-saturated.

        \item \textbf{Verify the two composites.}
        On fixed points,
        \[
            V\longmapsto R_V\longmapsto V_{R_V}=V
        \]
        and
        \[
            R\longmapsto V_R\longmapsto R_{V_R}=R.
        \]
        Thus the restrictions are mutually inverse.

        \item \textbf{Verify order reversal and completeness.}
        Both assignments are antitone. The complete-lattice structures are
        those of Corollary~\ref{cor:to-fixed-point-lattices}.
    \end{enumerate}
\end{proof}

The fixed-point duality is summarized by
\[
\begin{tikzcd}[column sep=huge]
    \mathsf{Loc}^{\mathrm{sat}}_{\mathbb A,\mathbb B}
        \arrow[r, shift left=1.2ex, "{V\mapsto R_V}"]
    &
    \mathsf{Cong}^{\mathrm{cl}}_{\mathrm{ioo}}
    (\mathbb A^{\mathrm{op}})^{\mathrm{op}}
        \arrow[l, shift left=1.2ex, "{R\mapsto V_R}"]
\end{tikzcd}
\]

\subsection{Subdirect operational representation}
\label{subsec:to-subdirect-operational-representation}

Let
\[
    \mathsf{Cat}_{A_0}^{\mathrm{ioo}}(\mathcal E)
\]
denote the category of internal categories with object object \(A_0\) and
identity-on-objects internal functors. Products in this category are formed
on arrow objects by products in
\[
    \mathcal E/(A_0\times A_0).
\]

\begin{theorem}[Subdirect product theorem]
\label{thm:to-subdirect-product}
    Let \((V_i)_i\) be a nonempty small family of local coequational theories.
    There is a canonical identity-on-objects monomorphism
    \[
        \operatorname{Op}\left(\bigvee_iV_i\right)
        \hookrightarrow
        \prod_i\operatorname{Op}(V_i)
    \]
    whose composites with the product projections are
    identity-on-objects regular epimorphisms. Thus
    \[
        \operatorname{Op}\left(\bigvee_iV_i\right)
    \]
    is a subdirect product of the categories
    \(\operatorname{Op}(V_i)\).
\end{theorem}

\begin{proof}
    \leavevmode
    \begin{enumerate}
        \item \textbf{Form the component quotient functors.}
        For every \(i\), write
        \[
            R_i:=R_{V_i}
        \]
        and let
        \[
            q_i:
            \mathbb A^{\mathrm{op}}
            \twoheadrightarrow
            \mathbb A^{\mathrm{op}}/R_i
        \]
        be the effective identity-on-objects quotient.

        By
        Theorem~\ref{thm:to-transition-output-quotient},
        \[
            \mathbb A^{\mathrm{op}}/R_i
            \cong
            \operatorname{Op}(V_i).
        \]

        \item \textbf{Form the product representation.}
        Products in
        \(\mathsf{Cat}_{A_0}^{\mathrm{ioo}}(\mathcal E)\)
        are formed on arrow objects by products in
        \(\mathcal E/(A_0\times A_0)\).
        Hence the quotient functors assemble into an
        identity-on-objects internal functor
        \[
            q
            :=
            \langle q_i\rangle_i:
            \mathbb A^{\mathrm{op}}
            \longrightarrow
            \prod_i
            \mathbb A^{\mathrm{op}}/R_i.
        \]

        \item \textbf{Compute the kernel congruence of the product map.}
        On arrow objects, equality after applying \(q_1\) is equality after
        applying every component \((q_i)_1\). Therefore, in the subobject
        lattice of
        \[
            \operatorname{Par}(\mathbb A),
        \]
        one has
        \[
        \begin{aligned}
            \ker(q_1)
            &=
            \bigwedge_i
            \ker\bigl((q_i)_1\bigr)
            \\
            &=
            \bigwedge_iR_i
            \\
            &=
            \bigwedge_iR_{V_i}.
        \end{aligned}
        \]
        By
        Theorem~\ref{thm:to-galois-join-formulas},
        \[
            \bigwedge_iR_{V_i}
            =
            R_{\bigvee_iV_i}.
        \]
        Consequently,
        \[
            \ker(q_1)
            =
            R_{\bigvee_iV_i}.
        \]

        \item \textbf{Take the identity-on-objects regular image.}
        Factor \(q\) on arrow objects as
        \[
        \begin{tikzcd}[column sep=large]
            \mathbb A^{\mathrm{op}}
                \arrow[r, two heads, "e"]
                \arrow[dr, "q"']
            &
            I
                \arrow[d, hook, "m"]
            \\
            &
            \displaystyle
            \prod_i
            \mathbb A^{\mathrm{op}}/R_i.
        \end{tikzcd}
        \]
        By
        Lemma~\ref{lem:to-create-ioo-regular-image},
        the arrow image carries a unique internal category structure for
        which \(e\) and \(m\) are identity-on-objects internal functors.

        Since \(m_1\) is monic,
        \[
            \ker(e_1)
            =
            \ker(q_1)
            =
            R_{\bigvee_iV_i}.
        \]
        Exactness therefore gives a canonical isomorphism
        \[
            I
            \cong
            \mathbb A^{\mathrm{op}}
            /
            R_{\bigvee_iV_i}.
        \]

        \item \textbf{Identify the regular image with the operational image.}
        Applying
        Theorem~\ref{thm:to-transition-output-quotient}
        to the local theory \(\bigvee_iV_i\) gives
        \[
            \mathbb A^{\mathrm{op}}
            /
            R_{\bigvee_iV_i}
            \cong
            \operatorname{Op}
            \left(
                \bigvee_iV_i
            \right).
        \]
        Transporting the monomorphism \(m\) across this canonical
        isomorphism gives
        \[
            \operatorname{Op}
            \left(
                \bigvee_iV_i
            \right)
            \hookrightarrow
            \prod_i
            \operatorname{Op}(V_i).
        \]

        \item \textbf{Prove that every projection is regular epic.}
        Let
        \[
            \pi_i:
            \prod_j
            \mathbb A^{\mathrm{op}}/R_j
            \to
            \mathbb A^{\mathrm{op}}/R_i
        \]
        be the \(i\)-th product projection, and put
        \[
            p_i:=\pi_i m:
            I\to
            \mathbb A^{\mathrm{op}}/R_i.
        \]
        Since
        \[
            p_i e
            =
            \pi_i m e
            =
            \pi_i q
            =
            q_i,
        \]
        the composite \(p_i e\) is regular epic on arrows.

        Factor the arrow map of \(p_i\) as
        \[
            (p_i)_1
            =
            n_i r_i,
        \]
        where \(r_i\) is regular epic and \(n_i\) is monic. Then
        \[
            (q_i)_1
            =
            n_i r_i e_1.
        \]
        The composite \(r_i e_1\) is regular epic, and
        \((q_i)_1\) is regular epic. Since a regular epimorphism is extremal,
        the monomorphism \(n_i\) through which \((q_i)_1\) factors is an
        isomorphism. Hence \((p_i)_1\) is a regular epimorphism.

        Transporting across the canonical quotient--image isomorphisms shows
        that every projection
        \[
            \operatorname{Op}
            \left(
                \bigvee_iV_i
            \right)
            \twoheadrightarrow
            \operatorname{Op}(V_i)
        \]
        is an identity-on-objects regular epimorphism.
    \end{enumerate}
\end{proof}

\begin{corollary}[Operational image of a generated theory]
\label{cor:to-generated-theory-subdirect}
    Let
    \[
        V=\bigvee_iV_i
    \]
    for a nonempty small family. Then every operation represented faithfully
    on \(V\) is determined by its family of represented operations on the
    \(V_i\), and every projection
    \[
        \operatorname{Op}(V)\twoheadrightarrow\operatorname{Op}(V_i)
    \]
    is a quotient.
\end{corollary}

\section{Specifications and syntactic transition--output categories}
\label{sec:to-syntactic-transition-output-categories}

\subsection{Generated derivative theories}
\label{subsec:to-generated-derivative-theories}

Let
\[
    k:K\to
    \mathbf{Beh}_{\mathbb A,\mathbb B,0}
\]
be a behaviour specification. Recall that
\[
    d_K:
    D_K\to
    \mathbf{Beh}_{\mathbb A,\mathbb B,0}
\]
is the residual evaluation map and that
\[
    \operatorname{Der}_0(K)
    \hookrightarrow
    \mathbf{Beh}_{\mathbb A,\mathbb B,0}
\]
is its regular image:
\[
\begin{tikzcd}[column sep=large]
    D_K
        \arrow[r, two heads, "\eta_K"]
        \arrow[dr, "d_K"']
    &
    \operatorname{Der}_0(K)
        \arrow[d, hook, "j_K"]
    \\
    &
    \mathbf{Beh}_{\mathbb A,\mathbb B,0}.
\end{tikzcd}
\]

Equivalently,
\[
    \operatorname{Der}_0(K)
    =
    \operatorname{Gen}_{\partial}
    \bigl(\operatorname{Im}(k)\bigr).
\]
Thus \(\operatorname{Der}_0(K)\) is the least local coequational theory
containing the specified behaviours.

\begin{proposition}[Formal derivative-orbit presentation]
\label{prop:to-formal-derivative-orbit-presentation}
    Let
    \[
        k:K\to
        \mathbf{Beh}_{\mathbb A,\mathbb B,0}
    \]
    be a behaviour specification, and write
    \[
        p_K
        :=
        p_{\mathbf{Beh}}k:
        K\to A_0.
    \]
    Define the formal derivative-orbit carrier
    \[
        D_K
        :=
        A_1
        \times_{t_A,A_0,p_K}
        K.
    \]
    A generalized element is written
    \[
        (r,\kappa),
        \qquad
        r:v\to p_K(\kappa).
    \]

    Give \(D_K\) the input anchor
    \[
        p_{D_K}(r,\kappa):=s_A(r)=v
    \]
    and output anchor
    \[
        q_{D_K}(r,\kappa)
        :=
        k(\kappa)(r),
    \]
    where \(k(\kappa)(r)\) denotes the object of \(\mathbb B\) assigned by the
    residual behaviour \(k(\kappa)\) to the object
    \[
        r:v\to p_K(\kappa)
    \]
    of the residual slice.

    Define the transition by
    \[
        \delta_{D_K}
        \bigl(
            a,(r,\kappa)
        \bigr)
        :=
        (r\circ a,\kappa)
    \]
    for
    \[
        a:u\to v,
        \qquad
        r:v\to p_K(\kappa),
    \]
    and define the output comparison by
    \[
        \omega_{D_K}
        \bigl(
            a,(r,\kappa)
        \bigr)
        :=
        k(\kappa)(r\circ a\to r).
    \]

    Then \(D_K\) is a Mealy morphism, and the residual evaluation
    \[
        d_K:
        D_K\to
        \mathbf{Beh}^{\mathrm{can}}_{\mathbb A,\mathbb B},
        \qquad
        d_K(r,\kappa)
        :=
        \partial_rk(\kappa),
    \]
    is a strict semantic homomorphism.

    Consequently, its regular-image factorization
    \[
    \begin{tikzcd}[column sep=large]
        D_K
            \arrow[r, two heads, "\eta_K"]
            \arrow[dr, "d_K"']
        &
        \operatorname{Der}_0(K)
            \arrow[d, hook, "j_K"]
        \\
        &
        \mathbf{Beh}^{\mathrm{can}}_{\mathbb A,\mathbb B}
    \end{tikzcd}
    \]
    is a strict regular-image factorization of Mealy morphisms.
\end{proposition}

\begin{proof}
    \leavevmode
    \begin{enumerate}
        \item \textbf{Verify the typing equations.}
        Let
        \[
            a:u\to v,
            \qquad
            r:v\to p_K(\kappa).
        \]
        Then
        \[
            r\circ a:u\to p_K(\kappa),
        \]
        so
        \[
            \delta_{D_K}
            \bigl(
                a,(r,\kappa)
            \bigr)
            =
            (r\circ a,\kappa)
        \]
        lies over \(u=s_A(a)\).

        Moreover,
        \[
            \omega_{D_K}
            \bigl(
                a,(r,\kappa)
            \bigr)
            =
            k(\kappa)(r\circ a\to r)
        \]
        is an arrow
        \[
            k(\kappa)(r\circ a)
            \to
            k(\kappa)(r),
        \]
        hence
        \[
            q_{D_K}
            \left(
                \delta_{D_K}
                \bigl(
                    a,(r,\kappa)
                \bigr)
            \right)
            \to
            q_{D_K}(r,\kappa).
        \]

        \item \textbf{Verify the unit equations.}
        For the identity arrow \(e_A(v)\),
        \[
        \begin{aligned}
            \delta_{D_K}
            \bigl(
                e_A(v),(r,\kappa)
            \bigr)
            &=
            (r\circ e_A(v),\kappa)
            \\
            &=
            (r,\kappa).
        \end{aligned}
        \]
        Since \(k(\kappa)\) is a functor on the residual slice,
        \[
            k(\kappa)(r\to r)
            =
            e_B(k(\kappa)(r)).
        \]
        Therefore
        \[
            \omega_{D_K}
            \bigl(
                e_A(v),(r,\kappa)
            \bigr)
            =
            e_B(q_{D_K}(r,\kappa)).
        \]

        \item \textbf{Verify transition multiplication.}
        Let
        \[
            a:u\to v,
            \qquad
            b:v\to w,
            \qquad
            r:w\to p_K(\kappa).
        \]
        Then
        \[
        \begin{aligned}
            \delta_{D_K}
            \bigl(
                b\circ a,(r,\kappa)
            \bigr)
            &=
            (r\circ(b\circ a),\kappa)
            \\
            &=
            ((r\circ b)\circ a,\kappa)
            \\
            &=
            \delta_{D_K}
            \left(
                a,
                \delta_{D_K}
                \bigl(
                    b,(r,\kappa)
                \bigr)
            \right).
        \end{aligned}
        \]

        \item \textbf{Verify output multiplication.}
        Functoriality of \(k(\kappa)\) on the residual slice gives
        \[
        \begin{aligned}
            &k(\kappa)
            \bigl(
                r\circ(b\circ a)\to r
            \bigr)
            \\
            &\qquad=
            k(\kappa)(r\circ b\to r)
            \circ
            k(\kappa)
            \bigl(
                (r\circ b)\circ a\to r\circ b
            \bigr).
        \end{aligned}
        \]
        Hence
        \[
        \begin{aligned}
            \omega_{D_K}
            \bigl(
                b\circ a,(r,\kappa)
            \bigr)
            &=
            \omega_{D_K}
            \bigl(
                b,(r,\kappa)
            \bigr)
            \\
            &\qquad\circ
            \omega_{D_K}
            \left(
                a,
                \delta_{D_K}
                \bigl(
                    b,(r,\kappa)
                \bigr)
            \right).
        \end{aligned}
        \]
        Thus \(D_K\) satisfies the Mealy multiplication law.

        \item \textbf{Verify strict preservation of transitions.}
        By the derivative composition law,
        \[
        \begin{aligned}
            d_K
            \left(
                \delta_{D_K}
                \bigl(
                    a,(r,\kappa)
                \bigr)
            \right)
            &=
            d_K(r\circ a,\kappa)
            \\
            &=
            \partial_{r\circ a}k(\kappa)
            \\
            &=
            \partial_a
            \bigl(
                \partial_rk(\kappa)
            \bigr)
            \\
            &=
            \delta_{\mathbf{Beh}}
            \bigl(
                a,d_K(r,\kappa)
            \bigr).
        \end{aligned}
        \]

        \item \textbf{Verify strict preservation of outputs.}
        The canonical output comparison produced by acting with \(a\) on
        \[
            \partial_rk(\kappa)
        \]
        is the value of \(k(\kappa)\) on the residual morphism
        \[
            r\circ a\to r.
        \]
        Therefore
        \[
            \omega_{\mathbf{Beh}}
            \bigl(
                a,d_K(r,\kappa)
            \bigr)
            =
            k(\kappa)(r\circ a\to r)
            =
            \omega_{D_K}
            \bigl(
                a,(r,\kappa)
            \bigr).
        \]
        Hence \(d_K\) is strict.

        \item \textbf{Pass to the regular image.}
        Regular images of strict Mealy homomorphisms are created on their
        carrier maps by the regular-image construction established in
        Chapter~\ref{ch:coequational-subobjects-birkhoff}. Therefore the
        carrier image of \(d_K\) inherits a unique Mealy structure for which
        \[
            \eta_K:
            D_K\twoheadrightarrow\operatorname{Der}_0(K)
        \]
        and
        \[
            j_K:
            \operatorname{Der}_0(K)
            \hookrightarrow
            \mathbf{Beh}^{\mathrm{can}}_{\mathbb A,\mathbb B}
        \]
        are strict.
    \end{enumerate}
\end{proof}

\begin{definition}[Syntactic operation congruence]
\label{def:to-syntactic-operation-congruence}
    The \textbf{syntactic operation congruence} of \(K\) is
    \[
        R_K^{\mathrm{syn}}
        :=
        R_{\operatorname{Der}_0(K)}.
    \]
\end{definition}

\begin{definition}[Syntactic transition--output category]
\label{def:to-syntactic-transition-output-category}
    The \textbf{syntactic transition--output category} of \(K\) is
    \[
        \operatorname{SynTO}_{\mathbb A,\mathbb B}(K)
        :=
        \operatorname{Op}_{\mathbb A,\mathbb B}
        \bigl(\operatorname{Der}_0(K)\bigr).
    \]
    We write
    \[
        q_K^{\mathrm{syn}}:
        \mathbb A^{\mathrm{op}}
        \twoheadrightarrow
        \operatorname{SynTO}_{\mathbb A,\mathbb B}(K)
    \]
    for the regular-epimorphic part of the image factorization of
    \[
        \chi_{\operatorname{Der}_0(K)}.
    \]
\end{definition}

By Theorem~\ref{thm:to-transition-output-quotient},
\[
    \operatorname{SynTO}_{\mathbb A,\mathbb B}(K)
    \cong
    \mathbb A^{\mathrm{op}}/
    R_K^{\mathrm{syn}},
\]
and under this isomorphism \(q_K^{\mathrm{syn}}\) corresponds to
\(q_{R_K^{\mathrm{syn}}}\). The defining diagram is
\[
\begin{tikzcd}[column sep=large, row sep=large]
    \mathbb A^{\mathrm{op}}
        \arrow[rr,
        "\chi_{\operatorname{Der}_0(K)}"]
        \arrow[dr, two heads, "q_K^{\mathrm{syn}}"']
    &&
    \operatorname{KlEnd}_{\mathbb B}
    \bigl(\operatorname{Der}_0(K)\bigr)
    \\
    &
    \operatorname{SynTO}_{\mathbb A,\mathbb B}(K)
        \arrow[ur, hook]
    &
\end{tikzcd}
\]

\subsection{Regular presentations of operational equations}
\label{subsec:to-regular-presentations}

\begin{theorem}[Regular-presentation detection]
\label{thm:to-regular-presentation-detection}
    Let
    \[
        e:G\twoheadrightarrow V
    \]
    be a strict regular epimorphism of Mealy morphisms, with \(V\) a local
    coequational theory. Let
    \[
        U\to
        \operatorname{Par}(\mathbb A)
    \]
    classify a pair of parallel input-arrow families
    \[
        a,a':U\to A_1.
    \]
    Pulling \(e\) back to the corresponding source-state families gives a
    regular epimorphism
    \[
        e_U:
        G_{\mathrm{src},U}
        \twoheadrightarrow
        V_{\mathrm{src},U}.
    \]

    Then the operations induced by \(a\) and \(a'\) on \(V\) are equal if and
    only if their evaluated transition and output components agree after
    precomposition with \(e_U\).
\end{theorem}

\begin{proof}
    \leavevmode
    \begin{enumerate}
        \item \textbf{Prove the forward implication.}
        Equality of the two operations on \(V\) implies equality after
        precomposition with any map, in particular with
        \[
            e_U:
            G_{\mathrm{src},U}
            \twoheadrightarrow
            V_{\mathrm{src},U}.
        \]

        \item \textbf{Represent the two operation families.}
        The parallel input-arrow families determine maps
        \[
            U\rightrightarrows
            \operatorname{KlEnd}_{\mathbb B}(V)_1.
        \]
        Their evaluations fit into
        \[
        \begin{tikzcd}[column sep=large]
            G_{\mathrm{src},U}
                \arrow[r, two heads, "e_U"]
                \arrow[dr, shift left=.7ex]
                \arrow[dr, shift right=.7ex]
            &
            V_{\mathrm{src},U}
                \arrow[d, shift left=.7ex]
                \arrow[d, shift right=.7ex]
            \\
            &
            E_{[2],U}.
        \end{tikzcd}
        \]

        \item \textbf{Apply epimorphic extensionality.}
        If the two evaluated maps agree after precomposition with \(e_U\),
        Lemma~\ref{lem:to-regular-epi-extensionality} implies equality of the
        two maps
        \[
            U\rightrightarrows
            \operatorname{KlEnd}_{\mathbb B}(V)_1.
        \]
        Hence the two operations on \(V\) are equal.
    \end{enumerate}
\end{proof}

\subsection{Context expansion}
\label{subsec:to-context-expansion}

Reasoning in the internal language, let
\[
    \kappa\in K
\]
and let
\[
    r:v\to p_K(\kappa)
\]
be a residual context. For
\[
    a:u\to v,
\]
the composite
\[
    r\circ a:u\to p_K(\kappa)
\]
is an object of the residual slice. We write
\[
    k(\kappa)(a,r)
    :=
    k(\kappa)(r\circ a\to r)
\]
for the output comparison assigned to the corresponding residual morphism:
\[
\begin{tikzcd}[column sep=large]
    r\circ a
        \arrow[r]
    &
    r
        \arrow[r]
    &
    1_{p_K(\kappa)}.
\end{tikzcd}
\]

\begin{theorem}[Context expansion]
\label{thm:to-context-expansion}
    Let
    \[
        a,a':u\to v
    \]
    be parallel arrows of \(\mathbb A\). Then
    \[
        a\,R_K^{\mathrm{syn}}\,a'
    \]
    if and only if, internally, for every
    \[
        \kappa\in K
    \]
    and every residual context
    \[
        r:v\to p_K(\kappa),
    \]
    one has
    \[
        \partial_{r\circ a}k(\kappa)
        =
        \partial_{r\circ a'}k(\kappa)
    \]
    and
    \[
        k(\kappa)(a,r)
        =
        k(\kappa)(a',r).
    \]
\end{theorem}

\begin{proof}
    \leavevmode
    \begin{enumerate}
        \item \textbf{Use the regular presentation of the derivative theory.}
        The derivative theory is presented by the strict regular epimorphism
        \[
            D_K\twoheadrightarrow\operatorname{Der}_0(K).
        \]

        \item \textbf{Restrict to the relevant source fibre.}
        Pulling the presentation back to the fibre over \(v\) gives a
        regular-epimorphic cover of
        \[
            \operatorname{Der}_0(K)_v
        \]
        by generated residual states
        \[
            \partial_rk(\kappa).
        \]
        The pullback is
        \[
        \begin{tikzcd}
            (D_K)_v
                \arrow[r, two heads]
                \arrow[d]
                \arrow[dr, phantom, "\lrcorner", very near start]
            &
            \operatorname{Der}_0(K)_v
                \arrow[d]
            \\
            D_K
                \arrow[r, two heads]
            &
            \operatorname{Der}_0(K).
        \end{tikzcd}
        \]

        \item \textbf{Evaluate the action of \(a\).}
        Acting by \(a\) on the generated residual state gives
        \[
            \partial_a(\partial_rk(\kappa))
            =
            \partial_{r\circ a}k(\kappa),
        \]
        and the output comparison is
        \[
            \omega_{\mathbf{Beh}}
            \bigl(a,\partial_rk(\kappa)\bigr)
            =
            k(\kappa)(r\circ a\to r).
        \]

        \item \textbf{Evaluate the action of \(a'\).}
        Similarly,
        \[
            \partial_{a'}(\partial_rk(\kappa))
            =
            \partial_{r\circ a'}k(\kappa)
        \]
        and
        \[
            \omega_{\mathbf{Beh}}
            \bigl(a',\partial_rk(\kappa)\bigr)
            =
            k(\kappa)(r\circ a'\to r).
        \]

        \item \textbf{Apply regular-presentation detection.}
        Theorem~\ref{thm:to-regular-presentation-detection} says that the
        operations induced by \(a\) and \(a'\) agree on
        \(\operatorname{Der}_0(K)\) exactly when the displayed transition and
        output components agree on this regular-epimorphic presentation.
        This is exactly the asserted contextual condition.
    \end{enumerate}
\end{proof}

\subsection{The syntactic universal property}
\label{subsec:to-syntactic-universal-property}

\begin{definition}[Recognition of a specification]
\label{def:to-recognition-specification}
    An operation quotient
    \[
        q_R:
        \mathbb A^{\mathrm{op}}
        \twoheadrightarrow
        \mathbb A^{\mathrm{op}}/R
    \]
    \textbf{recognizes} \(K\) if it recognizes the generated local theory
    \[
        \operatorname{Der}_0(K).
    \]
\end{definition}

\begin{theorem}[Syntactic transition--output universal property]
\label{thm:to-syntactic-universal-property}
    Let
    \[
        q_R:
        \mathbb A^{\mathrm{op}}
        \twoheadrightarrow
        \mathbb A^{\mathrm{op}}/R
    \]
    be an operation quotient. The following are equivalent:
    \begin{enumerate}
        \item \(q_R\) recognizes \(K\);
        \item
        \[
            R\leq R_K^{\mathrm{syn}};
        \]
        \item there is a unique identity-on-objects regular epimorphism
        \[
            h:
            \mathbb A^{\mathrm{op}}/R
            \twoheadrightarrow
            \operatorname{SynTO}_{\mathbb A,\mathbb B}(K)
        \]
        satisfying
        \[
            hq_R=q_K^{\mathrm{syn}}.
        \]
    \end{enumerate}

    Moreover,
    \[
        \operatorname{SynTO}_{\mathbb A,\mathbb B}(K)
    \]
    faithfully recognizes \(\operatorname{Der}_0(K)\).
\end{theorem}

\begin{proof}
    \leavevmode
    \begin{enumerate}
        \item \textbf{Apply the recognition criterion.}
        Proposition~\ref{prop:to-recognition-criterion}, applied to
        \[
            V=\operatorname{Der}_0(K),
        \]
        gives
        \[
            q_R\text{ recognizes }K
            \quad\Longleftrightarrow\quad
            R\leq R_{\operatorname{Der}_0(K)}
            =
            R_K^{\mathrm{syn}}.
        \]

        \item \textbf{Apply universal faithful recognition.}
        Theorem~\ref{thm:to-universal-faithful-recognition} gives a unique
        regular-epimorphic comparison
        \[
        \begin{tikzcd}[column sep=large, row sep=large]
            \mathbb A^{\mathrm{op}}
                \arrow[r, two heads, "q_R"]
                \arrow[dr, two heads,
                "q_K^{\mathrm{syn}}"']
            &
            \mathbb A^{\mathrm{op}}/R
                \arrow[d, dashed, two heads, "h"]
            \\
            &
            \operatorname{SynTO}_{\mathbb A,\mathbb B}(K).
        \end{tikzcd}
        \]

        \item \textbf{Verify faithfulness of the syntactic quotient.}
        The quotient by
        \[
            R_K^{\mathrm{syn}}
            =
            \ker
            \left(
                \chi_{\operatorname{Der}_0(K),1}
            \right)
        \]
        is the regular image of the residual representation. Its induced map
        into
        \[
            \operatorname{KlEnd}_{\mathbb B}
            \bigl(\operatorname{Der}_0(K)\bigr)
        \]
        is therefore monic on arrows.
    \end{enumerate}
\end{proof}

\begin{corollary}[Coarsest recognizing quotient]
\label{cor:to-coarsest-recognizing-quotient}
    The syntactic transition--output category is the coarsest
    identity-on-objects quotient of \(\mathbb A^{\mathrm{op}}\) recognizing
    \(K\), and the unique such quotient whose induced action on the generated
    derivative theory is faithful.

    Equivalently, it is terminal among recognizing quotients ordered by
    factorization: every recognizing quotient admits a unique
    identity-on-objects regular-epimorphic comparison onto
    \(\operatorname{SynTO}_{\mathbb A,\mathbb B}(K)\).
\end{corollary}

\subsection{Two-stage operational Myhill--Nerode theory}
\label{subsec:to-two-stage-operational-myhill-nerode}

\begin{definition}[Raw action image]
\label{def:to-raw-action-image}
    For a Mealy morphism \(P\), define
    \[
        \operatorname{ActIm}(P)
        :=
        \operatorname{RegIm}
        \left(
            \chi_P:
            \mathbb A^{\mathrm{op}}
            \to
            \operatorname{KlEnd}_{\mathbb B}(P)
        \right).
    \]
    Define also
    \[
        \operatorname{Op}(P)
        :=
        \operatorname{Op}(\operatorname{Beh}(P)).
    \]
\end{definition}

The raw action image records the action on the original state object:
\[
\begin{tikzcd}[column sep=large, row sep=large]
    \mathbb A^{\mathrm{op}}
        \arrow[rr, "\chi_P"]
        \arrow[dr, two heads]
    &&
    \operatorname{KlEnd}_{\mathbb B}(P)
    \\
    &
    \operatorname{ActIm}(P)
        \arrow[ur, hook]
    &
\end{tikzcd}
\]
The semantic operation image records the action after states with identical
residual behaviour have been identified.

\begin{theorem}[Two-stage operational quotient]
\label{thm:to-two-stage-operational-quotient}
    Let \(P\) be a Mealy morphism realizing a specification \(K\). Then there
    is a canonical chain of identity-on-objects regular epimorphisms
    \[
        \operatorname{ActIm}(P)
        \twoheadrightarrow
        \operatorname{Op}(P)
        \twoheadrightarrow
        \operatorname{SynTO}_{\mathbb A,\mathbb B}(K).
    \]
\end{theorem}

\begin{proof}
    \leavevmode
    \begin{enumerate}
        \item \textbf{Define the raw action congruence.}
        Put
        \[
            R_P^{\mathrm{act}}
            :=
            \ker(\chi_{P,1}).
        \]
        By Lemma~\ref{lem:to-create-ioo-regular-image},
        \[
            \operatorname{ActIm}(P)
            \cong
            \mathbb A^{\mathrm{op}}/R_P^{\mathrm{act}}.
        \]

        \item \textbf{Descend operational equality through state minimization.}
        The semantic quotient
        \[
            \eta_P:
            P\twoheadrightarrow\operatorname{Beh}(P)
        \]
        is strict and regular epic. By
        Proposition~\ref{prop:to-operational-descent},
        \[
            R_P^{\mathrm{act}}
            \leq
            R_{\operatorname{Beh}(P)}.
        \]

        \item \textbf{Form the first quotient comparison.}
        The preceding inclusion gives
        \[
        \begin{tikzcd}[column sep=large]
            \operatorname{ActIm}(P)
                \arrow[r, two heads]
            &
            \operatorname{Op}(P).
        \end{tikzcd}
        \]

        \item \textbf{Use realization of the specification.}
        Since \(P\) realizes \(K\),
        \[
            \operatorname{Der}_0(K)
            \leq
            \operatorname{Beh}(P).
        \]
        Contravariance of valid equations gives
        \[
            R_{\operatorname{Beh}(P)}
            \leq
            R_{\operatorname{Der}_0(K)}.
        \]

        \item \textbf{Form the second quotient comparison.}
        Passing to effective quotients gives
        \[
        \begin{tikzcd}[column sep=large]
            \operatorname{Op}(P)
                \arrow[r, two heads]
            &
            \operatorname{SynTO}_{\mathbb A,\mathbb B}(K).
        \end{tikzcd}
        \]

        \item \textbf{Assemble the two stages.}
        The complete chain is
        \[
        \begin{tikzcd}[column sep=large]
            \mathbb A^{\mathrm{op}}
                \arrow[r, two heads]
            &
            \operatorname{ActIm}(P)
                \arrow[r, two heads]
            &
            \operatorname{Op}(P)
                \arrow[r, two heads]
            &
            \operatorname{SynTO}(K).
        \end{tikzcd}
        \]
    \end{enumerate}
\end{proof}

\begin{corollary}[Separated realizations]
\label{cor:to-separated-action-image}
    If \(P\) is behaviourally separated, then
    \[
        \operatorname{ActIm}(P)
        \cong
        \operatorname{Op}(P).
    \]
\end{corollary}

\begin{proof}
    \leavevmode
    \begin{enumerate}
        \item \textbf{Use monicity of the semantic map.}
        Behavioural separatedness means that
        \[
            \beta_P:P\to\mathbf{Beh}^{\mathrm{can}}
        \]
        is monic.

        \item \textbf{Inspect the semantic image factorization.}
        Write
        \[
        \begin{tikzcd}[column sep=large]
            P
                \arrow[r, two heads, "\eta_P"]
                \arrow[dr, "\beta_P"']
            &
            \operatorname{Beh}(P)
                \arrow[d, hook, "m_P"]
            \\
            &
            \mathbf{Beh}^{\mathrm{can}}.
        \end{tikzcd}
        \]
        The composite
        \[
            m_P\eta_P=\beta_P
        \]
        is monic.

        \item \textbf{Prove that \(\eta_P\) is monic.}
        Suppose
        \[
            \eta_Pg=\eta_Ph.
        \]
        Then
        \[
        \begin{aligned}
            \beta_Pg
            &=
            m_P\eta_Pg
            \\
            &=
            m_P\eta_Ph
            \\
            &=
            \beta_Ph.
        \end{aligned}
        \]
        Since \(\beta_P\) is monic,
        \[
            g=h.
        \]
        Thus \(\eta_P\) is monic.

        \item \textbf{Conclude that \(\eta_P\) is an isomorphism.}
        Since \(\eta_P\) is both regular epic and monic, it is an isomorphism.

        \item \textbf{Transport the operation representation.}
        Transporting the transition--output representation across this strict
        isomorphism identifies
        \[
            \operatorname{ActIm}(P)
            \cong
            \operatorname{Op}(\operatorname{Beh}(P))
            =
            \operatorname{Op}(P).
        \]
    \end{enumerate}
\end{proof}

\begin{corollary}[Reachable realizations]
\label{cor:to-reachable-operation-image}
    If \(P\) is reachable as a realization of \(K\), then
    \[
        \operatorname{Op}(P)
        \cong
        \operatorname{SynTO}_{\mathbb A,\mathbb B}(K).
    \]
\end{corollary}

\begin{proof}
    \leavevmode
    \begin{enumerate}
        \item \textbf{Use the reachability quotient.}
        Reachability gives a regular epimorphism
        \[
            \rho_{P,K}:D_K\twoheadrightarrow P
        \]
        satisfying
        \[
            \beta_P\rho_{P,K}=d_K:
        \]
        \[
        \begin{tikzcd}[column sep=large]
            D_K
                \arrow[r, two heads, "\rho_{P,K}"]
                \arrow[dr, "d_K"']
            &
            P
                \arrow[d, "\beta_P"]
            \\
            &
            \mathbf{Beh}_0.
        \end{tikzcd}
        \]

        \item \textbf{Compare regular images.}
        Precomposition with a regular epimorphism does not change a regular
        image. Hence
        \[
        \begin{aligned}
            \operatorname{Im}(\beta_P)
            &=
            \operatorname{Im}(\beta_P\rho_{P,K})
            \\
            &=
            \operatorname{Im}(d_K).
        \end{aligned}
        \]

        \item \textbf{Identify the behavioural theories.}
        Therefore
        \[
            \operatorname{Beh}(P)
            \cong
            \operatorname{Der}_0(K)
        \]
        as strict subobjects of the canonical behaviour object.

        \item \textbf{Identify the operation images.}
        The strict isomorphism identifies their transition--output
        representations and hence their regular images:
        \[
            \operatorname{Op}(P)
            \cong
            \operatorname{SynTO}(K).
        \]
    \end{enumerate}
\end{proof}

\begin{corollary}[Operational Myhill--Nerode theorem]
\label{cor:to-operational-myhill-nerode}
    If \(P\) is reachable and behaviourally separated as a realization of
    \(K\), then
    \[
        \operatorname{ActIm}(P)
        \cong
        \operatorname{SynTO}_{\mathbb A,\mathbb B}(K).
    \]
\end{corollary}

\begin{proof}
    \leavevmode
    \begin{enumerate}
        \item \textbf{Use separatedness.}
        Corollary~\ref{cor:to-separated-action-image} gives
        \[
            \operatorname{ActIm}(P)
            \cong
            \operatorname{Op}(P).
        \]

        \item \textbf{Use reachability.}
        Corollary~\ref{cor:to-reachable-operation-image} gives
        \[
            \operatorname{Op}(P)
            \cong
            \operatorname{SynTO}(K).
        \]

        \item \textbf{Compose the canonical isomorphisms.}
        Combining the two isomorphisms proves the result.
    \end{enumerate}
\end{proof}

\subsection{A finiteness criterion}
\label{subsec:to-finiteness-criterion}

\begin{proposition}[Finiteness of syntactic operation categories]
\label{prop:to-finiteness-syntactic-category}
    Work in \(\mathbf{Set}\). Let \(V\) be a local coequational theory. Suppose:
    \begin{enumerate}
        \item \(A_0\) is finite;
        \item every fibre \(V_v\) is finite;
        \item every hom-set of \(\mathbb B\) between output objects occurring in
        \(V\) is finite.
    \end{enumerate}
    Then \(\operatorname{Op}(V)\) is a finite category.

    In particular, if \(\operatorname{Der}_0(K)\) has finite fibres and the
    relevant output hom-sets are finite, then
    \[
        \operatorname{SynTO}(K)
    \]
    is finite.
\end{proposition}

\begin{proof}
    \leavevmode
    \begin{enumerate}
        \item \textbf{Count state transformations.}
        A Kleisli operation
        \[
            V_v\rightsquigarrow V_u
        \]
        has an underlying function
        \[
            V_v\to V_u.
        \]
        Since both sets are finite, there are finitely many such functions.

        \item \textbf{Count output comparisons.}
        For every \(x\in V_v\), the operation also selects an arrow
        \[
            q(Fx)\to q(x)
        \]
        of \(\mathbb B\). The relevant hom-sets are finite and \(V_v\) is
        finite, so there are finitely many possible comparison families.

        \item \textbf{Count Kleisli hom-sets.}
        Hence every hom-set of
        \[
            \operatorname{KlEnd}_{\mathbb B}(V)
        \]
        is finite.

        \item \textbf{Count the operational image.}
        The object set \(A_0\) is finite, and
        \(\operatorname{Op}(V)\) is an internal subcategory of
        \(\operatorname{KlEnd}_{\mathbb B}(V)\). Therefore
        \(\operatorname{Op}(V)\) is finite.

        \item \textbf{Apply the result to a specification.}
        Taking
        \[
            V=\operatorname{Der}_0(K)
        \]
        gives the final statement.
    \end{enumerate}
\end{proof}

\section{Structured output and automatic algebraic closure}
\label{sec:to-structured-output}

\subsection{Pointwise output operations}
\label{subsec:to-pointwise-output-operations}

Let
\[
    \theta:\mathbb B^n\to\mathbb B
\]
be an internal functor. For Mealy morphisms
\[
    P_1,\dots,P_n:
    \mathbb A\rightsquigarrow\mathbb B,
\]
write
\[
    \theta(P_1,\dots,P_n)
\]
for the carrier
\[
    P_1\times_{A_0}\cdots\times_{A_0}P_n
\]
with input anchor the common \(A_0\)-anchor and output anchor
\[
    \theta_0\circ
    \langle q_1,\dots,q_n\rangle.
\]
Its transition is componentwise, and its output comparison is obtained by
applying \(\theta_1\) componentwise.

The carrier map is
\[
\begin{tikzcd}[column sep=large]
    P_1\times_{A_0}\cdots\times_{A_0}P_n
        \arrow[r, "{\langle p,q_\theta\rangle}"]
    &
    A_0\times B_0,
\end{tikzcd}
\]
where
\[
    q_\theta
    :=
    \theta_0\langle q_1,\dots,q_n\rangle.
\]

Pointwise application defines a map of residual behaviour objects
\[
    \theta_{\mathbf{Beh}}:
    \underbrace{
    \mathbf{Beh}_0
    \times_{A_0}\cdots\times_{A_0}
    \mathbf{Beh}_0
    }_{n\text{ factors}}
    \to
    \mathbf{Beh}_0.
\]
Its input anchor is unchanged, and its output anchor is obtained by applying
\(\theta_0\) to the \(n\) output anchors. The domain is regarded over
\(A_0\times B_0\) using
\[
    1_{A_0}\times\theta_0:
    A_0\times B_0^n\to A_0\times B_0.
\]

For \(n=0\), the functor
\[
    \theta:\mathbb B^0\to\mathbb B
\]
selects a constant output object and its identity comparison. The associated
carrier is \(A_0\) with the induced map to \(A_0\times B_0\).

\begin{proposition}[Pointwise operations intertwine residual actions]
\label{prop:to-pointwise-intertwining}
    Let
    \[
        H_1,\dots,H_n
    \]
    be residual behaviours over the same input object. For every input arrow
    \(a\),
    \[
        \partial_a
        \theta_{\mathbf{Beh}}(H_1,\dots,H_n)
        =
        \theta_{\mathbf{Beh}}
        (\partial_aH_1,\dots,\partial_aH_n),
    \]
    and
    \[
        \omega_{\mathbf{Beh}}
        \bigl(
            a,
            \theta_{\mathbf{Beh}}(H_1,\dots,H_n)
        \bigr)
        =
        \theta_1
        \bigl(
            \omega_{\mathbf{Beh}}(a,H_1),
            \dots,
            \omega_{\mathbf{Beh}}(a,H_n)
        \bigr).
    \]
    Consequently,
    \[
        \theta_{\mathbf{Beh}}:
        \theta
        \bigl(
            \mathbf{Beh}^{\mathrm{can}},
            \dots,
            \mathbf{Beh}^{\mathrm{can}}
        \bigr)
        \to
        \mathbf{Beh}^{\mathrm{can}}
    \]
    is a strict semantic homomorphism.
\end{proposition}

\begin{proof}
    \leavevmode
    \begin{enumerate}
        \item \textbf{Use pointwise residual functors.}
        Residual behaviours are internal functors on residual slices.
        Pointwise application of \(\theta\) produces the composite
        \[
        \begin{tikzcd}[column sep=large]
            \mathbb A/v
                \arrow[r, "{\langle H_1,\dots,H_n\rangle}"]
            &
            \mathbb B^n
                \arrow[r, "\theta"]
            &
            \mathbb B.
        \end{tikzcd}
        \]

        \item \textbf{Commute with residual precomposition.}
        Derivation by \(a\) is precomposition with the residual-slice functor
        induced by \(a\). Associativity of functor composition gives
        \[
            \partial_a
            \theta_{\mathbf{Beh}}(H_1,\dots,H_n)
            =
            \theta_{\mathbf{Beh}}
            (\partial_aH_1,\dots,\partial_aH_n).
        \]
        Diagrammatically,
        \[
        \begin{tikzcd}[column sep=large, row sep=large]
            \mathbb A/u
                \arrow[r]
                \arrow[d, "{\langle\partial_aH_i\rangle_i}"']
            &
            \mathbb A/v
                \arrow[d, "{\langle H_i\rangle_i}"]
            \\
            \mathbb B^n
                \arrow[r, equal]
            &
            \mathbb B^n
                \arrow[r, "\theta"]
            &
            \mathbb B.
        \end{tikzcd}
        \]

        \item \textbf{Compute the canonical output comparison.}
        The canonical output is the value on the residual morphism
        \[
            a\to1_v.
        \]
        Applying \(\theta\) to the tuple of assigned arrows gives
        \[
            \theta_1
            \bigl(
                \omega_{\mathbf{Beh}}(a,H_1),
                \dots,
                \omega_{\mathbf{Beh}}(a,H_n)
            \bigr).
        \]
        The output diagram is
        \[
        \begin{tikzcd}[column sep=large]
            \theta_0(q(\partial_aH_i))_i
                \arrow[r,
                "{\theta_1(\omega_{\mathbf{Beh}}(a,H_i))_i}"]
            &
            \theta_0(q(H_i))_i.
        \end{tikzcd}
        \]

        \item \textbf{Conclude strictness.}
        The derivative and output equations are exactly the strict
        transition- and output-preservation equations for
        \(\theta_{\mathbf{Beh}}\).

        \item \textbf{Handle the nullary case.}
        For \(n=0\), the pointwise behaviour is constant on the object
        selected by \(\theta\), every derivative remains constant, and every
        residual morphism is sent to the selected identity arrow. The same
        strictness equations therefore hold.
    \end{enumerate}
\end{proof}

\subsection{Recognized theories are automatically structured}
\label{subsec:to-recognized-theories-structured}

\begin{theorem}[Automatic structured closure]
\label{thm:to-automatic-structured-closure}
    Let \(R\) be an operation congruence. For every internal functor
    \[
        \theta:\mathbb B^n\to\mathbb B,
    \]
    pointwise application restricts to a strict map
    \[
        \theta_{V_R}:
        \theta(V_R,\dots,V_R)
        \to
        V_R.
    \]
    Thus \(V_R\) is closed under every pointwise finitary operation induced by
    an internal functor on \(\mathbb B\).
\end{theorem}

\begin{proof}
    \leavevmode
    \begin{enumerate}
        \item \textbf{Choose behaviours in the safety core.}
        Let
        \[
            H_1,\dots,H_n\in V_R
        \]
        and let \(r\) be an arbitrary residual context. Then
        \[
            \partial_rH_i\in\operatorname{Eq}(R)
        \]
        for every \(i\).

        \item \textbf{Apply an operational equation componentwise.}
        If
        \[
            a\,R\,a',
        \]
        then for every \(i\),
        \[
            \partial_a(\partial_rH_i)
            =
            \partial_{a'}(\partial_rH_i)
        \]
        and
        \[
            \omega_{\mathbf{Beh}}(a,\partial_rH_i)
            =
            \omega_{\mathbf{Beh}}(a',\partial_rH_i).
        \]

        \item \textbf{Apply the pointwise output operation.}
        Proposition~\ref{prop:to-pointwise-intertwining} gives
        \[
        \begin{aligned}
            \partial_a
            \theta_{\mathbf{Beh}}(\partial_rH_i)_i
            &=
            \theta_{\mathbf{Beh}}
            (\partial_a\partial_rH_i)_i
            \\
            &=
            \theta_{\mathbf{Beh}}
            (\partial_{a'}\partial_rH_i)_i
            \\
            &=
            \partial_{a'}
            \theta_{\mathbf{Beh}}(\partial_rH_i)_i.
        \end{aligned}
        \]
        It also gives equality of the corresponding output comparisons after
        applying \(\theta_1\).

        \item \textbf{Use derivative compatibility.}
        Since derivatives commute with pointwise \(\theta\),
        \[
            \partial_r
            \theta_{\mathbf{Beh}}(H_1,\dots,H_n)
            =
            \theta_{\mathbf{Beh}}
            (\partial_rH_1,\dots,\partial_rH_n).
        \]
        Hence every derivative of
        \[
            \theta_{\mathbf{Beh}}(H_1,\dots,H_n)
        \]
        lies in \(\operatorname{Eq}(R)\).

        \item \textbf{Conclude membership in the safety core.}
        Therefore
        \[
            \theta_{\mathbf{Beh}}(H_1,\dots,H_n)
            \in V_R.
        \]
        The factorization is
        \[
        \begin{tikzcd}[column sep=large]
            \theta(V_R,\dots,V_R)
                \arrow[r, "\theta_{\mathbf{Beh}}"]
                \arrow[dr, dashed, "\theta_{V_R}"']
            &
            \mathbf{Beh}^{\mathrm{can}}
            \\
            &
            V_R
                \arrow[u, hook].
        \end{tikzcd}
        \]

        \item \textbf{Prove strictness.}
        The factorization is the restriction of the strict semantic
        homomorphism \(\theta_{\mathbf{Beh}}\), so
        \(\theta_{V_R}\) is strict.

        \item \textbf{Handle nullary operations.}
        For \(n=0\), the pointwise behaviour is constant. Its derivatives are
        equal to itself, and every \(R\)-related pair acts by the same
        identity comparison. Hence the constant behaviour lies in \(V_R\),
        giving a strict map
        \[
            A_0\to V_R.
        \]
    \end{enumerate}
\end{proof}

\begin{corollary}[Saturated theories inherit output structure]
\label{cor:to-saturated-theories-structured}
    If \(V\) is operation-saturated, then \(V\) is closed under every
    pointwise operation induced by an internal functor
    \[
        \mathbb B^n\to\mathbb B.
    \]
    In particular, if \(\mathbb B\) carries the operations of a finite-product
    theory \(\mathbb T\), then every operation-saturated local theory is
    automatically a local \(\mathbb T\)-coequational theory.
\end{corollary}

\begin{proof}
    \leavevmode
    \begin{enumerate}
        \item \textbf{Represent the saturated theory by its equations.}
        Operation saturation gives
        \[
            V=V_{R_V}.
        \]

        \item \textbf{Apply automatic structured closure.}
        Theorem~\ref{thm:to-automatic-structured-closure}, applied to
        \(R_V\), gives closure of \(V_{R_V}=V\) under every pointwise
        operation induced by \(\mathbb B\).
    \end{enumerate}
\end{proof}

\begin{remark}[Pointwise structure versus synchronized machine lifting]
\label{rem:to-pointwise-versus-synchronized-structure}
    The pointwise operations on the canonical behaviour object and their
    restrictions to \(V_R\) should be distinguished from the synchronized
    lifting of an output operation to arbitrary Mealy morphisms.

    On the canonical behaviour object, an internal functor
    \[
        \theta:\mathbb B^n\to\mathbb B
    \]
    induces an actual pointwise operation
    \[
        \theta_{\mathbf{Beh}}:
        \mathbf{Beh}_0
        \times_{A_0}\cdots\times_{A_0}
        \mathbf{Beh}_0
        \to
        \mathbf{Beh}_0.
    \]
    When \(V_R\) is invariant under this operation, the factorization
    \[
        \theta_{V_R}:
        V_R
        \times_{A_0}\cdots\times_{A_0}
        V_R
        \to
        V_R
    \]
    is a genuine operation on the behavioural subobject. If the chosen
    operations on \(\mathbb B\) interpret a finite-product theory
    \(\mathbb T\), then the induced pointwise operations on
    \(\mathbf{Beh}_0\) and on \(V_R\) satisfy the equations of
    \(\mathbb T\), subject to the standing coherence convention for chosen
    finite products.

    By contrast, for arbitrary Mealy morphisms
    \[
        P_1,\dots,P_n:
        \mathbb A\rightsquigarrow\mathbb B,
    \]
    the synchronized lifting
    \[
        \theta(P_1,\dots,P_n)
    \]
    has carrier
    \[
        P_1\times_{A_0}\cdots\times_{A_0}P_n.
    \]
    Even when \(\theta\) is a projection operation on \(\mathbb B\), the
    lifted machine retains this full synchronized product carrier rather than
    becoming literally one of the factors \(P_i\).

    Consequently, the synchronized lifting does not by itself make the
    entire category of Mealy morphisms into a strict algebra for
    \(\mathbb T\). The genuine algebraic conclusion used here concerns the
    pointwise behavioural object and the invariant behavioural subobjects
    \(V_R\), not a strict finite-product-theory structure on the whole Mealy
    category.
\end{remark}

\section{Globalization as an indexed operational polarity}
\label{sec:to-global-indexed-polarity}

\subsection{Residual reindexing}
\label{subsec:to-residual-reindexing}

Let
\[
    \mathsf{Inp}
\]
be a chosen category of internal input categories and internal functors. Fix
the output category \(\mathbb B\).

Let
\[
    F:\mathbb A'\to\mathbb A
\]
be an internal functor. Residual reindexing is the map
\[
    F^\sharp:
    A'_0
    \times_{F_0,A_0,p_{\mathbf{Beh}}}
    \mathbf{Beh}_{\mathbb A,\mathbb B,0}
    \to
    \mathbf{Beh}_{\mathbb A',\mathbb B,0}
\]
given by
\[
    F^\sharp(v',H)
    =
    H\circ F_{/v'}.
\]
The residual-slice comparison is
\[
\begin{tikzcd}[column sep=large]
    \mathbb A'/v'
        \arrow[r, "F_{/v'}"]
    &
    \mathbb A/F_0(v')
        \arrow[r, "H"]
    &
    \mathbb B.
\end{tikzcd}
\]

Write
\[
    \overline F
    :=
    F_0\times1_{B_0}:
    A'_0\times B_0
    \to
    A_0\times B_0.
\]

\begin{proposition}[Ambient computation of joins of local theories]
\label{prop:to-local-theory-ambient-joins}
    Let
    \[
        (V_i)_i
    \]
    be a small family of local coequational theories in
    \[
        \mathbf{Beh}^{\mathrm{can}}_{\mathbb A,\mathbb B}.
    \]
    Their join in
    \[
        \mathsf{Loc}_{\mathbb A,\mathbb B}
    \]
    is computed as their join in the ambient subobject lattice
    \[
        \operatorname{Sub}_{\mathcal E/Z}
        \bigl(
            \mathbf{Beh}_{\mathbb A,\mathbb B,0}
        \bigr).
    \]
    In particular, the ambient join
    \[
        V:=\bigvee_iV_i
    \]
    is derivative-closed.
\end{proposition}

\begin{proof}
    Let
    \[
        \delta_{\mathbf{Beh}}:
        A_1
        \times_{t_A,A_0,p_{\mathbf{Beh}}}
        \mathbf{Beh}_0
        \to
        \mathbf{Beh}_0
    \]
    be the derivative transition of the canonical behaviour Mealy morphism.

    For a subobject
    \[
        C\hookrightarrow\mathbf{Beh}_0,
    \]
    write
    \[
        \mathsf D(C)
        :=
        A_1
        \times_{t_A,A_0,p_{\mathbf{Beh}}}
        C
    \]
    for its pulled-back action domain, regarded as a subobject of
    \[
        A_1
        \times_{t_A,A_0,p_{\mathbf{Beh}}}
        \mathbf{Beh}_0.
    \]

    Pullback on subobjects preserves small joins in a topos. Therefore
    \[
    \begin{aligned}
        \mathsf D(V)
        &=
        \mathsf D\left(\bigvee_iV_i\right)
        \\
        &=
        \bigvee_i\mathsf D(V_i).
    \end{aligned}
    \]

    Let
    \[
        \exists_{\delta_{\mathbf{Beh}}}
    \]
    denote the direct-image operation on subobjects along
    \(\delta_{\mathbf{Beh}}\). Since it is a left adjoint, it preserves small
    joins. Hence
    \[
    \begin{aligned}
        \exists_{\delta_{\mathbf{Beh}}}
        \mathsf D(V)
        &=
        \exists_{\delta_{\mathbf{Beh}}}
        \left(
            \bigvee_i\mathsf D(V_i)
        \right)
        \\
        &=
        \bigvee_i
        \exists_{\delta_{\mathbf{Beh}}}
        \mathsf D(V_i).
    \end{aligned}
    \]

    Each \(V_i\) is derivative-closed, so
    \[
        \exists_{\delta_{\mathbf{Beh}}}
        \mathsf D(V_i)
        \leq
        V_i.
    \]
    Consequently,
    \[
    \begin{aligned}
        \exists_{\delta_{\mathbf{Beh}}}
        \mathsf D(V)
        &\leq
        \bigvee_iV_i
        \\
        &=
        V.
    \end{aligned}
    \]
    This inequality says precisely that
    \(\delta_{\mathbf{Beh}}\) restricts to a transition
    \[
        A_1
        \times_{A_0}
        V
        \to
        V.
    \]
    Therefore \(V\) is derivative-closed and hence is a local coequational
    theory.

    Since \(V\) is the least ambient subobject containing all \(V_i\), it is
    also the least local theory containing all \(V_i\). Thus it is their join
    in
    \(\mathsf{Loc}_{\mathbb A,\mathbb B}\).
\end{proof}

\begin{definition}[Direct reindexing of local theories]
\label{def:to-Fbullet}
    For a local theory
    \[
        V\hookrightarrow
        \mathbf{Beh}_{\mathbb A,\mathbb B,0},
    \]
    define
    \[
        F^\bullet V
        :=
        \operatorname{Im}
        \left(
            \overline F^*V
            \xrightarrow{F^\sharp|_V}
            \mathbf{Beh}_{\mathbb A',\mathbb B,0}
        \right).
    \]
    Equivalently,
    \[
        F^\bullet V
        =
        \operatorname{Beh}(F^*V)
    \]
    as a subobject of
    \(\mathbf{Beh}_{\mathbb A',\mathbb B,0}\).
\end{definition}

The defining regular image is
\[
\begin{tikzcd}[column sep=large]
    \overline F^*V
        \arrow[r, two heads]
        \arrow[dr, "F^\sharp|_V"']
    &
    F^\bullet V
        \arrow[d, hook]
    \\
    &
    \mathbf{Beh}_{\mathbb A',\mathbb B,0}.
\end{tikzcd}
\]

\begin{proposition}[Basic properties of \(F^\bullet\)]
\label{prop:to-Fbullet-properties}
    The assignment
    \[
        F^\bullet:
        \mathsf{Loc}_{\mathbb A,\mathbb B}
        \to
        \mathsf{Loc}_{\mathbb A',\mathbb B}
    \]
    is monotone and preserves small joins.

    For the identity input functor there is a canonical equality in the
    subobject lattice
    \[
        (1_{\mathbb A})^\bullet V=V.
    \]
    For composable input functors
    \[
        \mathbb A''
        \xrightarrow{G}
        \mathbb A'
        \xrightarrow{F}
        \mathbb A,
    \]
    there is a canonical coherent isomorphism of chosen representatives
    \[
        G^\bullet(F^\bullet V)
        \cong
        (F\circ G)^\bullet V.
    \]
\end{proposition}

\begin{proof}
    \leavevmode
    \begin{enumerate}
        \item \textbf{Identify \(F^\bullet V\) as a semantic image.}
        The restricted residual-reindexing map
        \[
            F^\sharp|_V:
            \overline F^*V
            \to
            \mathbf{Beh}_{\mathbb A',\mathbb B,0}
        \]
        is the strict semantic map of the input-pullback Mealy morphism
        \(F^*V\). Therefore its regular image
        \[
            F^\bullet V
            =
            \operatorname{Beh}(F^*V)
        \]
        is a local coequational theory.

        \item \textbf{Express \(F^\bullet\) in subobject calculus.}
        On ambient subobject lattices,
        \[
            F^\bullet
            =
            \exists_{F^\sharp}
            \circ
            \overline F^{\,*}.
        \]
        Here
        \[
            \overline F^{\,*}
        \]
        is pullback along
        \[
            \overline F
            =
            F_0\times1_{B_0},
        \]
        and
        \[
            \exists_{F^\sharp}
        \]
        is the direct-image operation along the residual-reindexing map.

        Both operations are monotone, so \(F^\bullet\) is monotone.

        \item \textbf{Prove preservation of small joins.}
        Let \((V_i)_i\) be a small family of local theories. By
        Proposition~\ref{prop:to-local-theory-ambient-joins}, their join in
        \(\mathsf{Loc}_{\mathbb A,\mathbb B}\) is their ambient subobject
        join:
        \[
            \bigvee\nolimits_{\mathsf{Loc}}V_i
            =
            \bigvee\nolimits_{\operatorname{Sub}}V_i.
        \]

        Pullback on subobjects in a topos preserves small joins, and
        \(\exists_{F^\sharp}\) preserves small joins because it is a left
        adjoint. Hence
        \[
        \begin{aligned}
            F^\bullet
            \left(
                \bigvee_iV_i
            \right)
            &=
            \exists_{F^\sharp}
            \left(
                \overline F^{\,*}
                \left(
                    \bigvee_iV_i
                \right)
            \right)
            \\
            &=
            \exists_{F^\sharp}
            \left(
                \bigvee_i
                \overline F^{\,*}V_i
            \right)
            \\
            &=
            \bigvee_i
            \exists_{F^\sharp}
            \left(
                \overline F^{\,*}V_i
            \right)
            \\
            &=
            \bigvee_iF^\bullet V_i.
        \end{aligned}
        \]
        The join on the right is again the ambient join and therefore, by
        Proposition~\ref{prop:to-local-theory-ambient-joins}, the join in
        \(\mathsf{Loc}_{\mathbb A',\mathbb B}\).

        \item \textbf{Prove the identity law.}
        For the identity functor \(1_{\mathbb A}\), the residual-slice
        comparison is canonically the identity. Hence
        \[
            (1_{\mathbb A})^\sharp
        \]
        is canonically the identity map on the behaviour object. Its
        restriction to \(V\) has image \(V\), so
        \[
            (1_{\mathbb A})^\bullet V=V
        \]
        in the subobject lattice.

        \item \textbf{Factor the first-stage semantic map.}
        Factor the semantic map of \(F^*V\) as
        \[
        \begin{tikzcd}[column sep=large]
            F^*V
                \arrow[r, two heads, "e_F"]
                \arrow[dr]
            &
            F^\bullet V
                \arrow[d, hook]
            \\
            &
            \mathbf{Beh}_{\mathbb A',\mathbb B}.
        \end{tikzcd}
        \]
        The map \(e_F\) is a strict regular epimorphism of Mealy morphisms.

        \item \textbf{Pull the first-stage factorization back along \(G\).}
        Pullback stability of regular epimorphisms gives a strict regular
        epimorphism
        \[
            G^*e_F:
            G^*F^*V
            \twoheadrightarrow
            G^*(F^\bullet V).
        \]

        \item \textbf{Compare the second-stage semantic images.}
        The semantic map of
        \[
            G^*(F^\bullet V)
        \]
        precomposed with \(G^*e_F\) is the semantic map of \(G^*F^*V\).
        Precomposition with a regular epimorphism does not change a regular
        image. Therefore
        \[
            \operatorname{Beh}
            \bigl(
                G^*(F^\bullet V)
            \bigr)
            \cong
            \operatorname{Beh}
            \bigl(
                G^*F^*V
            \bigr).
        \]
        By definition, the left-hand side is
        \[
            G^\bullet(F^\bullet V).
        \]

        \item \textbf{Use pullback coherence.}
        The chosen pullbacks give a canonical coherent strict isomorphism of
        input-pulled-back Mealy morphisms
        \[
            G^*F^*V
            \cong
            (F\circ G)^*V.
        \]
        Strictly isomorphic Mealy morphisms have canonically isomorphic
        semantic images. Hence
        \[
        \begin{aligned}
            G^\bullet(F^\bullet V)
            &=
            \operatorname{Beh}
            \bigl(
                G^*(F^\bullet V)
            \bigr)
            \\
            &\cong
            \operatorname{Beh}
            \bigl(
                G^*F^*V
            \bigr)
            \\
            &\cong
            \operatorname{Beh}
            \bigl(
                (F\circ G)^*V
            \bigr)
            \\
            &=
            (F\circ G)^\bullet V.
        \end{aligned}
        \]

        The coherence of these isomorphisms follows from the coherence of
        iterated chosen pullbacks and the uniqueness of regular-image
        factorizations up to canonical isomorphism.
    \end{enumerate}
\end{proof}

\subsection{Operational inverse image}
\label{subsec:to-operational-inverse-image}

\begin{definition}[Operational inverse image]
\label{def:to-Fflat}
    Let
    \[
        R\in
        \mathsf{Cong}_{\mathrm{ioo}}(\mathbb A^{\mathrm{op}}).
    \]
    Define its \textbf{operational inverse image} along \(F\) by
    \[
        F^\flat R
        :=
        R_{F^\bullet(V_R)}.
    \]
\end{definition}

Thus
\[
    (\mathbb A')^{\mathrm{op}}/F^\flat R
\]
is the faithful syntactic quotient of the reindexed theory
\[
    F^\bullet(V_R):
\]
\[
\begin{tikzcd}[column sep=large]
    (\mathbb A')^{\mathrm{op}}
        \arrow[r, two heads]
    &
    (\mathbb A')^{\mathrm{op}}/F^\flat R
        \arrow[r, hook]
    &
    \operatorname{KlEnd}_{\mathbb B}
    \bigl(F^\bullet(V_R)\bigr).
\end{tikzcd}
\]

\begin{proposition}[Monotonicity of \(F^\flat\)]
\label{prop:to-Fflat-monotone}
    If
    \[
        R\leq S,
    \]
    then
    \[
        F^\flat R
        \leq
        F^\flat S.
    \]
\end{proposition}

\begin{proof}
    \leavevmode
    \begin{enumerate}
        \item \textbf{Reverse the congruence inclusion.}
        Since \(R\leq S\), antitonicity gives
        \[
            V_S\leq V_R.
        \]

        \item \textbf{Apply direct reindexing.}
        Monotonicity of \(F^\bullet\) gives
        \[
            F^\bullet V_S
            \leq
            F^\bullet V_R.
        \]

        \item \textbf{Apply valid-equation contravariance.}
        Applying the antitone map \(R_{(-)}\) gives
        \[
            R_{F^\bullet V_R}
            \leq
            R_{F^\bullet V_S}.
        \]
        By definition, this is
        \[
            F^\flat R
            \leq
            F^\flat S.
        \]
    \end{enumerate}
\end{proof}

\subsection{Raw pullback and semantic completion}
\label{subsec:to-raw-pullback-semantic-completion}

Let
\[
    R\hookrightarrow\operatorname{Par}(\mathbb A)
\]
be an operation congruence. The arrow map \(F_1\) induces
\[
    F_{\operatorname{Par}}:
    \operatorname{Par}(\mathbb A')
    \to
    \operatorname{Par}(\mathbb A),
    \qquad
    (a',b')
    \longmapsto
    (F_1a',F_1b').
\]
The defining diagram is
\[
\begin{tikzcd}
    \operatorname{Par}(\mathbb A')
        \arrow[r, "F_{\operatorname{Par}}"]
        \arrow[d]
    &
    \operatorname{Par}(\mathbb A)
        \arrow[d]
    \\
    A'_1\times A'_1
        \arrow[r, "F_1\times F_1"']
    &
    A_1\times A_1.
\end{tikzcd}
\]

Define the literal inverse-image relation
\[
    F^{-1}R
    :=
    F_{\operatorname{Par}}^*R.
\]
Equivalently, \(F^{-1}R\) is the kernel pair of
\[
    A'_1
    \longrightarrow
    (A'_0\times A'_0)
    \times_{A_0\times A_0}
    (A_1/R)
\]
induced by \(F_1\) and the quotient map \(A_1\to A_1/R\).

\begin{proposition}[Raw pullback is contained in operational inverse image]
\label{prop:to-raw-pullback-contained}
    The relation \(F^{-1}R\) is an identity-on-objects operation congruence on
    \((\mathbb A')^{\mathrm{op}}\), and
    \[
        F^{-1}R
        \leq
        F^\flat R.
    \]
    Consequently,
    \[
        d(F^{-1}R)
        \leq
        F^\flat R.
    \]
\end{proposition}

\begin{proof}
    \leavevmode
    \begin{enumerate}
        \item \textbf{Pull back the effective congruence.}
        The relation \(F^{-1}R\) is the pullback
        \[
        \begin{tikzcd}
            F^{-1}R
                \arrow[r]
                \arrow[d, hook]
                \arrow[dr, phantom, "\lrcorner", very near start]
            &
            R
                \arrow[d, hook]
            \\
            \operatorname{Par}(\mathbb A')
                \arrow[r, "F_{\operatorname{Par}}"']
            &
            \operatorname{Par}(\mathbb A).
        \end{tikzcd}
        \]
        Effectiveness of equivalence relations is stable under pullback in a
        topos.

        \item \textbf{Verify operational compatibility.}
        Compatibility with identities and composition follows from
        functoriality of \(F\). Hence \(F^{-1}R\) is an operation congruence.

        \item \textbf{Choose a raw pulled-back equation.}
        Let
        \[
            a',b':u'\to v'
        \]
        satisfy
        \[
            F_1(a')\,R\,F_1(b').
        \]

        \item \textbf{Use satisfaction of \(R\) by \(V_R\).}
        Since \(V_R\) satisfies \(R\), the arrows \(F_1(a')\) and \(F_1(b')\)
        induce the same operation on \(V_R\).

        \item \textbf{Pull the operation equality back along \(F\).}
        Their input-pulled-back operations on \(F^*V_R\) are equal. The
        equality is represented by
        \[
        \begin{tikzcd}[column sep=large]
            F^*V_R
                \arrow[r, shift left=.7ex, dashed, "a'"]
                \arrow[r, shift right=.7ex, dashed, "b'"']
            &
            F^*V_R.
        \end{tikzcd}
        \]

        \item \textbf{Descend through the semantic image.}
        The semantic image factorization gives a strict regular epimorphism
        \[
            F^*V_R
            \twoheadrightarrow
            F^\bullet(V_R).
        \]
        Proposition~\ref{prop:to-operational-descent} therefore implies that
        \(a'\) and \(b'\) induce the same operation on
        \(F^\bullet(V_R)\).

        \item \textbf{Conclude the first inclusion.}
        Thus
        \[
            a'\,F^\flat R\,b',
        \]
        and hence
        \[
            F^{-1}R\leq F^\flat R.
        \]

        \item \textbf{Apply semantic closure.}
        The congruence \(F^\flat R\) is of the form \(R_W\), so it is
        semantically closed. Monotonicity of \(d\) gives
        \[
        \begin{aligned}
            d(F^{-1}R)
            &\leq
            d(F^\flat R)
            \\
            &=
            F^\flat R.
        \end{aligned}
        \]
    \end{enumerate}
\end{proof}

Thus
\[
    F^{-1}R
    \leq
    d(F^{-1}R)
    \leq
    F^\flat R.
\]
The inequality may be strict because the reindexed semantic theory
\[
    F^\bullet(V_R)
\]
can be strictly smaller than the largest local theory
\[
    V_{F^{-1}R}
\]
satisfying the literal pullback congruence. Passing to the valid-equation
congruence of the smaller theory can therefore create additional operational
equations.

\subsection{Lax functoriality}
\label{subsec:to-Fflat-lax-functoriality}

\begin{proposition}[Lax composition of operational inverse image]
\label{prop:to-Fflat-lax-composition}
    For every operation congruence \(R\),
    \[
        (1_{\mathbb A})^\flat R
        =
        R_{V_R}
        =
        d(R).
    \]
    For composable input functors
    \[
        \mathbb A''
        \xrightarrow{G}
        \mathbb A'
        \xrightarrow{F}
        \mathbb A,
    \]
    one has
    \[
        G^\flat(F^\flat R)
        \leq
        (F\circ G)^\flat R.
    \]
    In particular, \(F^\flat\) is strictly unital on semantically closed
    congruences.
\end{proposition}

\begin{proof}
    \leavevmode
    \begin{enumerate}
        \item \textbf{Compute the identity case.}
        Since
        \[
            (1_{\mathbb A})^\bullet V_R=V_R,
        \]
        one has
        \[
            (1_{\mathbb A})^\flat R
            =
            R_{V_R}
            =
            d(R).
        \]

        \item \textbf{Introduce the first reindexed theory.}
        Put
        \[
            W:=F^\bullet(V_R).
        \]
        Then
        \[
            F^\flat R=R_W.
        \]

        \item \textbf{Apply inflationary saturation.}
        One has
        \[
            W\leq V_{R_W}.
        \]

        \item \textbf{Reindex along \(G\).}
        Monotonicity gives
        \[
            G^\bullet W
            \leq
            G^\bullet(V_{R_W}).
        \]

        \item \textbf{Reverse the inclusion by valid equations.}
        Contravariance of \(R_{(-)}\) gives
        \[
            R_{G^\bullet(V_{R_W})}
            \leq
            R_{G^\bullet W}.
        \]
        The left-hand side is
        \[
            G^\flat(F^\flat R).
        \]

        \item \textbf{Use composition coherence of \(F^\bullet\).}
        The right-hand side is canonically identified with
        \[
            R_{(F\circ G)^\bullet(V_R)}
            =
            (F\circ G)^\flat R.
        \]
        Thus
        \[
            G^\flat(F^\flat R)
            \leq
            (F\circ G)^\flat R.
        \]

        \item \textbf{Restrict to closed congruences.}
        If \(R=d(R)\), then
        \[
            (1_{\mathbb A})^\flat R=R,
        \]
        proving strict unitality on semantically closed congruences.
    \end{enumerate}
\end{proof}

\subsection{Saturation-adjusted reindexing}
\label{subsec:to-saturation-adjusted-reindexing}

\begin{definition}[Saturated reindexing]
\label{def:to-Fstar}
    For a local theory \(V\), define
    \[
        F^\star V
        :=
        V_{R_{F^\bullet V}}
        =
        c(F^\bullet V).
    \]
\end{definition}

Thus \(F^\star V\) is the least operation-saturated local theory containing
the raw reindexing \(F^\bullet V\):
\[
\begin{tikzcd}[column sep=large]
    F^\bullet V
        \arrow[r, hook]
    &
    F^\star V
        \arrow[r, hook]
    &
    \mathbf{Beh}_{\mathbb A',\mathbb B,0}.
\end{tikzcd}
\]

\begin{proposition}[Compatibility with local duality]
\label{prop:to-Fstar-Fflat-compatibility}
    If \(V\) is operation-saturated, then
    \[
        R_{F^\star V}
        =
        F^\flat(R_V).
    \]
    For every operation congruence \(R\),
    \[
        V_{F^\flat R}
        =
        F^\star(V_R).
    \]
\end{proposition}

\begin{proof}
    \leavevmode
    \begin{enumerate}
        \item \textbf{Assume \(V\) is saturated.}
        Since
        \[
            V=V_{R_V},
        \]
        one has
        \[
        \begin{aligned}
            F^\flat(R_V)
            &=
            R_{F^\bullet(V_{R_V})}
            \\
            &=
            R_{F^\bullet V}.
        \end{aligned}
        \]

        \item \textbf{Apply the triangle identity.}
        By definition,
        \[
            F^\star V
            =
            V_{R_{F^\bullet V}}.
        \]
        Hence
        \[
        \begin{aligned}
            R_{F^\star V}
            &=
            R_{V_{R_{F^\bullet V}}}
            \\
            &=
            R_{F^\bullet V}
            \\
            &=
            F^\flat(R_V).
        \end{aligned}
        \]

        \item \textbf{Compute the theory of an inverse image.}
        By definition,
        \[
        \begin{aligned}
            V_{F^\flat R}
            &=
            V_{R_{F^\bullet(V_R)}}
            \\
            &=
            c(F^\bullet(V_R))
            \\
            &=
            F^\star(V_R).
        \end{aligned}
        \]
    \end{enumerate}
\end{proof}

\subsection{Global theories and global congruence systems}
\label{subsec:to-global-theories-congruences}

\begin{definition}[Global coequational theory]
\label{def:to-global-coequational-theory}
    A \textbf{global coequational theory} is a family
    \[
        V=(V_{\mathbb A})_{\mathbb A\in\mathsf{Inp}}
    \]
    where
    \[
        V_{\mathbb A}
        \in
        \mathsf{Loc}_{\mathbb A,\mathbb B},
    \]
    such that for every
    \[
        F:\mathbb A'\to\mathbb A
    \]
    one has
    \[
        F^\bullet V_{\mathbb A}
        \leq
        V_{\mathbb A'}.
    \]
\end{definition}

The globality condition is represented by
\[
\begin{tikzcd}[column sep=large]
    \overline F^*V_{\mathbb A}
        \arrow[r, "F^\sharp|_{V_{\mathbb A}}"]
        \arrow[dr]
    &
    \mathbf{Beh}_{\mathbb A',\mathbb B,0}
    \\
    &
    V_{\mathbb A'}
        \arrow[u, hook].
\end{tikzcd}
\]

If every \(V_{\mathbb A'}\) is operation-saturated, the globality condition is
equivalent to
\[
    F^\star V_{\mathbb A}
    \leq
    V_{\mathbb A'}.
\]

\begin{definition}[Global operation congruence system]
\label{def:to-global-congruence-system}
    A \textbf{global operation congruence system} is a family
    \[
        R=(R_{\mathbb A})_{\mathbb A\in\mathsf{Inp}}
    \]
    with
    \[
        R_{\mathbb A}
        \in
        \mathsf{Cong}_{\mathrm{ioo}}
        (\mathbb A^{\mathrm{op}})
    \]
    such that, for every
    \[
        F:\mathbb A'\to\mathbb A,
    \]
    one has
    \[
        R_{\mathbb A'}
        \leq
        F^\flat R_{\mathbb A}.
    \]
\end{definition}

The lax-composition inequality of
Proposition~\ref{prop:to-Fflat-lax-composition} makes this condition coherent
under composite input functors. Indeed, if
\[
    R_{\mathbb A''}
    \leq
    G^\flat R_{\mathbb A'}
\]
and
\[
    R_{\mathbb A'}
    \leq
    F^\flat R_{\mathbb A},
\]
then monotonicity of \(G^\flat\) and lax composition give
\[
\begin{tikzcd}[column sep=large]
    R_{\mathbb A''}
        \arrow[r, hook]
    &
    G^\flat R_{\mathbb A'}
        \arrow[r, hook]
    &
    G^\flat(F^\flat R_{\mathbb A})
        \arrow[r, hook]
    &
    (F\circ G)^\flat R_{\mathbb A}.
\end{tikzcd}
\]

A global theory is operation-saturated when every component is
operation-saturated. A global congruence system is semantically closed when
every component is semantically closed.

\begin{proposition}[Saturated global theories induce closed congruence systems]
\label{prop:to-global-V-to-R}
    Let \(V\) be an operation-saturated global coequational theory. Then
    \[
        R^V_{\mathbb A}
        :=
        R_{V_{\mathbb A}}
    \]
    defines a semantically closed global operation congruence system.
\end{proposition}

\begin{proof}
    \leavevmode
    \begin{enumerate}
        \item \textbf{Prove local semantic closure.}
        Each \(R_{V_{\mathbb A}}\) is semantically closed by
        Corollary~\ref{cor:to-galois-triangle-identities}.

        \item \textbf{Apply globality of \(V\).}
        Let
        \[
            F:\mathbb A'\to\mathbb A.
        \]
        Globality gives
        \[
            F^\bullet V_{\mathbb A}
            \leq
            V_{\mathbb A'}.
        \]

        \item \textbf{Reverse the inclusion by valid equations.}
        Contravariance gives
        \[
            R_{V_{\mathbb A'}}
            \leq
            R_{F^\bullet V_{\mathbb A}}.
        \]

        \item \textbf{Use saturation of the source component.}
        Since
        \[
            V_{\mathbb A}
            =
            V_{R_{V_{\mathbb A}}},
        \]
        one has
        \[
            R_{F^\bullet V_{\mathbb A}}
            =
            F^\flat R_{V_{\mathbb A}}.
        \]

        \item \textbf{Conclude the global congruence inequality.}
        Therefore
        \[
        \begin{aligned}
            R^V_{\mathbb A'}
            &=
            R_{V_{\mathbb A'}}
            \\
            &\leq
            R_{F^\bullet V_{\mathbb A}}
            \\
            &=
            F^\flat R_{V_{\mathbb A}}
            \\
            &=
            F^\flat R^V_{\mathbb A}.
        \end{aligned}
        \]
        Thus \(R^V\) is a global operation congruence system.
    \end{enumerate}
\end{proof}

\begin{proposition}[Congruence systems induce saturated global theories]
\label{prop:to-global-R-to-V}
    Let \(R\) be a global operation congruence system. Then
    \[
        V^R_{\mathbb A}
        :=
        V_{R_{\mathbb A}}
    \]
    defines an operation-saturated global coequational theory.
\end{proposition}

\begin{proof}
    \leavevmode
    \begin{enumerate}
        \item \textbf{Prove local saturation.}
        Each theory
        \[
            V_{R_{\mathbb A}}
        \]
        is operation-saturated by
        Corollary~\ref{cor:to-galois-triangle-identities}.

        \item \textbf{Apply globality of \(R\).}
        For
        \[
            F:\mathbb A'\to\mathbb A,
        \]
        one has
        \[
            R_{\mathbb A'}
            \leq
            F^\flat R_{\mathbb A}.
        \]

        \item \textbf{Expand the operational inverse image.}
        By definition,
        \[
            F^\flat R_{\mathbb A}
            =
            R_{F^\bullet(V_{R_{\mathbb A}})}.
        \]

        \item \textbf{Apply operational soundness and completeness.}
        The inclusion
        \[
            R_{\mathbb A'}
            \leq
            R_{F^\bullet(V_{R_{\mathbb A}})}
        \]
        is equivalent to
        \[
            F^\bullet(V_{R_{\mathbb A}})
            \leq
            V_{R_{\mathbb A'}}.
        \]
        This is precisely the globality condition for \(V^R\).
    \end{enumerate}
\end{proof}

\begin{theorem}[Global algebraic Eilenberg duality]
\label{thm:to-global-algebraic-eilenberg}
    The assignments
    \[
        V
        \longmapsto
        R^V,
        \qquad
        R
        \longmapsto
        V^R
    \]
    restrict to an order-reversing isomorphism between:
    \begin{enumerate}
        \item operation-saturated global coequational theories;
        \item semantically closed global operation congruence systems.
    \end{enumerate}
\end{theorem}

\begin{proof}
    \leavevmode
    \begin{enumerate}
        \item \textbf{Verify well-definedness from theories to congruences.}
        Proposition~\ref{prop:to-global-V-to-R} shows that a saturated global
        theory determines a semantically closed global congruence system.

        \item \textbf{Verify well-definedness from congruences to theories.}
        Proposition~\ref{prop:to-global-R-to-V} shows that a closed global
        congruence system determines a saturated global theory.

        \item \textbf{Verify the first composite pointwise.}
        For every input category \(\mathbb A\),
        \[
            V^{R^V}_{\mathbb A}
            =
            V_{R_{V_{\mathbb A}}}
            =
            V_{\mathbb A},
        \]
        because \(V_{\mathbb A}\) is saturated.

        \item \textbf{Verify the second composite pointwise.}
        Similarly,
        \[
            R^{V^R}_{\mathbb A}
            =
            R_{V_{R_{\mathbb A}}}
            =
            R_{\mathbb A},
        \]
        because \(R_{\mathbb A}\) is semantically closed.

        \item \textbf{Conclude order reversal.}
        Both componentwise assignments are antitone, so the resulting global
        correspondence is an order-reversing isomorphism.
    \end{enumerate}
\end{proof}

\section{One-object and stateless specializations}
\label{sec:to-one-object-stateless}

\subsection{One-object input categories}
\label{subsec:to-one-object-input}

Let \(\mathbb A\) be the one-object category determined by a monoid object
\(M\). With the standing convention
\[
    b\circ a=ba,
\]
a Mealy transition is a right action
\[
    x\cdot m:=\delta(m,x)
\]
satisfying
\[
    x\cdot(mn)
    =
    (x\cdot m)\cdot n.
\]

Suppose also that \(\mathbb B\) is one-object, determined by a monoid object
\(N\). The output comparison is a cocycle
\[
    \omega:M\times P\to N
\]
satisfying
\[
    \omega(1,x)=1
\]
and
\[
    \omega(mn,x)
    =
    \omega(m,x)
    \omega(n,x\cdot m).
\]
The cocycle law is the one-object form of
\[
\begin{tikzcd}[column sep=large]
    q\bigl(x\cdot(mn)\bigr)
        \arrow[r, "{\omega(n,x\cdot m)}"]
        \arrow[rr, bend left=20, "{\omega(mn,x)}"']
    &
    q(x\cdot m)
        \arrow[r, "{\omega(m,x)}"]
    &
    q(x).
\end{tikzcd}
\]

The transition--output representation is the monoid homomorphism
\[
    \chi_P:
    M^{\mathrm{op}}
    \to
    \operatorname{KlEnd}_{\mathbb B}(P)
\]
given by
\[
    m\longmapsto
    \bigl(
        x\mapsto x\cdot m,
        \omega(m,x)
    \bigr).
\]

\subsection{Arrow monoids and execution monoids}
\label{subsec:to-arrow-execution-monoids}

For a one-object category \(\mathbb C\), let
\[
    \operatorname{ArrMon}(\mathbb C)
\]
denote its ordinary arrow monoid, whose multiplication is categorical
composition:
\[
    x\cdot_{\operatorname{Arr}}y
    :=
    x\circ_{\mathbb C}y.
\]

Define the \textbf{execution monoid} by
\[
    \operatorname{ExecMon}(\mathbb C)
    :=
    \operatorname{ArrMon}(\mathbb C)^{\mathrm{op}}.
\]
Equivalently, its multiplication is
\[
    x\star y
    :=
    y\circ_{\mathbb C}x.
\]

For the one-object category \(\mathbb A_M\) determined by \(M\) under the
standing convention,
\[
    \operatorname{ArrMon}(\mathbb A_M^{\mathrm{op}})
    \cong
    M^{\mathrm{op}},
\]
whereas
\[
    \operatorname{ExecMon}(\mathbb A_M^{\mathrm{op}})
    \cong
    M.
\]
The execution monoid is therefore the monoid recording input operations in
their temporal execution order.

A two-sided congruence on \(M^{\mathrm{op}}\) is the same underlying
equivalence relation as a two-sided congruence on \(M\). Thus operation
congruences may be regarded as ordinary two-sided monoid congruences when no
multiplication-order ambiguity is present.

\subsection{Residual cocycles}
\label{subsec:to-residual-cocycles}

A residual behaviour is an \(N\)-valued cocycle
\[
    \lambda:M\times M\to N
\]
satisfying
\[
    \lambda(r,1)=1
\]
and
\[
    \lambda(r,mn)
    =
    \lambda(r,m)\lambda(rm,n).
\]
The derivative is
\[
    (\partial_s\lambda)(r,m)
    =
    \lambda(sr,m),
\]
and the canonical output on input \(s\) is
\[
    \lambda(1,s).
\]

For a local theory \(V\), the residual representation is
\[
    \chi_V:
    M^{\mathrm{op}}
    \to
    \operatorname{KlEnd}_{\mathbb B}(V),
\]
with
\[
    \chi_V(s)(\lambda)
    =
    \bigl(
        \partial_s\lambda,
        \lambda(1,s)
    \bigr).
\]
The action is
\[
\begin{tikzcd}[column sep=huge]
    V
        \arrow[r, dashed,
        "{\lambda\mapsto(\partial_s\lambda,\lambda(1,s))}"]
    &
    V.
\end{tikzcd}
\]

Thus
\[
    m\,R_V\,n
\]
if and only if, for every \(\lambda\in V\),
\[
    \partial_m\lambda
    =
    \partial_n\lambda
\]
and
\[
    \lambda(1,m)=\lambda(1,n).
\]

\subsection{Stateless functors}
\label{subsec:to-stateless-functors}

Let
\[
    F:\mathbb A\to\mathbb B
\]
be an internal functor. Its stateless Mealy carrier is
\[
    P_F
    :=
    \left(
        A_0\xleftarrow{1_{A_0}}A_0\xrightarrow{F_0}B_0
    \right),
\]
displayed by
\[
\begin{tikzcd}
    &
    A_0
        \arrow[dl, "1_{A_0}"']
        \arrow[dr, "F_0"]
    &
    \\
    A_0
    &&
    B_0.
\end{tikzcd}
\]
The fibre over every \(v\in A_0\) is terminal.

A Kleisli operation
\[
    (P_F)_v\rightsquigarrow(P_F)_u
\]
is therefore determined by a single arrow
\[
    F_0(u)\to F_0(v)
\]
of \(\mathbb B\).

Consequently,
\[
    \operatorname{KlEnd}_{\mathbb B}(P_F)
\]
is the identity-on-objects pullback of \(\mathbb B^{\mathrm{op}}\) along
\(F_0\). Its arrow object consists internally of triples
\[
    (v,u,\alpha)
\]
with
\[
    \alpha:F_0(u)\to F_0(v).
\]
The pullback description is
\[
\begin{tikzcd}
    \operatorname{KlEnd}_{\mathbb B}(P_F)_1
        \arrow[r]
        \arrow[d]
        \arrow[dr, phantom, "\lrcorner", very near start]
    &
    B_1
        \arrow[d, "{(t_B,s_B)}"]
    \\
    A_0\times A_0
        \arrow[r, "{F_0\times F_0}"']
    &
    B_0\times B_0.
\end{tikzcd}
\]

\begin{definition}[Stateless action congruence]
\label{def:to-stateless-action-congruence}
    Define
    \[
        R_F^{\mathrm{act}}
        :=
        \ker(\chi_{F,1}),
    \]
    where
    \[
        \chi_F:
        \mathbb A^{\mathrm{op}}
        \to
        \operatorname{KlEnd}_{\mathbb B}(P_F)
    \]
    is the transition--output representation of the stateless Mealy
    morphism.
\end{definition}

\begin{proposition}[Stateless representation]
\label{prop:to-stateless-representation}
    Under the preceding identification, the transition--output representation
    \[
        \chi_F:
        \mathbb A^{\mathrm{op}}
        \to
        \operatorname{KlEnd}_{\mathbb B}(P_F)
    \]
    is the identity-on-objects lift of
    \[
        F^{\mathrm{op}}:
        \mathbb A^{\mathrm{op}}
        \to
        \mathbb B^{\mathrm{op}}.
    \]
    Hence
    \[
        a\,R_F^{\mathrm{act}}\,a'
        \quad\Longleftrightarrow\quad
        F_1(a)=F_1(a')
    \]
    for parallel input arrows \(a,a'\).
\end{proposition}

\begin{proof}
    \leavevmode
    \begin{enumerate}
        \item \textbf{Identify the state transition.}
        The state fibres are terminal, so the state transition
        \[
            (P_F)_v\to(P_F)_u
        \]
        is unique.

        \item \textbf{Identify the output comparison.}
        For
        \[
            a:u\to v,
        \]
        the output comparison is
        \[
            F_1(a):F_0(u)\to F_0(v).
        \]

        \item \textbf{Identify the representation.}
        The arrow map is therefore the lift
        \[
        \begin{tikzcd}
            A_1
                \arrow[r, "\chi_{F,1}"]
                \arrow[d, "F_1"']
            &
            \operatorname{KlEnd}_{\mathbb B}(P_F)_1
                \arrow[d]
            \\
            B_1
                \arrow[r, equal]
            &
            B_1.
        \end{tikzcd}
        \]
        This is the identity-on-objects lift of \(F^{\mathrm{op}}\).

        \item \textbf{Compute the kernel congruence.}
        Two parallel arrows induce the same operation exactly when their
        output arrows under \(F_1\) are equal. Hence
        \[
            a\,R_F^{\mathrm{act}}\,a'
            \quad\Longleftrightarrow\quad
            F_1(a)=F_1(a').
        \]
    \end{enumerate}
\end{proof}

The residual behaviour of \(F\) at \(v\) is the functor
\[
    H_v^F:\mathbb A/v\to\mathbb B
\]
defined by
\[
    H_v^F(r:u\to v)=F_0(u)
\]
and
\[
    H_v^F(r\circ a\to r)=F_1(a).
\]
Thus residual stateless semantics records every slice-restricted view of the
functor \(F\):
\[
\begin{tikzcd}[column sep=large]
    \mathbb A/v
        \arrow[r]
        \arrow[dr, "H_v^F"']
    &
    \mathbb A
        \arrow[d, "F"]
    \\
    &
    \mathbb B.
\end{tikzcd}
\]

\section{Boolean languages and classical syntactic monoids}
\label{sec:to-boolean-classical}

From this section onward, and throughout the comparison with classical
language theory and the classical Eilenberg correspondence, we specialize to
\[
    \mathcal E=\mathbf{Set}.
\]
Accordingly, the input and output categories, monoids, congruences, quotients,
languages, Boolean algebras, syntactic monoids, and pseudovarieties appearing
below are ordinary set-based structures. The internal constructions developed
above specialize along the canonical identification of internal categories,
internal functors, regular images, and effective equivalence relations in
\(\mathbf{Set}\) with their ordinary counterparts.

\subsection{Boolean residual behaviour}
\label{subsec:to-boolean-residual-behaviour}

Let
\[
    \mathbb B=2_{\mathrm{ch}}
\]
be the chaotic category on the Boolean algebra \(2\). Between any two Boolean
values there is a unique arrow, so output-comparison equality is determined
entirely by equality of endpoints.

Let
\[
    \mathbb A_\Sigma
\]
be the one-object category associated to the free monoid \(\Sigma^*\). A
residual behaviour is determined by a language
\[
    L:\Sigma^*\to2.
\]
The derivative convention is
\[
    (\partial_uL)(w)=L(uw).
\]

We identify a language
\[
    L\in
    \mathbf{Beh}_{\mathbb A_\Sigma,2_{\mathrm{ch}},0}
\]
with the singleton specification
\[
    L:
    1\to
    \mathbf{Beh}_{\mathbb A_\Sigma,2_{\mathrm{ch}},0}.
\]
Accordingly, we write
\[
    \operatorname{Der}_0(L)
    \quad\text{and}\quad
    \operatorname{SynTO}(L)
\]
for the constructions associated with this singleton specification.

For a monoid congruence \(R\) on \(\Sigma^*\), write
\[
    \eta_R:
    \Sigma^*\twoheadrightarrow\Sigma^*/R
\]
for the quotient morphism.

\begin{proposition}[Languages recognized by a congruence]
\label{prop:to-languages-recognized-congruence}
    A language \(L\subseteq\Sigma^*\) lies in \(V_R\) if and only if it is
    recognized by the quotient
    \[
        \eta_R:
        \Sigma^*\twoheadrightarrow\Sigma^*/R.
    \]
    Equivalently,
    \[
        V_R
        =
        \left\{
            \eta_R^{-1}(S)
            \mid
            S\subseteq\Sigma^*/R
        \right\}.
    \]
\end{proposition}

\begin{proof}
    \leavevmode
    \begin{enumerate}
        \item \textbf{Assume recognition by the quotient.}
        Suppose
        \[
            L=\eta_R^{-1}(S)
        \]
        for some \(S\subseteq\Sigma^*/R\).

        \item \textbf{Apply two-sided compatibility of \(R\).}
        If
        \[
            u\,R\,v,
        \]
        then for all contexts \(r,s\in\Sigma^*\),
        \[
            rus\,R\,rvs.
        \]
        Therefore
        \[
            L(rus)=L(rvs).
        \]

        \item \textbf{Apply the contextual description of \(V_R\).}
        In the chaotic output category, output comparisons are unique once
        their endpoints agree. The preceding equality therefore gives the
        contextual transition--output condition of
        Proposition~\ref{prop:to-contextual-form-VR}. Hence
        \[
            L\in V_R.
        \]

        \item \textbf{Assume membership in \(V_R\).}
        Conversely, suppose
        \[
            L\in V_R.
        \]
        Apply the contextual condition with the identity residual context and
        evaluate the derivative equality at the empty word. Then
        \[
            u\,R\,v
            \Longrightarrow
            L(u)=L(v).
        \]

        \item \textbf{Factor through the quotient.}
        The language \(L\) is constant on \(R\)-classes, so it is the inverse
        image of a subset of \(\Sigma^*/R\). Thus it is recognized by
        \(\eta_R\).
    \end{enumerate}
\end{proof}

\subsection{Recovery of the syntactic monoid}
\label{subsec:to-recovery-syntactic-monoid}

\begin{theorem}[Classical syntactic monoid recovery]
\label{thm:to-classical-syntactic-monoid}
    Let
    \[
        L\subseteq\Sigma^*
    \]
    be a language. There is a canonical isomorphism
    \[
        \operatorname{ExecMon}
        \left(
            \operatorname{SynTO}_{\mathbb A_\Sigma,2_{\mathrm{ch}}}(L)
        \right)
        \cong
        \operatorname{Syn}(L),
    \]
    where \(\operatorname{Syn}(L)\) is the classical syntactic monoid of
    \(L\).

    Equivalently,
    \[
        \operatorname{ArrMon}
        \left(
            \operatorname{SynTO}_{\mathbb A_\Sigma,2_{\mathrm{ch}}}(L)
        \right)
        \cong
        \operatorname{Syn}(L)^{\mathrm{op}}.
    \]

    If \(L\) is regular, these monoids are finite.
\end{theorem}

\begin{proof}
    \leavevmode
    \begin{enumerate}
        \item \textbf{Describe the generated derivative states.}
        The derivative states generated by \(L\) are
        \[
            \partial_rL,
            \qquad
            r\in\Sigma^*.
        \]

        \item \textbf{Compute the action of a word.}
        Acting by \(u\) sends
        \[
            \partial_rL
            \longmapsto
            \partial_u(\partial_rL)
            =
            \partial_{ru}L.
        \]
        The action diagram is
        \[
        \begin{tikzcd}[column sep=large]
            \partial_rL
                \arrow[r, dashed, "u"]
            &
            \partial_{ru}L.
        \end{tikzcd}
        \]

        \item \textbf{Discard redundant output-comparison data.}
        Since \(2_{\mathrm{ch}}\) is chaotic, equality of output comparisons
        is automatic once the derivative states agree.

        \item \textbf{Compute the syntactic congruence.}
        Therefore
        \[
            u\,R_L^{\mathrm{syn}}\,v
        \]
        if and only if
        \[
            \partial_{ru}L
            =
            \partial_{rv}L
            \qquad
            \text{for every }r\in\Sigma^*.
        \]

        \item \textbf{Evaluate the derivative equality.}
        Evaluating at \(s\in\Sigma^*\) gives
        \[
            L(rus)=L(rvs)
            \qquad
            \text{for every }r,s\in\Sigma^*.
        \]
        Hence
        \[
            u\,R_L^{\mathrm{syn}}\,v
            \quad\Longleftrightarrow\quad
            \forall r,s,\;
            rus\in L
            \Longleftrightarrow
            rvs\in L.
        \]
        This is the classical syntactic congruence.

        \item \textbf{Identify the categorical quotient.}
        Consequently,
        \[
            \operatorname{SynTO}(L)
            \cong
            \mathbb A_\Sigma^{\mathrm{op}}/R_L^{\mathrm{syn}}.
        \]

        \item \textbf{Identify the ordinary arrow monoid.}
        The ordinary arrow monoid of the quotient category has opposite
        multiplication:
        \[
            \operatorname{ArrMon}(\operatorname{SynTO}(L))
            \cong
            \operatorname{Syn}(L)^{\mathrm{op}}.
        \]

        \item \textbf{Pass to temporal execution order.}
        Taking the execution monoid reverses multiplication once more:
        \[
            \operatorname{ExecMon}(\operatorname{SynTO}(L))
            \cong
            \operatorname{Syn}(L).
        \]
        The comparison is
        \[
        \begin{tikzcd}[column sep=large]
            \Sigma^*
                \arrow[r, two heads]
            &
            \operatorname{Syn}(L)
                \arrow[r, "\cong"]
            &
            \operatorname{ExecMon}(\operatorname{SynTO}(L)).
        \end{tikzcd}
        \]

        \item \textbf{Prove finiteness for regular languages.}
        If \(L\) is regular, it has finitely many derivatives.
        Proposition~\ref{prop:to-finiteness-syntactic-category} therefore
        applies, so the syntactic transition--output category and its arrow
        and execution monoids are finite.
    \end{enumerate}
\end{proof}

\subsection{Finite Boolean saturation and right quotients}
\label{subsec:to-boolean-saturation-right-quotients}

For a language \(L\) and a word \(u\), write
\[
    Lu^{-1}
    :=
    \{w\in\Sigma^*\mid wu\in L\}
\]
for the right quotient.

\begin{theorem}[Finite Boolean saturation theorem]
\label{thm:to-finite-boolean-saturation}
    Let \(V\) be a finite Boolean algebra of languages over \(\Sigma\), closed
    under left derivatives. Then every language in \(V\) is regular, and the
    following are equivalent:
    \begin{enumerate}
        \item \(V\) is operation-saturated:
        \[
            V=V_{R_V};
        \]
        \item \(V\) is closed under right quotients.
    \end{enumerate}
\end{theorem}

\begin{proof}
    \leavevmode
    \begin{enumerate}
        \item \textbf{Establish regularity.}
        Let \(L\in V\). Since \(V\) is closed under left derivatives, every
        left derivative
        \[
            \partial_uL,
            \qquad
            u\in\Sigma^*,
        \]
        belongs to \(V\). Because \(V\) is finite, \(L\) has only finitely
        many distinct left derivatives. Hence \(L\) is regular.

        \item \textbf{Assume operation saturation.}
        Suppose
        \[
            V=V_{R_V}.
        \]
        By
        Proposition~\ref{prop:to-languages-recognized-congruence},
        every language \(L\in V\) is recognized by the quotient
        \[
            \eta_{R_V}:
            \Sigma^*\twoheadrightarrow
            \Sigma^*/R_V.
        \]
        Thus there is a subset
        \[
            S\subseteq\Sigma^*/R_V
        \]
        such that
        \[
            L=\eta_{R_V}^{-1}(S).
        \]

        \item \textbf{Show that the same quotient recognizes right quotients.}
        Let \(u\in\Sigma^*\). Define
        \[
            Su^{-1}
            :=
            \left\{
                c\in\Sigma^*/R_V
                \mid
                c\,\eta_{R_V}(u)\in S
            \right\}.
        \]
        Then
        \[
        \begin{aligned}
            w\in Lu^{-1}
            &\Longleftrightarrow
            wu\in L
            \\
            &\Longleftrightarrow
            \eta_{R_V}(wu)\in S
            \\
            &\Longleftrightarrow
            \eta_{R_V}(w)\eta_{R_V}(u)\in S
            \\
            &\Longleftrightarrow
            \eta_{R_V}(w)\in Su^{-1}.
        \end{aligned}
        \]
        Therefore
        \[
            Lu^{-1}
            =
            \eta_{R_V}^{-1}(Su^{-1})
        \]
        is recognized by the same quotient. Hence
        \[
            Lu^{-1}\in V_{R_V}=V.
        \]

        Notice that this implication does not use finiteness of \(V\):
        every operation-saturated Boolean language theory is closed under
        right quotients.

        \item \textbf{Assume closure under right quotients.}
        Conversely, suppose that \(V\) is closed under right quotients.

        For words \(x,w\in\Sigma^*\), the definition of \(R_V\) gives
        \[
            x\,R_V\,w
        \]
        if and only if
        \[
            \partial_xL
            =
            \partial_wL
            \qquad
            \text{for every }L\in V.
        \]
        Equivalently,
        \[
            L(xs)=L(ws)
            \qquad
            \text{for every }L\in V
            \text{ and every }s\in\Sigma^*.
        \]

        \item \textbf{Relate operational equivalence to membership profiles.}
        We claim that
        \[
            x\,R_V\,w
            \quad\Longleftrightarrow\quad
            \forall J\in V,\;
            J(x)=J(w).
        \]

        For the forward implication, take the empty suffix
        \(s=\varepsilon\). Then
        \[
            J(x)=J(w)
        \]
        for every \(J\in V\).

        Conversely, suppose
        \[
            J(x)=J(w)
            \qquad
            \text{for every }J\in V.
        \]
        Let \(L\in V\) and \(s\in\Sigma^*\). Since \(V\) is closed under right
        quotients,
        \[
            Ls^{-1}\in V.
        \]
        Applying membership-profile equality to \(Ls^{-1}\) gives
        \[
        \begin{aligned}
            L(xs)
            &=
            (Ls^{-1})(x)
            \\
            &=
            (Ls^{-1})(w)
            \\
            &=
            L(ws).
        \end{aligned}
        \]
        Hence
        \[
            \partial_xL=\partial_wL
        \]
        for every \(L\in V\), and therefore
        \[
            x\,R_V\,w.
        \]

        \item \textbf{Express each congruence class Booleanly.}
        Fix \(w\in\Sigma^*\). By the membership-profile characterization,
        \[
            x\in[w]_{R_V}
        \]
        if and only if \(x\) and \(w\) have the same membership value in
        every \(J\in V\). Therefore
        \[
            [w]_{R_V}
            =
            \bigcap_{\substack{J\in V\\w\in J}}J
            \cap
            \bigcap_{\substack{J\in V\\w\notin J}}
            (\Sigma^*\setminus J).
        \]
        Because \(V\) is finite and Boolean, both intersections are finite
        Boolean combinations of members of \(V\). Hence
        \[
            [w]_{R_V}\in V.
        \]

        \item \textbf{Establish finite index.}
        The \(R_V\)-class of a word is determined by its membership profile
        \[
            \bigl(J(w)\bigr)_{J\in V}\in 2^V.
        \]
        Since \(V\) is finite, there are only finitely many such profiles.
        Thus \(R_V\) has finite index.

        \item \textbf{Recover every language recognized by the quotient.}
        Let
        \[
            L\in V_{R_V}.
        \]
        By
        Proposition~\ref{prop:to-languages-recognized-congruence},
        \(L\) is a union of \(R_V\)-classes. Since \(R_V\) has finite index,
        this is a finite union.

        Every \(R_V\)-class belongs to \(V\), and \(V\) is closed under finite
        unions. Therefore
        \[
            L\in V.
        \]
        Hence
        \[
            V_{R_V}\leq V.
        \]

        \item \textbf{Use inflationarity.}
        The general operational Galois connection gives
        \[
            V\leq V_{R_V}.
        \]
        Combining the two inclusions yields
        \[
            V=V_{R_V}.
        \]
        Thus \(V\) is operation-saturated.
    \end{enumerate}
\end{proof}

\begin{remark}[Regularity and operation saturation]
\label{rem:to-regularity-operation-saturation}
    The fixed-point closure
    \[
        V\longmapsto V_{R_V}
    \]
    is formed in the complete lattice of all local behavioural theories and
    need not preserve regularity when \(R_V\) has infinite index. Therefore
    the classical Eilenberg correspondence is not obtained by simply
    restricting the local fixed-point duality to arbitrary varieties of
    regular languages. Instead, it is recovered by taking the finite
    syntactic quotients of the individual regular languages in the variety
    and closing those finite monoids under products and divisors.
\end{remark}

\begin{remark}[Categorical origins of the classical closure conditions]
\label{rem:to-classical-closure-origins}
    In the Boolean specialization:
    \begin{itemize}
        \item derivative closure is left-quotient closure;
        \item pointwise output structure gives Boolean closure;
        \item operation saturation gives right-quotient closure in the finite
        local setting;
        \item global input reindexing gives inverse morphic preimage.
    \end{itemize}
    Thus the familiar closure conditions of language varieties arise from
    distinct categorical parts of the general construction.
\end{remark}

\subsection{Residual reindexing as inverse morphic preimage}
\label{subsec:to-boolean-residual-reindexing}

\begin{proposition}[Residual reindexing of languages]
\label{prop:to-residual-reindexing-languages}
    Let
    \[
        h:\Gamma^*\to\Sigma^*
    \]
    be a monoid homomorphism, regarded as an input functor
    \[
        F_h:\mathbb A_\Gamma\to\mathbb A_\Sigma.
    \]
    For every language
    \[
        L\subseteq\Sigma^*,
    \]
    one has
    \[
        F_h^\sharp(L)=h^{-1}(L).
    \]
    Consequently, for every local theory \(V\) over \(\Sigma\),
    \[
        F_h^\bullet V
        =
        \operatorname{Im}
        \left(
            V\longrightarrow
            \mathbf{Beh}_{\mathbb A_\Gamma,2_{\mathrm{ch}},0},
            \quad
            L\longmapsto h^{-1}(L)
        \right).
    \]
\end{proposition}

\begin{proof}
    \leavevmode
    \begin{enumerate}
        \item \textbf{Compute residual precomposition.}
        For every \(w\in\Gamma^*\),
        \[
        \begin{aligned}
            F_h^\sharp(L)(w)
            &=
            L(F_h(w))
            \\
            &=
            L(h(w)).
        \end{aligned}
        \]

        \item \textbf{Identify the inverse image.}
        The characteristic function of \(h^{-1}(L)\) is exactly
        \[
            w\longmapsto L(h(w)).
        \]
        Hence
        \[
            F_h^\sharp(L)=h^{-1}(L).
        \]

        \item \textbf{Pass to local theories.}
        By definition, \(F_h^\bullet V\) is the regular image of the
        restriction of \(F_h^\sharp\) to \(V\). This is precisely the image of
        the map
        \[
            L\longmapsto h^{-1}(L).
        \]
        The factorization is
        \[
        \begin{tikzcd}[column sep=large]
            V
                \arrow[r, two heads]
                \arrow[dr, "{L\mapsto h^{-1}(L)}"']
            &
            F_h^\bullet V
                \arrow[d, hook]
            \\
            &
            2^{\Gamma^*}.
        \end{tikzcd}
        \]
    \end{enumerate}
\end{proof}

Thus the globality condition
\[
    F_h^\bullet V_\Sigma\leq V_\Gamma
\]
is exactly closure under inverse morphic preimages.

\section{Compatibility with the classical Eilenberg correspondence}
\label{sec:to-classical-eilenberg}

\subsection{Varieties of regular languages and pseudovarieties}
\label{subsec:to-varieties-pseudovarieties}

\begin{definition}[Global Boolean regular-language theory]
\label{def:to-global-boolean-regular-theory}
    A \textbf{global Boolean regular-language theory} is a family
    \[
        \mathcal V=(\mathcal V_\Sigma)_\Sigma
    \]
    indexed by finite alphabets, such that:
    \begin{enumerate}
        \item each \(\mathcal V_\Sigma\) is a Boolean algebra of regular
        languages over \(\Sigma\);
        \item each \(\mathcal V_\Sigma\) is closed under left and right
        quotients;
        \item for every monoid homomorphism
        \[
            h:\Gamma^*\to\Sigma^*,
        \]
        one has
        \[
            h^{-1}(\mathcal V_\Sigma)
            \subseteq
            \mathcal V_\Gamma.
        \]
    \end{enumerate}
\end{definition}

A \textbf{pseudovariety of finite monoids} is a nonempty,
isomorphism-closed class of finite monoids closed under finite products and
divisors. Equivalently, it is closed under finite products, submonoids, and
quotient monoids. In particular, every pseudovariety contains the trivial
monoid.

For a global Boolean regular-language theory \(\mathcal V\), define
\[
    \mathsf P_{\mathcal V}
\]
to be the pseudovariety generated by all classical syntactic monoids
\[
    \operatorname{Syn}(L),
    \qquad
    L\in\mathcal V_\Sigma,
\]
as \(\Sigma\) varies. By
Theorem~\ref{thm:to-classical-syntactic-monoid}, these are equivalently the
execution monoids
\[
    \operatorname{ExecMon}(\operatorname{SynTO}(L)).
\]

For a pseudovariety \(\mathsf P\), define
\[
    \mathcal V^{\mathsf P}_\Sigma
    :=
    \left\{
        L\subseteq\Sigma^*
        \mid
        \operatorname{Syn}(L)\in\mathsf P
    \right\}.
\]

\subsection{Syntactic separation}
\label{subsec:to-syntactic-separation}

\begin{lemma}[Syntactic fibres]
\label{lem:to-syntactic-fibres}
    Let \(L\subseteq\Sigma^*\) be regular. Every fibre of its syntactic
    morphism
    \[
        \eta_L:
        \Sigma^*\twoheadrightarrow\operatorname{Syn}(L)
    \]
    is a Boolean combination of two-sided quotients
    \[
        x^{-1}Ly^{-1}
        :=
        \{w\in\Sigma^*\mid xwy\in L\}.
    \]
\end{lemma}

\begin{proof}
    \leavevmode
    \begin{enumerate}
        \item \textbf{Use syntactic separation.}
        Two words lie in the same syntactic congruence class precisely when no
        two-sided context
        \[
            x(-)y
        \]
        separates them with respect to membership in \(L\).

        \item \textbf{Choose finitely many separating contexts.}
        Since \(\operatorname{Syn}(L)\) is finite, fix a syntactic class
        \(C\). For every distinct class \(D\), choose a context
        \[
            x_D(-)y_D
        \]
        separating \(C\) from \(D\).

        \item \textbf{Express the class as a Boolean combination.}
        For each \(D\), the chosen quotient
        \[
            x_D^{-1}Ly_D^{-1}
        \]
        either contains \(C\) and excludes \(D\), or its complement does.
        The finite intersection of these selected quotient conditions is
        exactly \(C\).

        \item \textbf{Conclude.}
        Hence each fibre of \(\eta_L\) is a Boolean combination of two-sided
        quotients of \(L\).
    \end{enumerate}
\end{proof}

\subsection{Finite syntactic generation}
\label{subsec:to-finite-syntactic-generation}

\begin{lemma}[Finite monoids from syntactic quotients]
\label{lem:to-finite-monoid-syntactic-generation}
    Let \(M\) be a finite monoid generated by a finite alphabet \(\Sigma\).
    Then \(M\) divides a finite product of syntactic monoids of languages
    recognized by \(M\).
\end{lemma}

\begin{proof}
    \leavevmode
    \begin{enumerate}
        \item \textbf{Choose a finite generating morphism.}
        Since \(M\) is generated by the finite alphabet \(\Sigma\), choose a
        surjective monoid morphism
        \[
            h:\Sigma^*\twoheadrightarrow M.
        \]

        \item \textbf{Form languages recognized by \(M\).}
        For every subset \(S\subseteq M\), define
        \[
            L_S:=h^{-1}(S).
        \]
        The language \(L_S\) is recognized by \(M\) through the morphism
        \(h\).

        \item \textbf{Compare the kernel of \(h\) with the syntactic
        congruence.}
        Let
        \[
            x,y\in\Sigma^*
        \]
        satisfy
        \[
            h(x)=h(y).
        \]
        For every pair of two-sided contexts
        \[
            p,q\in\Sigma^*,
        \]
        one has
        \[
        \begin{aligned}
            h(pxq)
            &=
            h(p)h(x)h(q)
            \\
            &=
            h(p)h(y)h(q)
            \\
            &=
            h(pyq).
        \end{aligned}
        \]
        Consequently,
        \[
            pxq\in L_S
            \quad\Longleftrightarrow\quad
            h(pxq)\in S
            \quad\Longleftrightarrow\quad
            h(pyq)\in S
            \quad\Longleftrightarrow\quad
            pyq\in L_S.
        \]
        Thus
        \[
            x\equiv_{L_S}y,
        \]
        where \(\equiv_{L_S}\) is the classical syntactic congruence of
        \(L_S\). Hence
        \[
            \ker(h)
            \leq
            {\equiv}_{L_S}.
        \]

        \item \textbf{Factor the syntactic morphism through \(M\).}
        Let
        \[
            \eta_{L_S}:
            \Sigma^*
            \twoheadrightarrow
            \operatorname{Syn}(L_S)
        \]
        be the syntactic morphism.

        Since \(\eta_{L_S}\) coequalizes \(\ker(h)\), there is a unique monoid
        morphism
        \[
            \rho_S:
            M
            \to
            \operatorname{Syn}(L_S)
        \]
        satisfying
        \[
            \rho_Sh=\eta_{L_S}.
        \]
        The factorization is
        \[
        \begin{tikzcd}[column sep=large]
            \Sigma^*
                \arrow[r, two heads, "h"]
                \arrow[dr, two heads, "\eta_{L_S}"']
            &
            M
                \arrow[d, dashed, "\rho_S"]
            \\
            &
            \operatorname{Syn}(L_S).
        \end{tikzcd}
        \]

        Because both \(h\) and \(\eta_{L_S}\) are surjective,
        \(\rho_S\) is surjective. Indeed, if
        \(z\in\operatorname{Syn}(L_S)\), choose \(x\in\Sigma^*\) with
        \(\eta_{L_S}(x)=z\). Then
        \[
            z
            =
            \eta_{L_S}(x)
            =
            \rho_S(h(x)).
        \]

        \item \textbf{Choose a finite separating family of subsets.}
        Since \(M\) is finite, the family of singleton subsets
        \[
            \bigl\{\{m\}\mid m\in M\bigr\}
        \]
        is finite and separates distinct elements of \(M\).

        For each \(m\in M\), write
        \[
            L_m:=h^{-1}(\{m\})
        \]
        and
        \[
            \rho_m:
            M\twoheadrightarrow
            \operatorname{Syn}(L_m)
        \]
        for the preceding quotient morphism.

        \item \textbf{Prove joint injectivity.}
        Form the product morphism
        \[
            \rho
            :=
            \langle\rho_m\rangle_{m\in M}:
            M
            \to
            \prod_{m\in M}
            \operatorname{Syn}(L_m).
        \]

        Suppose \(a,b\in M\) with \(a\neq b\). Choose words
        \[
            x,y\in\Sigma^*
        \]
        such that
        \[
            h(x)=a,
            \qquad
            h(y)=b.
        \]
        Consider the singleton language
        \[
            L_a=h^{-1}(\{a\}).
        \]
        Then
        \[
            x\in L_a,
            \qquad
            y\notin L_a.
        \]
        The empty two-sided context therefore separates \(x\) and \(y\), so
        \[
            \eta_{L_a}(x)
            \neq
            \eta_{L_a}(y).
        \]
        Using
        \[
            \eta_{L_a}
            =
            \rho_a h,
        \]
        one obtains
        \[
            \rho_a(a)
            \neq
            \rho_a(b).
        \]
        Hence
        \[
            \rho(a)\neq\rho(b).
        \]
        Thus \(\rho\) is injective.

        \item \textbf{Conclude divisibility.}
        The morphism \(\rho\) identifies \(M\) with a submonoid of the finite
        product
        \[
            \prod_{m\in M}
            \operatorname{Syn}(L_m).
        \]
        Each language \(L_m\) is recognized by \(M\). Therefore \(M\) divides
        a finite product of syntactic monoids of languages recognized by
        \(M\).
    \end{enumerate}
\end{proof}

\subsection{Recognition through divisors of products}
\label{subsec:to-recognition-divisors-products}

\begin{lemma}[Recognition by products of syntactic monoids]
\label{lem:to-recognition-products-syntactic}
    Let \(\mathcal V\) be a global Boolean regular-language theory. Suppose
    \(L\subseteq\Sigma^*\) is recognized by a finite monoid \(D\) that divides
    a finite product
    \[
        \prod_{i=1}^n\operatorname{Syn}(L_i),
        \qquad
        L_i\in\mathcal V_{\Sigma_i},
    \]
    where \(n\geq0\), and where the case \(n=0\) is interpreted as the trivial
    monoid. Then
    \[
        L\in\mathcal V_\Sigma.
    \]
\end{lemma}

\begin{proof}
    \leavevmode
    \begin{enumerate}
        \item \textbf{Handle the empty-product case.}
        Suppose first that
        \[
            n=0.
        \]
        Then
        \[
            \prod_{i=1}^0\operatorname{Syn}(L_i)
        \]
        is the trivial monoid. Every divisor of the trivial monoid is
        trivial, so \(D\) is the trivial monoid.

        A language recognized by the trivial monoid is either
        \[
            \varnothing
        \]
        or
        \[
            \Sigma^*.
        \]
        Since \(\mathcal V_\Sigma\) is a Boolean algebra, both languages
        belong to \(\mathcal V_\Sigma\). Hence the result holds when \(n=0\).

        For the remainder of the proof, assume
        \[
            n\geq1.
        \]

        \item \textbf{Choose divisor data.}
        Choose a submonoid
        \[
            j:S\hookrightarrow
            \prod_{i=1}^n\operatorname{Syn}(L_i)
        \]
        and a surjective homomorphism
        \[
            p:S\twoheadrightarrow D.
        \]

        \item \textbf{Choose a recognizing morphism.}
        Let
        \[
            \varphi:\Sigma^*\to D
        \]
        recognize \(L\). Thus there is a subset
        \[
            T\subseteq D
        \]
        such that
        \[
            L=\varphi^{-1}(T).
        \]

        \item \textbf{Lift through the divisor quotient.}
        Since \(\Sigma^*\) is free and \(p\) is surjective, choose, for each
        generator \(a\in\Sigma\), an element
        \[
            \widetilde a\in S
        \]
        such that
        \[
            p(\widetilde a)=\varphi(a).
        \]
        These choices extend uniquely to a monoid homomorphism
        \[
            \widetilde\varphi:\Sigma^*\to S
        \]
        satisfying
        \[
            p\widetilde\varphi=\varphi:
        \]
        \[
        \begin{tikzcd}[column sep=large]
            \Sigma^*
                \arrow[r, "\widetilde\varphi"]
                \arrow[dr, "\varphi"']
            &
            S
                \arrow[d, two heads, "p"]
            \\
            &
            D.
        \end{tikzcd}
        \]

        \item \textbf{Map into the product of syntactic monoids.}
        Put
        \[
            \psi:=j\widetilde\varphi:
            \Sigma^*\to
            \prod_{i=1}^n\operatorname{Syn}(L_i),
        \]
        and write
        \[
            \varphi_i:
            \Sigma^*\to\operatorname{Syn}(L_i)
        \]
        for its \(i\)-th component.

        \item \textbf{Lift each component to a free monoid.}
        Since
        \[
            \eta_{L_i}:
            \Sigma_i^*\twoheadrightarrow
            \operatorname{Syn}(L_i)
        \]
        is surjective and \(\Sigma^*\) is free, each \(\varphi_i\) lifts to a
        monoid homomorphism
        \[
            h_i:\Sigma^*\to\Sigma_i^*
        \]
        satisfying
        \[
            \eta_{L_i}h_i=\varphi_i:
        \]
        \[
        \begin{tikzcd}[column sep=large]
            \Sigma^*
                \arrow[r, "h_i"]
                \arrow[dr, "\varphi_i"']
            &
            \Sigma_i^*
                \arrow[d, two heads, "\eta_{L_i}"]
            \\
            &
            \operatorname{Syn}(L_i).
        \end{tikzcd}
        \]

        \item \textbf{Show that component fibres lie in the variety.}
        Fix
        \[
            c\in\operatorname{Syn}(L_i).
        \]
        One has
        \[
            \varphi_i^{-1}(c)
            =
            h_i^{-1}
            \left(
                \eta_{L_i}^{-1}(c)
            \right).
        \]
        By Lemma~\ref{lem:to-syntactic-fibres}, the syntactic fibre
        \[
            \eta_{L_i}^{-1}(c)
        \]
        is a Boolean combination of two-sided quotients of \(L_i\).
        Since
        \[
            L_i\in\mathcal V_{\Sigma_i}
        \]
        and \(\mathcal V\) is closed under left and right quotients, Boolean
        operations, and inverse morphic preimages, it follows that
        \[
            \varphi_i^{-1}(c)
            \in
            \mathcal V_\Sigma.
        \]

        \item \textbf{Show that product fibres lie in the variety.}
        For a tuple
        \[
            c=(c_1,\dots,c_n)
            \in
            \prod_{i=1}^n\operatorname{Syn}(L_i),
        \]
        one has
        \[
            \psi^{-1}(c)
            =
            \bigcap_{i=1}^n
            \varphi_i^{-1}(c_i).
        \]
        Since \(\mathcal V_\Sigma\) is closed under finite intersections,
        every fibre of \(\psi\) lies in \(\mathcal V_\Sigma\).

        \item \textbf{Restrict to the submonoid.}
        Since \(j\) is injective, for every \(s\in S\),
        \[
            \widetilde\varphi^{-1}(s)
            =
            \psi^{-1}(j(s)).
        \]
        Thus every fibre of \(\widetilde\varphi\) lies in
        \(\mathcal V_\Sigma\).

        \item \textbf{Descend to the quotient monoid.}
        For every \(d\in D\),
        \[
            \varphi^{-1}(d)
            =
            \bigcup_{s\in p^{-1}(d)}
            \widetilde\varphi^{-1}(s).
        \]
        The monoid \(S\) is finite because it is a submonoid of a finite
        product of finite monoids. Hence each fibre \(p^{-1}(d)\) is finite,
        so
        \[
            \varphi^{-1}(d)
            \in
            \mathcal V_\Sigma.
        \]

        \item \textbf{Recover the recognized language.}
        Since
        \[
            L=\varphi^{-1}(T),
        \]
        one has
        \[
            L
            =
            \bigcup_{d\in T}
            \varphi^{-1}(d).
        \]
        The monoid \(D\) is finite, so this is a finite union of languages in
        \(\mathcal V_\Sigma\). Therefore
        \[
            L\in\mathcal V_\Sigma.
        \]
    \end{enumerate}
\end{proof}

\subsection{The classical correspondence}
\label{subsec:to-classical-correspondence}

\begin{lemma}[Normal form for generated pseudovarieties]
\label{lem:to-generated-pseudovariety-normal-form}
    Let \(\mathcal G\) be a class of finite monoids. The pseudovariety
    generated by \(\mathcal G\) consists exactly of the finite monoids that
    divide a finite product
    \[
        G_1\times\cdots\times G_n
    \]
    with
    \[
        G_i\in\mathcal G.
    \]
    The case \(n=0\) is interpreted as the trivial monoid.
\end{lemma}

\begin{proof}
    Let \(\mathsf D(\mathcal G)\) be the class of finite monoids that divide
    finite products of members of \(\mathcal G\).

    The class \(\mathsf D(\mathcal G)\) contains every member of
    \(\mathcal G\), since each \(G\in\mathcal G\) divides the one-factor
    product \(G\). It also contains the trivial monoid, viewed as the empty
    product.

    It is closed under finite products. Indeed, if \(M\) divides
    \[
        G_1\times\cdots\times G_m
    \]
    and \(N\) divides
    \[
        H_1\times\cdots\times H_n,
    \]
    then \(M\times N\) divides
    \[
        G_1\times\cdots\times G_m
        \times
        H_1\times\cdots\times H_n.
    \]

    It is also closed under divisors, because divisibility of finite monoids
    is transitive. Hence \(\mathsf D(\mathcal G)\) is a pseudovariety
    containing \(\mathcal G\).

    Conversely, every pseudovariety containing \(\mathcal G\) contains every
    finite product of members of \(\mathcal G\), and therefore contains every
    divisor of such a product. Thus every such pseudovariety contains
    \(\mathsf D(\mathcal G)\).

    It follows that \(\mathsf D(\mathcal G)\) is precisely the pseudovariety
    generated by \(\mathcal G\).
\end{proof}

\begin{theorem}[Compatibility with the classical Eilenberg theorem]
\label{thm:to-classical-eilenberg}
    The assignments
    \[
        \mathcal V
        \longmapsto
        \mathsf P_{\mathcal V},
        \qquad
        \mathsf P
        \longmapsto
        \mathcal V^{\mathsf P}
    \]
    are mutually inverse and give the classical Eilenberg correspondence
    between varieties of regular languages and pseudovarieties of finite
    monoids.
\end{theorem}

\begin{proof}
    \leavevmode
    \begin{enumerate}
        \item \textbf{Relate the categorical and classical syntactic objects.}
        By Theorem~\ref{thm:to-classical-syntactic-monoid}, for every regular
        language \(L\),
        \[
            \operatorname{ExecMon}
            \bigl(
                \operatorname{SynTO}(L)
            \bigr)
            \cong
            \operatorname{Syn}(L).
        \]
        Thus the finite monoids used in the classical construction are
        exactly the execution monoids of the categorical syntactic
        transition--output objects.

        \item \textbf{Verify that \(\mathsf P_{\mathcal V}\) is a
        pseudovariety.}
        By definition,
        \[
            \mathsf P_{\mathcal V}
        \]
        is the pseudovariety generated by the finite monoids
        \[
            \operatorname{Syn}(L),
            \qquad
            L\in\mathcal V_\Sigma,
        \]
        as \(\Sigma\) varies. Hence it is nonempty, isomorphism-closed, and
        closed under finite products and divisors.

        \item \textbf{Verify regularity of
        \(\mathcal V^{\mathsf P}\).}
        Let
        \[
            L\in\mathcal V^{\mathsf P}_\Sigma.
        \]
        Then
        \[
            \operatorname{Syn}(L)\in\mathsf P.
        \]
        Every member of a pseudovariety of finite monoids is finite, so
        \(\operatorname{Syn}(L)\) is finite. Therefore \(L\) is regular.

        \item \textbf{Verify Boolean closure.}
        The empty and full languages are recognized by the trivial monoid,
        which belongs to every pseudovariety.

        Let
        \[
            L,K\in\mathcal V^{\mathsf P}_\Sigma.
        \]
        Their syntactic monoids
        \[
            M:=\operatorname{Syn}(L),
            \qquad
            N:=\operatorname{Syn}(K)
        \]
        belong to \(\mathsf P\). Hence
        \[
            M\times N\in\mathsf P.
        \]
        Every Boolean combination of \(L\) and \(K\) is recognized by
        \(M\times N\). Its syntactic monoid therefore divides
        \(M\times N\), and consequently belongs to \(\mathsf P\). Thus
        \(\mathcal V^{\mathsf P}_\Sigma\) is a Boolean algebra.

        \item \textbf{Verify left- and right-quotient closure.}
        Let
        \[
            L\in\mathcal V^{\mathsf P}_\Sigma
        \]
        and let \(M=\operatorname{Syn}(L)\). Every left or right quotient of
        \(L\) is recognized by \(M\). The syntactic monoid of the quotient
        therefore divides \(M\), and belongs to \(\mathsf P\). Hence
        \(\mathcal V^{\mathsf P}_\Sigma\) is closed under left and right
        quotients.

        \item \textbf{Verify inverse-morphic-preimage closure.}
        Let
        \[
            h:\Gamma^*\to\Sigma^*
        \]
        be a monoid homomorphism and let
        \[
            L\in\mathcal V^{\mathsf P}_\Sigma.
        \]
        If
        \[
            \eta:\Sigma^*\to M
        \]
        recognizes \(L\), then
        \[
            \eta h:\Gamma^*\to M
        \]
        recognizes \(h^{-1}(L)\). Taking
        \[
            M=\operatorname{Syn}(L)\in\mathsf P,
        \]
        the syntactic monoid of \(h^{-1}(L)\) divides \(M\), and hence belongs
        to \(\mathsf P\). Therefore
        \[
            h^{-1}(L)
            \in
            \mathcal V^{\mathsf P}_\Gamma.
        \]

        \item \textbf{Conclude that
        \(\mathcal V^{\mathsf P}\) is a global Boolean regular-language
        theory.}
        The preceding steps establish regularity, Boolean closure, closure
        under both quotients, and closure under inverse morphic preimages.
        Thus
        \[
            \mathcal V^{\mathsf P}
        \]
        is a global Boolean regular-language theory.

        \item \textbf{Prove
        \(\mathsf P_{\mathcal V^{\mathsf P}}\subseteq\mathsf P\).}
        Every generator of
        \[
            \mathsf P_{\mathcal V^{\mathsf P}}
        \]
        has the form
        \[
            \operatorname{Syn}(L)
        \]
        with
        \[
            L\in\mathcal V^{\mathsf P}_\Sigma.
        \]
        By definition of \(\mathcal V^{\mathsf P}\),
        \[
            \operatorname{Syn}(L)\in\mathsf P.
        \]
        Since \(\mathsf P\) is closed under finite products and divisors, it
        contains the pseudovariety generated by these monoids. Hence
        \[
            \mathsf P_{\mathcal V^{\mathsf P}}
            \subseteq
            \mathsf P.
        \]

        \item \textbf{Prove
        \(\mathsf P\subseteq\mathsf P_{\mathcal V^{\mathsf P}}\).}
        Let
        \[
            M\in\mathsf P.
        \]
        Since \(M\) is finite, it is generated by some finite alphabet.
        Lemma~\ref{lem:to-finite-monoid-syntactic-generation} gives languages
        \[
            L_1,\dots,L_n
        \]
        recognized by \(M\) such that \(M\) divides
        \[
            \prod_{i=1}^n\operatorname{Syn}(L_i).
        \]

        Each \(\operatorname{Syn}(L_i)\) divides \(M\), because \(M\)
        recognizes \(L_i\). Since \(M\in\mathsf P\) and \(\mathsf P\) is
        divisor-closed,
        \[
            \operatorname{Syn}(L_i)\in\mathsf P.
        \]
        Therefore
        \[
            L_i\in\mathcal V^{\mathsf P}.
        \]
        The monoids \(\operatorname{Syn}(L_i)\) are consequently generators
        of
        \[
            \mathsf P_{\mathcal V^{\mathsf P}}.
        \]
        Since that class is closed under products and divisors, it contains
        \(M\). Thus
        \[
            \mathsf P
            \subseteq
            \mathsf P_{\mathcal V^{\mathsf P}}.
        \]

        \item \textbf{Conclude the monoid equality.}
        Combining the two preceding inclusions gives
        \[
            \mathsf P_{\mathcal V^{\mathsf P}}
            =
            \mathsf P.
        \]

        \item \textbf{Prove
        \(\mathcal V\subseteq
        \mathcal V^{\mathsf P_{\mathcal V}}\).}
        Let
        \[
            L\in\mathcal V_\Sigma.
        \]
        Then
        \[
            \operatorname{Syn}(L)
        \]
        is one of the generators used to define
        \(\mathsf P_{\mathcal V}\). Hence
        \[
            \operatorname{Syn}(L)\in\mathsf P_{\mathcal V},
        \]
        and therefore
        \[
            L\in
            \mathcal V^{\mathsf P_{\mathcal V}}_\Sigma.
        \]

        \item \textbf{Prove
        \(\mathcal V^{\mathsf P_{\mathcal V}}\subseteq\mathcal V\).}
        Let
        \[
            L\in
            \mathcal V^{\mathsf P_{\mathcal V}}_\Sigma.
        \]
        Then
        \[
            \operatorname{Syn}(L)
            \in
            \mathsf P_{\mathcal V}.
        \]
        By
        Lemma~\ref{lem:to-generated-pseudovariety-normal-form}, there exist
        languages
        \[
            L_i\in\mathcal V_{\Sigma_i},
            \qquad
            1\leq i\leq n,
        \]
        such that
        \[
            \operatorname{Syn}(L)
        \]
        divides the finite product
        \[
            \prod_{i=1}^n
            \operatorname{Syn}(L_i).
        \]

        The language \(L\) is recognized by its own syntactic monoid.
        Therefore it is recognized by a finite monoid dividing the displayed
        product. Lemma~\ref{lem:to-recognition-products-syntactic} now gives
        \[
            L\in\mathcal V_\Sigma.
        \]
        Hence
        \[
            \mathcal V^{\mathsf P_{\mathcal V}}
            \subseteq
            \mathcal V.
        \]

        \item \textbf{Conclude the language equality.}
        The two language inclusions give
        \[
            \mathcal V^{\mathsf P_{\mathcal V}}
            =
            \mathcal V.
        \]

        \item \textbf{Conclude the correspondence.}
        The assignments
        \[
            \mathcal V\longmapsto\mathsf P_{\mathcal V},
            \qquad
            \mathsf P\longmapsto\mathcal V^{\mathsf P}
        \]
        are mutually inverse. They are order-preserving in the ordinary
        inclusion orders, and therefore give the classical Eilenberg
        correspondence.
    \end{enumerate}
\end{proof}

The classical correspondence may be summarized by
\[
\begin{tikzcd}[column sep=huge]
    \left\{
    \begin{array}{c}
        \text{varieties of}\\
        \text{regular languages}
    \end{array}
    \right\}
        \arrow[r, shift left=1.2ex,
        "{\mathcal V\mapsto\mathsf P_{\mathcal V}}"]
    &
    \left\{
    \begin{array}{c}
        \text{pseudovarieties}\\
        \text{of finite monoids}
    \end{array}
    \right\}
        \arrow[l, shift left=1.2ex,
        "{\mathsf P\mapsto\mathcal V^{\mathsf P}}"]
\end{tikzcd}
\]

\begin{remark}[Fixed-point duality and pseudovarieties]
\label{rem:to-fixed-points-pseudovarieties}
    The local fixed-point duality concerns a theory and the single possibly
    non-finite quotient determined by all of its operational equations.
    The classical Eilenberg correspondence packages the finite syntactic
    quotients of individual languages into a pseudovariety closed under
    products and divisors. The subdirect-product theorem and the syntactic
    universal property explain categorically why those algebraic closure
    operations arise.
\end{remark}

\section{Synthesis: state minimization and operation minimization}
\label{sec:to-synthesis}

Residual semantics gives two related minimization processes.

At the state level, the semantic image factorization
\[
\begin{tikzcd}[column sep=large]
    P
        \arrow[r, two heads]
        \arrow[dr, "\beta_P"']
    &
    \operatorname{Beh}(P)
        \arrow[d, hook]
    \\
    &
    \mathbf{Beh}^{\mathrm{can}}
\end{tikzcd}
\]
identifies states having the same residual behaviour. For a specification
\(K\), reachability identifies the relevant semantic image with
\[
    \operatorname{Der}_0(K).
\]

At the raw operation level, the image factorization
\[
\begin{tikzcd}[column sep=large, row sep=large]
    \mathbb A^{\mathrm{op}}
        \arrow[rr, "\chi_P"]
        \arrow[dr, two heads]
    &&
    \operatorname{KlEnd}_{\mathbb B}(P)
    \\
    &
    \operatorname{ActIm}(P)
        \arrow[ur, hook]
    &
\end{tikzcd}
\]
identifies input arrows having the same effect on the original states of
\(P\).

At the semantic operation level, the representation
\[
    \chi_{\operatorname{Beh}(P)}:
    \mathbb A^{\mathrm{op}}
    \to
    \operatorname{KlEnd}_{\mathbb B}(\operatorname{Beh}(P))
\]
has image factorization
\[
\begin{tikzcd}[column sep=large, row sep=large]
    \mathbb A^{\mathrm{op}}
        \arrow[rr, "\chi_{\operatorname{Beh}(P)}"]
        \arrow[dr, two heads]
    &&
    \operatorname{KlEnd}_{\mathbb B}(\operatorname{Beh}(P))
    \\
    &
    \operatorname{Op}(P)
        \arrow[ur, hook]
    &
\end{tikzcd}
\]
It identifies input arrows having the same effect after states with identical
residual behaviour have been identified.

Let \(P\) be a realization of a specification \(K\). For \(K\), the
corresponding faithful quotient is
\[
    \operatorname{SynTO}_{\mathbb A,\mathbb B}(K).
\]
The complete two-level reduction is
\[
\begin{tikzcd}[column sep=large, row sep=large]
    \text{states}
    &
    P
        \arrow[r, two heads]
    &
    \operatorname{Beh}(P)
        \arrow[r, hook]
    &
    \mathbf{Beh}^{\mathrm{can}}
    \\
    \text{operations}
    &
    \mathbb A^{\mathrm{op}}
        \arrow[r, two heads]
    &
    \operatorname{ActIm}(P)
        \arrow[r, two heads]
    &
    \operatorname{Op}(P)
        \arrow[r, two heads]
    &
    \operatorname{SynTO}(K).
\end{tikzcd}
\]
These constructions are summarized by
\[
\begin{array}{c|c}
\text{state semantics}
&
\text{operation semantics}
\\
\hline
\beta_P
&
\chi_P
\\[0.4em]
\operatorname{Beh}(P)
&
\operatorname{ActIm}(P)\twoheadrightarrow\operatorname{Op}(P)
\\[0.4em]
\operatorname{Der}_0(K)
&
\operatorname{SynTO}(K)
\\[0.4em]
\text{behavioural kernel}
&
\text{operation congruence}
\\[0.4em]
\text{reachable and separated realization}
&
\text{faithful syntactic quotient}
\end{array}
\]

The operational satisfaction polarity
\[
    V\leq V_R
    \quad\Longleftrightarrow\quad
    R\leq R_V
\]
shows that behavioural theories and operational equations determine one
another up to semantic completion:
\[
\begin{tikzcd}[column sep=huge]
    V
        \arrow[r, mapsto]
    &
    R_V
        \arrow[r, mapsto]
    &
    V_{R_V}=c(V),
    \qquad
    R
        \arrow[r, mapsto]
    &
    V_R
        \arrow[r, mapsto]
    &
    R_{V_R}=d(R).
\end{tikzcd}
\]
Its fixed points yield the local and global algebraic Eilenberg dualities.

For a behaviour specification \(K\), the syntactic transition--output category
\[
    \operatorname{SynTO}(K)
\]
is therefore simultaneously:
\begin{enumerate}
    \item the regular image of the residual transition--output
    representation;
    \item the quotient of \(\mathbb A^{\mathrm{op}}\) by the equations valid
    in every generated residual context;
    \item the unique faithful quotient recognizing
    \(\operatorname{Der}_0(K)\);
    \item the operation-level counterpart of the minimal residual realization.
\end{enumerate}

These characterizations fit into
\[
\begin{tikzcd}[column sep=large, row sep=large]
    \mathbb A^{\mathrm{op}}
        \arrow[r, two heads, "q_K^{\mathrm{syn}}"]
        \arrow[dr, two heads,
        "{q_{R_K^{\mathrm{syn}}}}"']
    &
    \operatorname{SynTO}(K)
        \arrow[r, hook]
        \arrow[d, "\cong"]
    &
    \operatorname{KlEnd}_{\mathbb B}
    \bigl(\operatorname{Der}_0(K)\bigr)
    \\
    &
    \mathbb A^{\mathrm{op}}/R_K^{\mathrm{syn}}.
\end{tikzcd}
\]

In the one-object Boolean specialization, the ordinary arrow monoid of this
category is the opposite of the classical syntactic monoid, while its
execution monoid is canonically the classical syntactic monoid itself:
\[
    \operatorname{ExecMon}(\operatorname{SynTO}(L))
    \cong
    \operatorname{Syn}(L).
\]
Under finite-state recognizability and global closure, the resulting algebraic
theory recovers Eilenberg's correspondence between varieties of regular
languages and pseudovarieties of finite monoids.